% SIAM Article Template
\documentclass[review,hidelinks,onefignum,onetabnum]{siamart250211}


%%%%%%%%%%%%%%%%%%%%%%%%%%%%%%%%%%%%
% \usepackage{amsthm}
\usepackage[numbers,sort]{natbib} % load natbib package for references
%%%%%%%%%%%%%%%%%%%%%%%%%%%%%%
% \usepackage{xcolor}
% \usepackage[bookmarks,colorlinks]{hyperref}
\definecolor{darkblue}{rgb}{0,0.22,0.66}
\definecolor{linkgray}{rgb}{0.36, 0.36, 0.36}
\hypersetup{colorlinks=true,
citecolor=cyan,linkcolor=darkblue,urlcolor=magenta}
\hypersetup{allcolors=linkgray}
\interfootnotelinepenalty=10000
% %%%%%%%%%%%%%%%%%%%%%%%%%%
\input{sisc_preamble}
%%%%%%%%%%%%%%%%%%%%%%%%%%%%%%%%
% pre defined cmds
% Use A and alpha as examples
%%%%%%%%%%%%%%%%%%%%
% \calA  means \mathcal{A}
% \calalpha  means \mathcal{\alpha}

% \cal can be replaced by \bm, \scr, \tilde, \hat, \bmcal, \bmscr, \tildebm, \hatbm, and so on


%%%%%%%%%%%%%%%%%%%%%%%%%%%%%
\renewcommand{\footercopyright}{}
\nolinenumbers

%%%%%%%%%%%%%%%%%%%%%%%%%%%%
% SIAM Shared Information Template
% This is information that is shared between the main document and any
% supplement. If no supplement is required, then this information can
% be included directly in the main document.


% Packages and macros go here
\usepackage{lipsum}
\usepackage{amsfonts}
\usepackage{graphicx}
\usepackage{epstopdf}
\usepackage{algorithmic}
\ifpdf
  \DeclareGraphicsExtensions{.eps,.pdf,.png,.jpg}
\else
  \DeclareGraphicsExtensions{.eps}
\fi

% Add a serial/Oxford comma by default.
\newcommand{\creflastconjunction}{, and~}

% Used for creating new theorem and remark environments
\newsiamremark{remark}{Remark}
\newsiamremark{hypothesis}{Hypothesis}
\crefname{hypothesis}{Hypothesis}{Hypotheses}
\newsiamthm{claim}{Claim}
\newsiamremark{fact}{Fact}
\crefname{fact}{Fact}{Facts}
%%%%%%%%%%%%%%%%%%%%%%%%%%%%%%%
% Sets running headers as well as PDF title and authors
\headers{Multi-component and Multi-layer Neural Networks}{S. Zhang,
        H. Zhao,
        Y. Zhong, and H. Zhou}

\title{Structured and Balanced Multi-component\\ 
and Multi-layer Neural Networks
% \thanks{
% Submitted to the editors DATE.
% \funding{S. Zhang was partially supported by the 
% start-up fund P0053092
% from The Hong Kong Polytechnic University.
% H. Zhao was partially supported by NSF grants DMS-2309551, and DMS-2012860. Y. Zhong was partially supported by NSF grant DMS-2309530. H. Zhou was partially supported by NSF grant DMS-2307465.}}
}


\author{Shijun Zhang\thanks{Corresponding author. Department of Applied Mathematics, Hong Kong Polytechnic University  (\href{mailto:shijun.zhang@polyu.edu.hk}{shijun.zhang@polyu.edu.hk})
}
\and Hongkai Zhao\thanks{Department of Mathematics, Duke University (\href{mailto:zhao@math.duke.edu}{zhao@math.duke.edu})}
\and Yimin Zhong\thanks{Department of Mathematics and Statistics, Auburn University (\href{mailto:yimin.zhong@auburn.edu}{yimin.zhong@auburn.edu}})
\and Haomin Zhou\thanks{School of Mathematics,
 Georgia Institute of Technology (\href{mailto:hmzhou@math.gatech.edu}{hmzhou@math.gatech.edu})}
 }
 % \date{}
 \usepackage{amsopn}
\DeclareMathOperator{\diag}{diag}
%%%%%%%%%%%%%%%%%%%%%%%%%%%%%%%%%%%%%
%%%%%%%%%%%%%%%%%%%%%%%%%%%%%%%%%%%%%%%%%%%%  
\begin{document}

\maketitle

\begin{abstract}
    In this work, we propose a balanced multi-component and multi-layer neural network (MMNN) structure to accurately and efficiently approximate functions with complex features, in terms of both degrees of freedom and computational cost. The main idea is inspired by a multi-component approach, in which each component can be effectively approximated by a single-layer network, combined with a multi-layer decomposition strategy to capture the complexity of the target function. Although MMNNs can be viewed as a simple modification of fully connected neural networks (FCNNs) or multi-layer perceptrons (MLPs) by introducing balanced multi-component structures, they achieve a significant reduction in training parameters, a much more efficient training process, and improved accuracy compared to FCNNs or MLPs. Extensive numerical experiments demonstrate the effectiveness of MMNNs in approximating highly oscillatory functions and their ability to automatically adapt to localized features.
    Our code and implementations are available at \href{https://github.com/ShijunZhangMath/MMNN/}{GitHub}.
    % For printing purposes, the links should be explicit or use a footnote.
\end{abstract}

\begin{keywords}
structured decomposition, function compositions, 
        deep neural networks,
        rectified linear unit, Fourier analysis
\end{keywords}

% REQUIRED
\begin{MSCcodes}
65Y10, 65Y20, 68W20, 68W25
\end{MSCcodes}



%%%%%%%%%%%%%%%%%%%%%%%%%%%%%%%%%%%%%%%%%%%%%%%
%%%%%%%%%%%%%%%%%%%%%%%%%%%%%%%%%%%%%%%%%%%%%%
\section{Introduction}
\label{sec:intro}

The key use of neural networks is to approximate an input-to-output relation, i.e., a mapping or a function in mathematical terms. In this work, we continue our study of numerical understanding of neural network approximation of functions from representation to learning dynamics. In our earlier study \cite{ZZZZ-23}, we demonstrated that a one-hidden-layer (also known as a two-layer or shallow) network is essentially a ``low-pass filter" when approximating a function in practice. Due to the strong correlation among the family of activation functions (parameterized by the weight and bias), such as \ReLU\  (rectified linear unit), the Gram matrix, the element of which is the pairwise correlation (inner product) of the activation functions, has a fast spectral decay. If initialized randomly, the eigenvectors of the Gram matrix correspond to generalized Fourier modes from low frequency to high frequency ordered corresponding to decreasing eigenvalues. Due to the ill-conditioning of the representation, no matter how wide a one-hidden-layer network is, it can only learn and approximate smooth functions or sample low-frequency modes effectively and stably (with respect to noise or machine round-off errors). 

In this work, we propose a balanced multi-component and multi-layer neural network (MMNN) structure based on our previous understanding of a one-hidden-layer network. First, we show that a multi-layer network with a multi-component structure, each of which can be approximated well and effectively by a one-hidden-layer network, can overcome the limitation of a shallow network by smooth decomposition and transformation. Compared to a fully connected neural network of a similar structure, our proposed MMNN is much more effective in terms of representation, training, and accuracy in approximating functions, especially for functions containing complex features, e.g., high-frequency modes. The key idea of MMNNs is to view a linear combination of activation functions as randomly parameterized basis functions, called a \emph{component},  as a whole to represent a smooth function. Each layer has multiple components all sharing the common basis functions with different linear combinations. The number of components, called \emph{rank}, is typically much smaller than the layer's width and increases to enhance the flexibility of decomposition when dealing with more complex functions. These components are combined and composed (through layers) in a structured and balanced way in terms of network width, rank, and depth to approximate a complicated function effectively. 
%Through optimization a.k.a. learning or training, the final layer of the multi-layer network produce a set of learned basis functions, a linear combination of which can approximate the target function well. 
Another important feature we used in practice is that weights and biases inside each activation function are randomly assigned and fixed during the optimization while the linear combination weights of activation functions in each component are trained. This leads to more efficient training processes motivated by our finding that a one-hidden-layer neural network can be trained effectively to approximate a smooth function well using random basis functions. We also demonstrate interesting learning dynamics based on Adam optimizer~\cite{DBLP:journals/corr/KingmaB14}, which is crucial for the successful and efficient training of MMNNs. An important remark is that a balanced and holistic approach needs to consider both representation and optimization as well as their interplay altogether. 


The structure of this paper is as follows. Section~\ref{sec:MMNN} introduces and details the design of MMNNs. Section~\ref{sec:comparison} provides a comparison between FCNNs and MMNNs. In Section~\ref{sec:decomposition}, we present a mathematical framework for smooth decomposition and transformation based on the MMNN architecture, showing that each component can be effectively approximated by a single-hidden-layer network. Section~\ref{sec:experiments} presents extensive numerical experiments to validate our analysis and demonstrate the capability of MMNNs in approximating complex functions. Additional insights and implementation guidelines are discussed in Section~\ref{sec:further:discussion}. Finally, Section~\ref{sec:conclusion} concludes the paper with final remarks.

% The outline of this paper is summarized as follows. In Section~\ref{sec:MMNN}, the design of MMNNs is proposed and explained. Later in Section~\ref{sec:comparison}, we comapre FCnns vs MMNNs. Then, in Section~\ref{sec:decomposition}, a mathematical framework for smooth decomposition and transformation using the MMNN architecture is presented, demonstrating that each component can be effectively approximated by a one-hidden-layer network. Extensive numerical experiments are presented in Section~\ref{sec:experiments} to verify the analysis and demonstrate the effectiveness of MMNNs in the approximation of more complicated functions. Further discussion is presented in Section~\ref{sec:further:discussion}, where more insights and implementation guidelines of MMNNs are provided. Finally, remarks and conclusions are provided in Section~\ref{sec:conclusion}.


% \section{Notations and our new network architectures}

% \section{Notations and MMNN}

\section{Multi-component and multi-layer neural network (MMNN)}
\label{sec:MMNN}


In this section, we present a novel network architecture called the Multi-component and Multi-layer Neural Network (MMNN). 
Let's begin with some notations. Let $\R$ represent the set of real numbers. The indicator (or characteristic) function of a set $A$, denoted by $\one_{A}$, is a function that takes the value $1$ for elements in $A$ and $0$ for elements not in $A$. Vectors and matrices are denoted by bold lowercase and uppercase letters, respectively.
We use slicing notation for a vector $\bmx = (x_1, \cdots, x_d) \in \R^d$, where $\bmx{[n:m]}$ denotes a slice of $\bmx$ from its $n$-th to the $m$-th entries for any $n, m \in \{1, 2, \cdots, d\}$ with $n \le m$, and $\bmx{[n]}$ denotes the $n$-th entry of $\bmx$. For example, if $\bmx = (x_1, x_2, x_3) \in \R^3$, then $(5\bmx){[2:3]} = (5x_2, 5x_3)$ and $(6\bmx + 1){[3]} = 6x_3 + 1$. A similar notation is used for matrices. For instance, $\bmA[:, i]$ refers to the $i$-th column of $\bmA$, whereas $\bmA[i, :]$ indicates the $i$-th row of $\bmA$. Additionally, $\bmA[i, n:m]$ corresponds to $(\bmA[i, :])[n:m]$, extracting the entries from the $n$-th to the $m$-th in the $i$-th row.

Later in this section, we introduce the architecture of MMNNs in Section~\ref{sec:MMNN:architecture}. Following this, in Section~\ref{sec:MMNN:learning:stategy}, we outline the learning strategy of MMNN and highlight its advantages over other methods. 
%Finally, in Section~\ref{sec:comparison}, we compare the numerical performance of MMNNs and FCNNs.


\subsection{Architecture of MMNNs}
\label{sec:MMNN:architecture}

In this section, we introduce the architecture of our Multi-component and Multi-layer Neural Network (MMNN). Each layer of the MMNN is a (shallow) neural network of the form
%\subsubsection{Shallow network}
\begin{equation*}
    \bmh(\bmx)=\bmA \sigma(\bmW\bmx +\bmb)+\bmc
\end{equation*}
to approximate a vector-valued function 
$\bmf:\R^\din\to\R^\dout$
% $\bmf\in C([0,1]^\din;\R^\dout)$,
where $\bmW \in \R^{n\times \din}, \bmA\in \R^{\dout\times n}$, and $n$ is the width of this network.
Here, \(\sigma:\mathbb{R}\to\mathbb{R}\) represents the activation function that can be applied elementwise to vector inputs. Throughout this paper, the activation function is chosen as \ReLU, unless otherwise specified.
One can also write it in a more compact form, 
\begin{equation}\label{eq:shallow}
    \bmh=\bmA \sigma(\bmW\bmx +\bmb)+\bmc
    =\bmtildeA\left[\begin{matrix}
        \sigma(\bmtildeW\bmtildex) \\
         1
    \end{matrix}\right],
    % =\bmtildeA\sigma(\bmtildeW \bmtildex),
\end{equation}
where 
\begin{equation*}
    \bmtildeW=\left[\begin{matrix}
        \bmW, \bmb
    \end{matrix}\right],\quad 
    \bmtildeA=\left[\begin{matrix}
        \bmA, \bmc
    \end{matrix}\right],\quad
    \bmtildex=\left[\begin{matrix}
        \bmx\\  1
    \end{matrix}\right]
    .
\end{equation*}
We call each element of $\bmh$,  i.e., $\bmh{[i]}=\tilde{\bm{A}}{[i,:]}\cdot\left[\begin{matrix}
        \sigma(\bmtildeW\bmtildex) \\
         1
    \end{matrix}\right]$ for $i=1,2,\cdots,d_{\tn{out}}$, a component. Here are a few key features of $\bmh$:
\begin{enumerate}
    \item Each component is viewed as a linear combination of basis functions $\sigma(\bmW{[i,:]}\cdot\bmx+\bmb[i]),\, i=1,2,\cdots,n$, which is a function in $\bmx$, as a whole.
    \item Different components of $\bmh$ share the same set of basis with different coefficients $\tilde{\bm{A}}[i,:]$.
    \item Only $\bmtildeA$ is trained while $\bmtildeW$ is randomly assigned and fixed. 
    \item The output dimension $\dout$ and network width $n$ can be tuned according to the intrinsic dimension and complexity of the problem.
\end{enumerate}
In comparison, each layer in a typical deep FCNN takes the form $\sigma(\bmtildeW\bmtildex)$, and each hidden neuron is individually a function of the input $\bmx$ or each point $\bmx\in \R^{d_{\tn{in}}}$ is mapped to $\R^{n}$, where $n$ is the layer width. All weights $\bmtildeW$ are training parameters. In MMNNs, each layer is composed of multiple components $\bmtildeA\sigma(\bmtildeW\bmtildex)$. Each component is a linear combination of randomly parameterized hidden neurons $\sigma(\bmtildeW \bmtildex)$, which can be more effectively and stably trained through $\bmtildeA$ as a smooth decomposition/transformation. Typically the number of components $d_{\tn{out}}$, the dimension of intermediate feature space, is (much) smaller than the layer width $n$, the number of random neurons (or basis functions) in the MMNN. On one hand, the intermediate feature space is compressed, and on the other hand, there are diverse random basis whose linear combinations can have enough representation power in the feature space.


%\begin{equation*}
%    \{1\}\cup \big\{ \sigma(\bmtildeW{[i,:]}\cdot\bmtildex):i=1,2,\cdots,n\big\},
%\end{equation*}

% all components of $\bmh$ share the same ``bases" 
% \begin{equation*}
%     \{1\}\cup \big\{ \sigma(\bmtildeA_{[i,:]}\bmtildex:i=1,2,\cdots,n\big\},
% \end{equation*}
% but each component of $\bmh$, denoted as $\bmh_{[i]}$, has its own weights $\bmtildeW_{[i,:]}$ corresponding to the bases.

% \subsubsection{Our new network architecture designed via compositions of shallow ones}

% \subsubsection{Compositions of shallow networks}


An MMNN is a multi-layer composition of $\bmh_i$, i.e., $\bmh:\R^\din \mapsto\R^\dout$
\begin{equation*}
    \bmh=\bmh_m\circ \cdots \circ \bmh_2\circ \bmh_1,
\end{equation*}
where each $\bmh_i:\R^{d_{i-1}}\mapsto\R^{d_{i}}$ is a multi-component shallow network defined in~\eqref{eq:shallow} of width $n_i$, where 
\begin{equation*}
    d_0=\din,\qquad d_1,\cdots, d_{m-1}\ll   n_i,\qquad d_m=\dout.
\end{equation*}
The width of this MMNN is defined as $\max\{n_i: i=1,2,\cdots,m-1\}$, the rank as $\max\{d_i: i=1,2,\cdots,m-1\}$, and the depth as $m$.
MMNNs are designed to reduce the dimensionality of the intermediate feature space, thereby simplifying optimization while largely preserving the network's expressive power.
To simplify, we denote a network with width $w$, rank $r$, and depth $l$ using the compact notation $(w,r,l)$.
% \begin{equation*}
%     d_0=\din,\qquad d_1=\cdots=d_{m-1}=r\ll   n_i,\qquad d_m=\dout.
% \end{equation*}
See 
Figure~\ref{fig:MMNN:vs:FC}(a)
% Figure~\ref{subfig:MMNN}
% and \ref{fig:lowRankNetEg2}  
for an illustration of an MMNN of size $(4,2,2)$.
In contrast, an FCNN $\bmphi$ can be expressed in the following composition form
	% . Specifically, it can be written as
	\begin{equation*}
		\bmphi =\calbmL_L\circ\sigma\circ
		\calbmL_{L-1}\circ %\sigma\circ
		\ \cdots \  \circ 
		\sigma\circ
		\calbmL_1\circ\sigma\circ\calbmL_0,
	\end{equation*}
 where $\calbmL_i$ is an affine linear map given by $\calbmL_i(\bmy)=\bmW_i\cdot \bmy+\bmb_i$. Readers are referred to 
 % Figure~\ref{subfig:fc}
 Figure~\ref{fig:MMNN:vs:FC}(b)
 for an illustration and a comparison with the MMNN.


We clarify the structural differences between FCNNs and MMNNs, omitting bias terms for notational simplicity.
In MMNNs, each layer is defined by a composition
\[
\bmh: \mathbb{R}^r \to \mathbb{R}^r, \quad \bmh(\bmx) = \bmA\sigma(\bmW\bmx),
\]
while in standard FCNNs, layers are typically written as
\[
\bmh: \mathbb{R}^n \to \mathbb{R}^n, \quad \bmh(\bmx) = \sigma(\bmW\bmx),
\]
where \( r \ll n \), and \( r \) denotes the MMNN rank, an internal dimensionality that helps regulate network complexity.
The core principle behind MMNNs is to limit the dimensionality of intermediate space, which helps streamline the optimization process while maintaining the model's expressive power.

% The central design principle of MMNNs is to constrain the dimensionality of the intermediate representations. By fixing \(\bmW\) and only training \(\bmA\), we reduce the number of learnable parameters and simplify the optimization landscape, while still preserving expressivity through the nonlinear transformation \(\sigma(\bmW\bmx)\), which acts as a set of basis functions.

One might suggest that the product \(\bmA_j \bmW_{j+1}\) in MMNNs can be collapsed into a single matrix, making MMNNs a special case of FCNNs. However, this overlooks two critical distinctions:

\begin{enumerate}
    \item \textbf{Asymmetric Parameter Roles:} In MMNNs, \(\bmW\) is randomly initialized and fixed while \(\bmA\) is learnable. This asymmetry is central to our representation and training strategy and cannot be replicated by simply reinterpreting MMNNs as FCNNs with a low rank factorization of the weight matrix $\bmW$ which needs to be learned fully posing a challenging task for the optimization.
    
    \item \textbf{Architectural Interventions:} In practice, modern networks often employ techniques such as batch normalization, dropout, and residual connections, which fundamentally alter the layer-wise composition. For instance,
    \begin{itemize}
        \item With batch normalization layers $B_i$'s:
        \[
        \bmh_m \circ B_{m-1} \circ \bmh_{m-1} \circ \cdots \circ B_1 \circ \bmh_1;
        \]
        \item With dropout layers $D_i$'s:
        \[
        \bmh_m \circ D_{m-1} \circ \bmh_{m-1} \circ \cdots \circ D_1 \circ \bmh_1;
        \]
        \item With residual connections:
        \[
        \bmh_m \circ (\bmI + \bmh_{m-1}) \circ \cdots \circ (\bmI + \bmh_2)\circ   \bmh_1.
        \]
    \end{itemize}
    In these scenarios, the clean compositional structure required to merge adjacent matrices (e.g., \(\bmA_j\) and \(\bmW_{j+1}\)) is disrupted. Thus, MMNNs remain distinct in both form and training philosophy.
\end{enumerate}
For very deep MMNNs, one can borrow ideas from ResNets \cite{7780459} to address the gradient vanishing issue, making training more efficient. Incorporating this idea, we propose a new architecture given by a multi-layer composition of $\bmI+\bmh_i$, i.e., $\bmh:\R^\din \mapsto \R^\dout$
\begin{equation*}
    \bmh = \bmh_m \circ (\bmI + \bmh_{m-1}) \circ \cdots \circ (\bmI + \bmh_3) \circ (\bmI + \bmh_2) \circ \bmh_1,
\end{equation*}
where each $\bmh_i: \R^{d_{i-1}} \mapsto \R^{d_i}$ is a multi-component shallow network defined in \eqref{eq:shallow} with width $n_i$, 
\begin{equation*}
    d_0=\din,\qquad d_1=\cdots=d_{m-1}=r\ll   n_i,\qquad d_m=\dout,
\end{equation*}
and $\bmI$ is the identity map. We call this architecture ResMMNN.
See 
Figure~\ref{fig:MMNN:vs:FC}(c)
% Figure~\ref{subfig:ResMMNN}
% and \ref{fig:lowRankNetEg2}  
for an illustration of a ResMMNN of size (4,2,3).




The above definition of ResMMNNs requires \(d_1 = \cdots = d_{m-1} = r\). If this condition does not hold, we can alternatively define ResMMNNs via
\begin{equation*}
    \bmh = (\bmI \oplus \bmh_m) \circ (\bmI \oplus \bmh_{m-1}) \circ \cdots \circ (\bmI \oplus \bmh_3) \circ (\bmI \oplus \bmh_2) \circ (\bmI \oplus \bmh_1),
\end{equation*}
where \(\oplus\) is an operation defined as follows. For any functions \( \bmf:\mathbb{R}^d \mapsto \mathbb{R}^{d_\bmf} \) and \( \bmg:\mathbb{R}^d \mapsto \mathbb{R}^{d_\bmg} \), the \(\oplus\) operation is given by
\[
\bmf \oplus \bmg \coloneqq (\tildebmf + \tildebmg)[1:d_\bmg],
\]
where
\[
\tildebmf = \begin{bmatrix}
   \bmf \\  
   \bmzero 
\end{bmatrix}
\in \mathbb{R}^{\max\{d_\bmf,d_\bmg\}}
\quad \tn{and}\quad \tildebmg = \begin{bmatrix}
   \bmg \\  
   \bmzero 
\end{bmatrix}
\in \mathbb{R}^{\max\{d_\bmf,d_\bmg\}}.
\]
% \[
% \bmf \oplus \bmg \coloneqq \iota \bmf + \iota \bmg,
% \quad \tn{where}\quad 
% \iota \bmf := \begin{bmatrix}
%    \bmf \\  
%    \bmzero 
% \end{bmatrix}
% \in \mathbb{R}^{\max\{d_\bmf,d_\bmg\}},
% \quad \iota \bmg = \begin{bmatrix}
%    \bmg \\  
%    \bmzero 
% \end{bmatrix}
% \in \mathbb{R}^{\max\{d_\bmf,d_\bmg\}}.
% \]


% It is noteworthy that such an operation is both straightforward and cost-effective to implement. For example, in Python, one can use the following code:
% \begin{equation*}
%     \textsf{y = f(x);\qquad z = g(x);\qquad n = min(len(y), len(z));\qquad z[:n] = y[:n] + z[:n]}
% \end{equation*}
% After executing this code, \textsf{z} will be the result of the map \(\textsf{f} \oplus \textsf{g}\) at \textsf{x}.
% We remark that the definition of ResMMNN can be generalized to only adding identity maps to certain specific layers, which we still refer to as ResMMNN.


% \red{lr1: Learning rates for the parameters  of $h_i$ ``inside activation", i.e., $\bmtildeA_i$ }

% \red{lr2: Learning rates for the parameters  of $h_i$ ``outside activation", i.e., $\bmtildeW_i$ }

% \def\lr{\tn{lr}}
% \begin{itemize}
%     \item $0=\lr1<\lr2$: seems good for easy target functions
%     \item $0<\lr1\ll  \lr2$: seems better    \item $0<\lr1\approx \lr2$: seems best
% \end{itemize}


\begin{figure}%[htbp]%[h]%[H]
    \centering
    \begin{subfigure}[c]{0.98\textwidth}
    \centering                \includegraphics[width=0.758025\textwidth]{figures/Others/lowRankNetEg.pdf}
    \subcaption{MMNN of size $(4,2,2)$, i.e., width $4$, rank 2, and depth $2$.}
    		\label{subfig:MMNN}
    \end{subfigure}\\
    \vspace{6pt}
    \begin{subfigure}[c]{0.98\textwidth}
    \centering                \includegraphics[width=0.64\textwidth]{figures/Others/FCNetEg.pdf}
    \subcaption{FCNN of size (4, --, 2), i.e., width $4$ and depth $2$.}
    		 \label{subfig:fc}
    \end{subfigure}
    \\
    \vspace{6pt}
    \begin{subfigure}[c]{0.98\textwidth}
    \centering                \includegraphics[width=0.9985\textwidth]{figures/Others/lowRankResNetEg.pdf}
    \subcaption{ResMMNN of size $(4,2,3)$, i.e., width $4$, rank 2, and depth $3$.}
    		\label{subfig:ResMMNN}
    \end{subfigure}
\caption{Illustrations of $\sigma$-activated MMNN, FCNN, and ResMMNN. 
% Illustration
% of (a) MMNN, (b) fully connected network, and (c) ResMMNN. \\ 
% Illustration of (a) MMNN, (b) fully connected network, and (c) ResMMNN. \\ 
% (a) an MMNN of size $(4,2,2)$ - width $4$ (number of basis functions in each component), rank $2$ (number of components), and depth $2$ (number of layers) with activation function $\varrho$. (b) A FCNN of size (4,2) - width $4$ and depth $2$. (c) A ResMMNN of size $(4,2,3)$.
}
    \label{fig:MMNN:vs:FC}
\end{figure}


\subsection{Learning strategy of MMNNs}
\label{sec:MMNN:learning:stategy}


Our learning strategy is motivated by the following basic principle: a function can be decomposed in a multi-component and multi-layer structure each component of which can be approximated and trained effectively using a one-hidden-layer network, which is a linear combination of random basis functions (e.g., of the form $\sigma(\bmW_i\cdot \bmx + \bmb_i)$, see Section~\ref{sec:decomposition}).
Therefore, optimizing the linear combination weights of the random basis functions, namely $\bmA_i$'s and $\bmc_i$'s, is both computationally efficient and sufficiently expressive.
On the other hand, optimizing the weights (orientations of the basis functions) $\bmW_i$'s and biases $\bmb_i$'s to make the basis functions more adaptive to fine-tune features of the target function, which would require capturing high-frequency information by a single layer network, leads to not only significantly more parameters to optimize but also difficulties in training as shown in~\cite{ZZZZ-23}.
%inspired by two principal considerations: the simultaneous training of different layers in network learning, and the methodologies employed in the random feature method. 
Specifically, for each layer of an MMNN, we fix the activation function parameters (\(\bm{W}_i\)'s and \(\bm{b}_i\)'s) as per PyTorch's default setting during the training process. This entails initializing both weights and biases uniformly from the distribution \(\mathcal{U}(-\sqrt{k}, \sqrt{k})\), where \(k = \frac{1}{\text{in\_features}}\).\footnote{It is noteworthy that this initialization approach is similar to the widely used Xavier initialization \cite{pmlr-v9-glorot10a}, which draws weights from the distribution \(\mathcal{U}(-\sqrt{k}, \sqrt{k})\) with \(k = \frac{\sqrt{6}}{\text{in\_features} + \text{out\_features}}\) and sets the bias to $\bmzero$.}
% \footnote{It's worth noting that this initialization approach resembles the widely used Xavier initialization \cite{pmlr-v9-glorot10a}, which defines \(k\) as \(\frac{\sqrt{6}}{\text{in\_features} + \text{out\_features}}\).}
% Specifically, for each one-hidden-layer block, we fix the parameters within the activation function (i.e., \(\bmW_i\)'s and \(\bmb_i\)'s) using the default  initialization in PyTorch, i.e., both weights and biases are 
% initialized from the uniform distribution $\calU(-\sqrt{k}, \sqrt{k})$, \text{ where } k = \frac{1}{\text{in\_features}}. 
% We remark that this initilzaiton is similar to the popular one 
% --Xavier initialization \cite{pmlr-v9-glorot10a},
%  whose $k=\frac{\sqrt{6}}{\text{in\_features}+\text{out\_features}}$. 
% \begin{equation*}
%     \tn{Weight} \sim \calU\left(-\sqrt{\frac{6}{n_i + n_{i+1}}}, \sqrt{\frac{6}{n_i + n_{i+1}}}\right), \quad \tn{bias} = 0,
% \end{equation*}
The whole training process optimizes all \(\bmA_i\)'s and \(\bmc_i\)'s simultaneously using the Adam optimizer \cite{DBLP:journals/corr/KingmaB14}. Note that it is important to have a uniform sampling of orientations $\bmW_i$ and biases $\bmb_i$ for the random basis functions to be able to approximate an arbitrary smooth function well. Unless stated otherwise, parameter initialization adheres to the default settings provided by PyTorch in our experiments.


%\red{
% We would like to note that although the \(\bm{W}_i\)'s and \(\bm{b}_i\)'s are fixed, their initialization is crucial. The default PyTorch initialization is a good choice. Other types of initialization have also been experimented with, such as 1) drawing all values of the \(\bm{W}_i\)'s and \(\bm{b}_i\)'s from the default PyTorch random initialization and then incrementing them by 1, and 2) drawing all values of the \(\bm{W}_i\)'s and \(\bm{b}_i\)'s from the default PyTorch random initialization and then taking their absolute values. In both cases, the learning process failed, i.e., the training error did not decrease at all. Such results can be expected since \(\bm{W}_i\) and \(\bm{b}_i\) determine the basis functions \(\sigma(\bmW_i[j,:] \cdot \bmx + \bmb_i[j])\) at the \(i\)-th layer. We aim to ensure that the span of all these basis functions can approximate a wide variety of functions effectively.

To demonstrate the advantages of our training approach (labeled S1), we conduct a comparison with the typical strategy in deep neural networks, denoted as Strategy S2, which uses the default PyTorch initialization and optimizes all parameters during training. In our tests, we select an oscillatory target function $f(x)=\cos(36\pi x^2)-0.6\cos(12\pi x^2)$ and use fairly compact networks. The tests are performed on a total of $1000$ uniform samples in $[-1,1]$ with a mini-batch size of $100$ and a learning rate for epoch $k$ set at $0.001 \times 0.9^{\lfloor k/400 \rfloor}$ for $k = 1, 2, \cdots, 20000$, where $\lfloor \cdot \rfloor$ denotes the floor operation. The Adam optimizer \cite{DBLP:journals/corr/KingmaB14} is applied throughout the training process.
% All experiments are conducted on an \textit{NVIDIA RTX 3500 Ada Generation Laptop GPU (power cap 130W)}. 
% We would like to remark that learning rate scheduler can be a subtle and important issue in practice. For all our training process,
% the Step Learning Rate suffices. However, one could consider using other learning rate schedulers, such as the Cosine Scheduler \cite{loshchilov2017sgdr} or the gradual warm-up strategy \cite{2017arXiv170602677G}. Exploring and designing a more efficient learning rate scheduler with some automatic restart mechanism for MMNNs is a potential interesting topic to study for future work.



% \begin{table}[htbp!]
% 	\centering  
%  \setlength{\tabcolsep}{0.68em} 
%  \renewcommand{\arraystretch}{1.15}
% \caption{Comparison of test errors.}
% 	\label{tab:stategy:error:comparison}
% 	\resizebox{0.95\textwidth}{!}{ 
% 		\begin{tabular}{ccccccccc} 
% 			\toprule% \toprule[1.2pt]  
% 			    network &  (width, rank, depth) & {\#parameters (trained / all)}  &      {test error (MSE)} 
%     & test error (MAX) & training time \\
% 			 % \cmidrule{3-5}
% 			 % & & MSE & MAE & MAX\\
% 			\midrule
% 			 \rowcolor{mygray}
% 			  MMNN1 (S1) & (400, 20, 6) &  40501 / 83301 & 
% $2.12 \times 10^{-5}$  &  $4.42 \times 10^{-2}$
%     & 23.9s / 1000 epochs \\	
% 			% \midrule
%    			 MMNN1 (S2) & (400, 20, 6) & 83301 / 83301 & 
% $3.50 \times 10^{-5}$  &  $4.63 \times 10^{-2}$
%    & 30.2s / 1000 epochs
%    \\
% 		\rowcolor{mygray}	 MMNN2 (S1) & (590, 28, 6) &  83331 / 170061 & 
% $1.17 \times 10^{-5}$  &  $2.85 \times 10^{-2}$
%    & 25.2s / 1000 epochs\\			
% 			\bottomrule% \bottomrule[1.2pt] 
% 		\end{tabular} 
% 	}%%% \resizebox
% % 	\vskip -0.1in
% \end{table} 
\begin{table}[htbp!]
	\centering  
 \setlength{\tabcolsep}{0.68em} 
 \renewcommand{\arraystretch}{1.15}
\caption{
Comparison of test errors averaged over the last 100 epochs.
}
	\label{tab:stategy:error:comparison}
	\resizebox{0.95\textwidth}{!}{ 
		\begin{tabular}{ccccccccc} 
			\toprule% \toprule[1.2pt]  
			    network &  (width, rank, depth) & {\#parameters (trained / all)}  &      {test error (MSE)} 
    & test error (MAX) & training time \\
			 % \cmidrule{3-5}
			 % & & MSE & MAE & MAX\\
			\midrule
			 \rowcolor{mygray}
			  MMNN1 (S1) & (400, 20, 6) &  40501 / 83301 & 
$2.01 \times 10^{-5}$  &  $4.36 \times 10^{-2}$
    & 23.9s / 1000 epochs \\	
			% \midrule
   			 MMNN1 (S2) & (400, 20, 6) & 83301 / 83301 & 
$4.26 \times 10^{-5}$  &  $4.71 \times 10^{-2}$
   & 30.2s / 1000 epochs
   \\
		\rowcolor{mygray}	 MMNN2 (S1) & (590, 28, 6) &  83331 / 170061 & 
$\bm{1.39 \times 10^{-5}}$  &  $\bm{2.80 \times 10^{-2}}$
   & 25.2s / 1000 epochs\\			
			\bottomrule% \bottomrule[1.2pt] 
		\end{tabular} 
	}%%% \resizebox
% 	\vskip -0.1in
\end{table} 
% \red{Not perform well if Not Xavier initialization; e.g.,  W and b only cover half disk }

% \red{another eg; W and b learnable; similar results}

% We utilize a compact network for approximation, configured with a width of 800, a rank of 30, and a depth of 8 layers. The following hyperparameters are used during the training of our neural networks.  The learning rate for epoch-$k$ is set at $0.001 \times 0.9^{\lfloor k/300 \rfloor}$ for $k=1,2,\cdots,3000$.
% Mini-Batch Size is $100$.
% We employ uniform grid sampling, with $1000$ samples across the domain $[-1,1]$.

% \begin{itemize}
% \item \textbf{Learning Rate}: The learning rate for epoch-$k$ is set at $0.001 \times 0.9^{\lfloor k/300 \rfloor}$ for $k=1,2,\cdots,3000$.
% \item \textbf{Mini-Batch Size}: $100$.
% \item \textbf{Sampling}: We employ uniform grid sampling, with $1000$ samples across the domain $[-1,1]$.
% \end{itemize}

\begin{figure}%[htbp]%[h]%[htbp!]
    \centering	
        \begin{subfigure}[c]{0.32438\textwidth}
    \centering            \includegraphics[width=0.985\textwidth]{figures/Strategy/LSf1D.pdf}
    % \subcaption{}
    		% \label{subfig:f3}
    \end{subfigure}
    \hfill
\begin{subfigure}[c]{0.32438\textwidth}
	\centering            \includegraphics[width=0.9285\textwidth]{figures/Strategy/1Derrors_vsFCNN.pdf}
	% \subcaption{}
	% \label{subfig:f3}
\end{subfigure}
\hfill
\begin{subfigure}[c]{0.32438\textwidth}
    \centering            \includegraphics[width=0.9285\textwidth]{figures/Strategy/1Derrors_vsFCNN_aver.pdf}
    % \subcaption{MMNN (500,25,6).}
    		% \label{subfig:f3}
\end{subfigure}
\caption{Left: target function $f(x)=\cos(36\pi x^2)-0.6\cos(12\pi x^2)$. Middle: base-10 logarithm of test errors vs. epoch. Right: 
base-10 logarithm of ``test-error-aver'' vs. epoch, where ``test-error-aver'' for epoch \( k \) is calculated by averaging the errors in epochs \(\max\{1, k-100\}\) to \(\min\{k+100, \#\text{epochs}\}\).
}
    \label{fig:S1:vs:S2:errors}
    % 		\vskip 0.02in
\end{figure}

% \begin{figure}%[htbp]%[h]%[htbp!]
%     \centering	
%         \begin{subfigure}[c]{0.2438\textwidth}
%     \centering            \includegraphics[width=0.985\textwidth]{figures/Strategy/errors1D_w500r25d6Aver.pdf}
%     \subcaption{MMNN (500,25,6).}
%     		% \label{subfig:f3}
%     \end{subfigure}
%     \hfill
% \begin{subfigure}[c]{0.2438\textwidth}
% 	\centering            \includegraphics[width=0.985\textwidth]{figures/Strategy/errors1D_w1000r50d8Aver.pdf}
% 	\subcaption{MMNN (1000,50,8).}
% 	% \label{subfig:f3}
% \end{subfigure}
% \hfill
% \begin{subfigure}[c]{0.2438\textwidth}
%     \centering            \includegraphics[width=0.985\textwidth]{figures/Strategy/errors1D_w500r25d6Comparison.pdf}
%     \subcaption{MMNN (500,25,6).}
%     		% \label{subfig:f3}
% \end{subfigure}\hfill
% \begin{subfigure}[c]{0.2438\textwidth}
% 	\centering            \includegraphics[width=0.985\textwidth]{figures/Strategy/errors1D_w1000r50d8Comparison.pdf}
% 	\subcaption{MMNN (1000,50,8).}
% 	% \label{subfig:f3}
% \end{subfigure}
% \caption{
% % First row: MMNN of size (500,25,6). Second row: MMNN of size (1000,50,8).
% % Comparison of training and test errors vs. epoch.
% % The ``test-error-aver'' for each epoch is calculated by averaging the errors from the preceding and following 100 epochs (if applicable).
% The ``test-error-aver'' for epoch \( k \) is calculated by averaging the errors from the preceding and following 100 epochs (if applicable), i.e., the average of the errors in epochs \(\max\{1, k-100\}\) to \(\min\{k+100, \#\text{epochs}\}\).
% % (a, b, c, d) First row: MMNN of size (500,25,6). (e, f, g, h) Second row: MMNN of size (1000,50,8).
% }
%     \label{fig:S1:vs:S2:errors}
%     % 		\vskip 0.02in
% \end{figure}

% \begin{figure}%[htbp]%[h]%[htbp!]
%     \centering	
%     \begin{subfigure}[c]{0.2438\textwidth}
%     \centering            \includegraphics[width=0.985\textwidth]{figures/Strategy/errors1D_w500r25d6P1.pdf}
%     \subcaption{S1.}
%     		% \label{subfig:f3}
%     \end{subfigure}
%     \hfill
%         \begin{subfigure}[c]{0.2438\textwidth}
%     \centering            \includegraphics[width=0.985\textwidth]{figures/Strategy/errors1D_w500r25d6P2.pdf}
%     \subcaption{S2.}
%     		% \label{subfig:f3}
%     \end{subfigure}
%     \hfill
%         \begin{subfigure}[c]{0.2438\textwidth}
%     \centering            \includegraphics[width=0.985\textwidth]{figures/Strategy/errors1D_w500r25d6Aver.pdf}
%     \subcaption{S1 vs. S2.}
%     		% \label{subfig:f3}
%     \end{subfigure}
%     \hfill
%         \begin{subfigure}[c]{0.2438\textwidth}
%     \centering            \includegraphics[width=0.985\textwidth]{figures/Strategy/errors1D_w500r25d6Comparison.pdf}
%     \subcaption{S1 vs. S2.}
%     		% \label{subfig:f3}
%     \end{subfigure}\\
%         \begin{subfigure}[c]{0.2438\textwidth}
% 	\centering            \includegraphics[width=0.985\textwidth]{figures/Strategy/errors1D_w1000r50d8P1.pdf}
% 	\subcaption{S1.}
% 	% \label{subfig:f3}
% \end{subfigure}
% \hfill
% \begin{subfigure}[c]{0.2438\textwidth}
% 	\centering            \includegraphics[width=0.985\textwidth]{figures/Strategy/errors1D_w1000r50d8P2.pdf}
% 	\subcaption{S2.}
% 	% \label{subfig:f3}
% \end{subfigure}
% \hfill
% \begin{subfigure}[c]{0.2438\textwidth}
% 	\centering            \includegraphics[width=0.985\textwidth]{figures/Strategy/errors1D_w1000r50d8Aver.pdf}
% 	\subcaption{S1 vs. S2.}
% 	% \label{subfig:f3}
% \end{subfigure}
% \hfill
% \begin{subfigure}[c]{0.2438\textwidth}
% 	\centering            \includegraphics[width=0.985\textwidth]{figures/Strategy/errors1D_w1000r50d8Comparison.pdf}
% 	\subcaption{S1 vs. S2.}
% 	% \label{subfig:f3}
% \end{subfigure}
% \caption{
% First row: MMNN of size (500,25,6). Second row: MMNN of size (1000,50,8).
% % Comparison of training and test errors vs. epoch.
% % The ``test-error-aver'' for each epoch is calculated by averaging the errors from the preceding and following 100 epochs (if applicable).
% The ``test-error-aver'' for epoch \( k \) is calculated by averaging the errors from the preceding and following 100 epochs (if applicable), i.e., the average of the errors in epochs \(\max\{1, k-100\}\) to \(\min\{k+100, \#\text{epochs}\}\).
% % (a, b, c, d) First row: MMNN of size (500,25,6). (e, f, g, h) Second row: MMNN of size (1000,50,8).
% }
%     \label{fig:S1:vs:S2:errors}
%     % 		\vskip 0.02in
% \end{figure}

% \begin{figure}%[htbp]%[h]%[htbp!]
%     \centering	
%     \foreach \ssss in {1,2}{
%     \foreach \iiidx in {400,900,1400,3000}
%     {
%     \begin{subfigure}[c]{0.231\textwidth}
%     \centering            \includegraphics[width=0.998055\textwidth]{figures/Strategy/epoch\iiidx_idx\ssss_w800r30d8P\ssss.png}
%     \subcaption{Epoch \iiidx\, (S\ssss).}
%     		% \label{subfig:f3}
%     \end{subfigure}
%     }\\
%     }
% \caption{Comparison of training process.}   \label{fig:S1:vs:S2:plots}
%     % 		\vskip 0.02in
% \end{figure}

% \begin{figure}%[htbp]%[h]
%     \centering
%     \begin{subfigure}[c]{0.242\textwidth}
%         \centering
%         \includegraphics[width=0.998055\textwidth]{figures/Strategy/epoch100_idx1_w500r25d6P1.pdf}
%         \subcaption{Epoch 100 (S1).}
%     \end{subfigure}
%     \hfill
%     \begin{subfigure}[c]{0.242\textwidth}
%         \centering
%         \includegraphics[width=0.998055\textwidth]{figures/Strategy/epoch500_idx1_w500r25d6P1.pdf}
%         \subcaption{Epoch 500 (S1).}
%     \end{subfigure}
%     \hfill
%     \begin{subfigure}[c]{0.242\textwidth}
%         \centering
%         \includegraphics[width=0.998055\textwidth]{figures/Strategy/epoch1000_idx1_w500r25d6P1.pdf}
%         \subcaption{Epoch 1000 (S1).}
%     \end{subfigure}
%     \hfill
%     \begin{subfigure}[c]{0.242\textwidth}
%         \centering
%         \includegraphics[width=0.998055\textwidth]{figures/Strategy/epoch20000_idx1_w500r25d6P1.pdf}
%         \subcaption{Epoch 20000 (S1).}
%     \end{subfigure}\\   
%     \begin{subfigure}[c]{0.242\textwidth}
%         \centering       \includegraphics[width=0.998055\textwidth]{figures/Strategy/epoch100_idx2_w500r25d6P2.pdf}
%         \subcaption{Epoch 100 (S2).}
%     \end{subfigure}
%     \hfill
%     \begin{subfigure}[c]{0.242\textwidth}
%         \centering
%         \includegraphics[width=0.998055\textwidth]{figures/Strategy/epoch500_idx2_w500r25d6P2.pdf}
%         \subcaption{Epoch 500 (S2).}
%     \end{subfigure}
%     \hfill
%     \begin{subfigure}[c]{0.242\textwidth}
%         \centering
%         \includegraphics[width=0.998055\textwidth]{figures/Strategy/epoch1000_idx2_w500r25d6P2.pdf}
%         \subcaption{Epoch 1000 (S2).}
%     \end{subfigure}
%     \hfill
%     \begin{subfigure}[c]{0.242\textwidth}
%         \centering
%         \includegraphics[width=0.998055\textwidth]{figures/Strategy/epoch20000_idx2_w500r25d6P2.pdf}
%         \subcaption{Epoch 20000 (S2).}
%     \end{subfigure}
    
%     \caption{Comparison of training process for MMNN of size (500,25,6).}   
%     \label{fig:S1:vs:S2:plots}
% \end{figure}



As illustrated in Table~\ref{tab:stategy:error:comparison} and Figure~\ref{fig:S1:vs:S2:errors},  our learning strategy S\(1\) is significantly more effective than strategy S\(2\) with comparable accuracy. 
There are two main advantages of S1. First, S1 requires training only about half the number of parameters compared to S2, which results in time savings. Second, S1 converges more quickly and performs significantly better when the training is not sufficient. 
We would like to note that in certain specific cases, S2 may outperform S1, particularly when the network size is relatively small and S2 is well-trained. This is expected since S\(2\) trains all parameters, whereas S\(1\) only trains a subset.
% Therefore, a combination of S\(2\) and S\(3\) may be more effective; specifically, using S\(2\) for the initial epochs and S\(3\) for the later epochs.
Based on our experience, S1 is more effective in practice, particularly for sufficiently large networks.
% Alternatively, one might consider a hybrid learning strategy.
%between S\(1\) and S\(2\). This would involve training the \(\bm{A}_i\)'s and \(\bm{c}_i\)'s, as well as parts of the \(\bm{W}_i\) and \(\bm{b}_i\), while fixing the remaining parts of \(\bm{W}_i\) and \(\bm{b}_i\) fixed.

% \red{It is important to note that our standard activation function, \ReLU, is not differentiable at $0$. For applications requiring a smooth (e.g., \(C^\infty\)) function, one might consider using \GELU\ \cite{2016arXiv160608415H} or \Swish\ \cite{2017arXiv171005941R}, which are smooth alternatives to \ReLU\  and generally yield results comparable to \ReLU. In contrast, other popular S-shaped activation functions like \Sigmoid\ and \Tanh\ have shown poor performance in our tests. 
% The choice of MMNN size largely depends on the complexity of the target functions. 
% For most target functions, there is a wide range of acceptable sizes for the MMNN to effectively learn, although larger models may require more training time.
% Several strategies may be applied: 1) The number of parameters may be comparable to that used in an FCNN for the same target function. 2) The width of the MMNN should significantly exceed its rank. 3) For highly oscillatory target functions that FCNNs struggle to learn effectively,
% a deeper ResMMNN with relatively large width and small rank should be considered. We believe that with appropriate hyperparameters, such as a gradually decreasing small learning rate over increasing epochs, the MMNN can be effectively trained.
% }


% In contrast, strategy S\(3\) exhibits the poorest performance,
% % ; as depicted in Figure~\ref{fig:S3},
% the training error does not decrease at all. This underscores the critical importance of choosing the appropriate initialization for the \(\bm{W}_i\)'s and \(\bm{b}_i\)'s for the random basis.


% \begin{figure}%[htbp]%[h]%[htbp!]
%     \centering	
%      \begin{subfigure}[c]{0.431\textwidth}
%     \centering            \includegraphics[width=0.998055\textwidth]{figures/Strategy/epoch20000_idx3_w500r25d6P3.pdf}
%     \subcaption{Epoch 20000 (S3).}
%     		% \label{subfig:f3}
%     \end{subfigure}
%     \hspace{25pt}
%     \begin{subfigure}[c]{0.405\textwidth}
%     \centering            \includegraphics[width=0.998055\textwidth]{figures/Strategy/errors1D_w500r25d6P3.pdf}
%     \subcaption{Errors vs. epoch.}
%     		% \label{subfig:f3}
%     \end{subfigure}
% \caption{Illustrations of the learning strategy S3 for MMNN of size (500,25,6).}   \label{fig:S3}
%     % 		\vskip 0.02in
% \end{figure}


% Connection to KAN, Fully connected ...


% \paragraph{Low-rank:}
% \red{

% \cite{2022arXiv220913569R}
% \href{https://arxiv.org/abs/2209.13569}{Exploring Low Rank Training of Deep Neural Networks}

% \cite{9157223}
% \href{https://openaccess.thecvf.com/content_CVPR_2020/papers/Idelbayev_Low-Rank_Compression_of_Neural_Nets_Learning_the_Rank_of_Each_CVPR_2020_paper.pdf}{Low-Rank Compression of Neural Nets: Learning the Rank of Each Layer}

% \cite{6638949}
% \href{https://vikas.sindhwani.org/lowRank.pdf}{Low-rank matrix factorization for Deep Neural Network training with high-dimensional output targets}
% }


% \section{How neural network approximate and model high-frequency}

% In this section, we explore the ability of neural networks to manage high-frequency components.
% We first demonstrated in Section~\ref{sec:shallow:low:pass:filter} that shallow neural networks function as low-pass filters, indicating their difficulty in processing high-frequency signals. Next, we explore the ability of deep networks to manage high-frequency components, analyzing this capability from a decomposition viewpoint. In Section~\ref{sec:decomposition:1D}, we initiate our discussion with the one-dimensional scenario. Subsequently, in Section~\ref{sec:decomposition:highD}, we extend our findings to cases involving higher dimensions.



% \subsection{MMNN versus fully connected network}

\section{MMNNs versus FCNNs}\label{sec:comparison}

% Now, let us compare MMNNs with fully connected neural networks (FCNNs).
% In Section~\ref{sec:MMNN:architecture}, we discussed the difference between MMNNs and FCNNs in terms of representation and learning strategy. Here we compare their numerical performance
% for oscillatory function approximaiton in Section~\ref{sec:1D2Dfunction} and PDE solving in Section~\ref{sec:PDE}
% .
In Section~\ref{sec:MMNN}, we outlined the distinctions between MMNNs and FCNNs regarding their representation and learning approaches. Here, we evaluate their numerical performance for approximating oscillatory functions in Section~\ref{sec:1D2Dfunction} and solving PDEs in Section~\ref{sec:PDE}.
To ensure a fair comparison, we use networks with a similar number of parameters and ensure that all networks have sufficient parameters to learn the target function. Typically, when training an FCNN, all parameters are optimized. For a thorough comparison, we will employ two learning strategies for MMNNs as detailed in Section~\ref{sec:MMNN:learning:stategy}: S1 and S2. S1 involves training approximately half the number of parameters of the MMNN, while S2 involves training all parameters.


\subsection{Oscillatory function approximation}\label{sec:1D2Dfunction}

We consider  a one-dimensional function \( f_1(x) = \cos(20\pi |x|^{1.4}) + 0.5\cos(12\pi |x|^{1.6}) \) and a two-dimensional function 
\begin{equation*}
	f_{2}(x_1,x_2)= \sum_{i=1}^2\sum_{j=1}^2 a_{ij} \sin(s b_i x_i+s c_{i,j} x_i x_j) \cos (s b_j x_j + s d_{i,j} x_i^2),
\end{equation*}
 where $s=2$ and
\begin{equation*}
	(a_{i,j})=	\begin{bmatrix*}
		0.3 & 0.2 \\ 0.2 & 0.3
	\end{bmatrix*}, \qquad 
	(b_{i})=	\begin{bmatrix*}
		2\pi \\ 4\pi
	\end{bmatrix*},
	\qquad
	(c_{i,j})=	\begin{bmatrix*}
		2\pi & 4\pi \\ 8\pi & 4\pi
	\end{bmatrix*}, \quad 
	(d_{i,j})=	\begin{bmatrix*}
		4\pi & 6\pi \\ 8\pi & 6\pi
	\end{bmatrix*}.
\end{equation*}
% \( f_2(r, \theta) \) defined in polar coordinates \((r, \theta)\) as follows:
% \begin{equation*}
%     f(r,\theta)=g(0.5+3\rho_2-3r)-g(0.2+8\rho_1-8r),
% \end{equation*}
% where
% \begin{equation*}
%     g(z)= \begin{cases}
%         0 & \text{if } z \le 0, \\
%         1 & \text{if } z \ge 1, \\
%         z & \text{otherwise}
%     \end{cases}
%     \quad \tn{and} \quad 
%     \begin{cases}
%     \rho_1 = 0.2 + 0.02\cos(8\theta),\\[3pt]
%         \rho_2 = 0.66 + 0.08\cos(\pi^2 \theta^2).
%     \end{cases}
% \end{equation*}
Refer to Figures~\ref{fig:vsFCNN:f1D} and~\ref{fig:vsFCNN:f2D} for illustrations of \( f_1 \) and \( f_2 \), respectively.

\begin{figure}%[htbp]%[h]%[htbp!]
\begin{minipage}[b]{0.302\linewidth}
    \centering	
    \includegraphics[width=0.925\textwidth]{figures/vsFCNN/vsFCNNf1D.pdf}
    % \subcaption{}
    %		\label{subfig:f3}
        \caption{ Plot of $f_1$.}
    	\label{fig:vsFCNN:f1D}
    % 		\vskip 0.02in
    \end{minipage}
    \hfill
\begin{minipage}[b]{0.57\linewidth}
    \centering	
    \begin{subfigure}[c]{0.46\textwidth}
    \centering            \includegraphics[height=0.80\textwidth]{figures/vsFCNN/vsFCNNf2D.png}
    % \subcaption{}
    %		\label{subfig:f3}
    \end{subfigure}
    \hfill
    \begin{subfigure}[c]{0.46\textwidth}
    \centering            \includegraphics[height=0.80\textwidth]{figures/vsFCNN/vsFCNNf3D.pdf}
    % \subcaption{}
    %		\label{subfig:f3}
    \end{subfigure}
    \caption{Plots of $f_2$.}
    	\label{fig:vsFCNN:f2D}
    % 		\vskip 0.02in
    \end{minipage}
\end{figure}


% For the 1D function, the FCNN has a width of 100 and a depth of 6, resulting in 50801 parameters. our MMNN has a compact size of $(321, 15, 6)$ with 50794 parameters.  For this test, a total of 100 data points are sampled on a uniform grid in \([-1,1]\) with a mini-batch size of 100 and a learning rate of \(0.001 \times 0.9^{\lfloor k/400 \rfloor}\) for epochs \(k=1, 2, \cdots, 20,000\).
% For 2D function,The FCNN has a width of 200 and a depth of 6, resulting in 201801 parameters.
% Our MMNN is of rather compact size $(653, 30, 6)$ with 201928 param. For this test, a total of $600^2$ data are sampled on a uniform grid in $[-1,1]^2$ with a mini-batch size of $1000$ and a learning rate of $0.001 \times 0.9^{\lfloor k/8 \rfloor}$ for epochs $k=1,2,\cdots,800$.

% For the 1D function, the FCNN has a width of 100 and a depth of 6, resulting in 50,801 parameters. Our MMNN has a compact size of $(321, 15, 6)$ with 50,794 parameters. For this test, a total of 1000 data points are sampled on a uniform grid in \([-1,1]\) with a mini-batch size of 100 and a learning rate of \(0.001 \times 0.9^{\lfloor k/400 \rfloor}\) for epochs \(k = 1, 2, \cdots, 20,000\).
% For the 2D function, the FCNN has a width of 200 and a depth of 6, resulting in 201,801 parameters. Our MMNN has a compact size of $(653, 30, 6)$ with 201,928 parameters. For this test, a total of \(600^2\) data points are sampled on a uniform grid in \([-1,1]^2\) with a mini-batch size of 1,000 and a learning rate of \(0.001 \times 0.9^{\lfloor k/20 \rfloor}\) for epochs \(k = 1, 2, \cdots, 1000\).
% Figure~\ref{fig:vsFCNN:errors} shows the error plot. As illustrated in the figure, MMNNs outperform FCNNs when both have the same depth and a comparable number of parameters. With a relatively small number of parameters, MMNNs show a significant performance advantage over FCNNs (See Figure~\ref{subfig:vsFCNN:errors:1D}). Even with a larger or sufficient number of parameters, MMNNs continue to perform better than FCNNs (See Figure~\ref{subfig:vsFCNN:errors:2D}).

% For the 1D function, the FCNN has a width of 120 and a depth of 6, resulting in 72961 parameters. Our MMNN is of size $(388, 18, 6)$ with 73035 parameters. For this test, a total of 1000 data points are sampled on a uniform grid in \([-1,1]\) with a mini-batch size of 100 and a learning rate of \(0.001 \times 0.9^{\lfloor k/400 \rfloor}\) for epochs \(k = 1, 2, \cdots, 20000\).
% For the 2D function, the FCNN has a width of 240 and a depth of 12, resulting in 637201 parameters. Our MMNN is of size of $(789, 36, 12)$ with 637120 parameters. For this test, a total of \(600^2\) data points are sampled on a uniform grid in \([-1,1]^2\) with a mini-batch size of 1,000 and a learning rate of \(0.001 \times 0.9^{\lfloor k/16 \rfloor}\) for epochs \(k = 1, 2, \cdots, 800\).

Large network sizes (see Table~\ref{tab:vsFCNN:error:comparison}) are selected to ensure that all networks possess sufficient parameters to learn the target functions.\footnote{FCNNs perform poorly if the network size is small. For a fair comparison, we choose relatively large network sizes for FCNNs and MMNNs, where both perform reasonably well.
}
% The MMNN size is characterized by (width, rank, depth), and the FCNN size is characterized by (width, depth). 
For training the one-dimensional function, we sample a total of 1000 data points on a uniform grid within \([-1,1]\), using a mini-batch size of 100 and a learning rate of \(0.001 \times 0.9^{\lfloor k/400 \rfloor}\) for epochs \(k = 1, 2, \cdots, 20000\).
For training the two-dimensional function, we sample a total of \(600^2\) data points on a uniform grid within \([-1,1]^2\), using a mini-batch size of 1000 and a learning rate of \(0.001 \times 0.9^{\lfloor k/16 \rfloor}\) for epochs \(k = 1, 2, \cdots, 800\).  
The Adam optimizer is employed for both functions.


% \begin{table}[htbp!]
% %	\vskip 0.15in
% 	\centering  
%  \setlength{\tabcolsep}{0.8em} 
%  \renewcommand{\arraystretch}{1.15}
% \caption{Comparison of test errors: averages of test errors in the final 20 epochs.}
% 	\label{tab:vsFCNN:error:comparison}
% 	\resizebox{0.99\textwidth}{!}{ 
% 		\begin{tabular}{ccccccccc} 
% 			\toprule% \toprule[1.2pt]  
% 			 target function &   network & {\#parameters (trained / all)}  &      {test error (MSE)} 
%     & test error (MAX) & training time \\
% 			 % \cmidrule{3-5}
% 			 % & & MSE & MAE & MAX\\
% 			\midrule
% 			 % \rowcolor{mygray}
% 			 $f_1$ & MMNN (S1)  &  35399 / 73035 & 
% $2.16 \times 10^{-6}$  &  $9.78 \times 10^{-3}$
%     & 23.3s / 1000 epochs \\	
% 			% \midrule
% 			$f_1$ & MMNN (S2)  &  73035 / 73035 & 
% $2.27 \times 10^{-6}$  &  $1.93 \times 10^{-2}$ 
%    & 27.4s / 1000 epochs\\			
% 			% \midrule
% 			$f_1$ & FCNN  & 72961 / 72961 & 
% $1.73 \times 10^{-4}$  &  $1.14 \times 10^{-1}$
%    & 22.3s / 1000 epochs 
%    \\
% 			\midrule
% 			 % \rowcolor{mygray}
% 			 $f_2$ & MMNN (S1)  &  313630 / 637120 & 
% $4.29 \times 10^{-6}$  &  $1.45 \times 10^{-2}$
%     & 30.3s / 10 epochs \\	
% 			% \midrule
% 			$f_2$ & MMNN (S2)  &  637120 / 637120 & 
% $5.89 \times 10^{-6}$  &  $5.92 \times 10^{-2}$
%    & 35.8s / 10 epochs\\			
% 			% \midrule
% 			$f_2$ & FCNN  & 637201 / 637201 & 
% $3.20 \times 10^{-5}$  &  $1.37 \times 10^{-1}$
%    & 29.3s / 10 epochs \\
% %			Training time (s) & 690  & 922 & 984 &  1328  \\
% 			\bottomrule% \bottomrule[1.2pt] 
% 		\end{tabular} 
% 	}%%% \resizebox
% % 	\vskip -0.1in
% \end{table} 

% \begin{table}[htbp!]
% %	\vskip 0.15in
% 	\centering  
%  \setlength{\tabcolsep}{0.568em} 
%  \renewcommand{\arraystretch}{1.15}
% \caption{Comparison of test errors.}
% 	\label{tab:vsFCNN:error:comparison}
% 	\resizebox{0.999\textwidth}{!}{ 
% 		\begin{tabular}{ccccccccc} 
% 			\toprule% \toprule[1.2pt]  
% 			 target function &   network &  (width, rank, depth) & {\#parameters (trained / all)}  &      {test error (MSE)} 
%     & test error (MAX) & training time \\
% 			 % \cmidrule{3-5}
% 			 % & & MSE & MAE & MAX\\
% 			\midrule
% 			 \rowcolor{mygray}
% 			 $f_1$ & MMNN1 (S1) & (388, 18, 6) &  35399 / 73035 & 
% $1.95 \times 10^{-6}$  &  $1.02 \times 10^{-2}$
%     & 23.3s / 1000 epochs \\	
% 			% \midrule
%    			$f_1$ & FCNN1-1 & (83, --,  6) & 35110 / 35110 & 
% $2.42 \times 10^{-4}$  &  $1.89 \times 10^{-1}$
%    & 19.5s / 1000 epochs
%    \\
% 		\rowcolor{mygray}	$f_1$ & MMNN1 (S2) & (388, 18, 6) &  73035 / 73035 & 
% $2.23 \times 10^{-6}$  &  $2.11 \times 10^{-2}$
%    & 27.4s / 1000 epochs\\			
% 			% \midrule
% 			$f_1$ & FCNN1-2 & (120, --,  6)  & 72961 / 72961 & 
% $1.73 \times 10^{-4}$  &  $1.12 \times 10^{-1}$
%    & 22.3s / 1000 epochs 
%    \\
% 			\midrule
% 			 \rowcolor{mygray}
% 			 $f_2$ & MMNN2 (S1) & (789, 36, 12) &  313630 / 637120 & 
% $4.40 \times 10^{-6}$  &  $1.33 \times 10^{-2}$
%     & 30.3s / 10 epochs \\
%     			$f_2$ & FCNN2-1 & (168, --, 12) & 312985 / 312985 & 
% $2.38 \times 10^{-4}$  &  $2.76 \times 10^{-1}$
%    & 26.7s / 10 epochs \\
% 			% \midrule
% 		\rowcolor{mygray}	$f_2$ & MMNN2 (S2) & (789, 36, 12)  &  637120 / 637120 & 
% $5.77 \times 10^{-6}$  &  $5.80 \times 10^{-2}$
%    & 35.8s / 10 epochs\\			
% 			% \midrule
% 			$f_2$ & FCNN2-2 & (240, --, 12) & 637201 / 637201 & 
% $3.23 \times 10^{-5}$  &  $1.37 \times 10^{-1}$
%    & 29.3s / 10 epochs \\
% %			Training time (s) & 690  & 922 & 984 &  1328  \\
% 			\bottomrule% \bottomrule[1.2pt] 
% 		\end{tabular} 
% 	}%%% \resizebox
% % 	\vskip -0.1in
% \end{table} 

\begin{table}[htbp!]
%	\vskip 0.15in
	\centering  
 \setlength{\tabcolsep}{0.568em} 
 \renewcommand{\arraystretch}{1.15}
\caption{
Comparison of test errors averaged over the last 100 epochs.
}
	\label{tab:vsFCNN:error:comparison}
	\resizebox{0.999\textwidth}{!}{ 
		\begin{tabular}{ccccccccc} 
			\toprule% \toprule[1.2pt]  
			 target function &   network &  (width, rank, depth) & {\#parameters (trained / all)}  &      {test error (MSE)} 
    & test error (MAX) & training time \\
			 % \cmidrule{3-5}
			 % & & MSE & MAE & MAX\\
			\midrule
			 \rowcolor{mygray}
			 $f_1$ & MMNN1 (S1) & (388, 18, 6) &  35399 / 73035 & 
$2.49 \times 10^{-6}$  &  $\bm{9.93 \times 10^{-3}}$
    & 23.3s / 1000 epochs \\	
			% \midrule
   			$f_1$ & FCNN1-1 & (83, --,  6) & 35110 / 35110 & 
$2.43 \times 10^{-4}$  &  $1.87 \times 10^{-1}$
   & 19.5s / 1000 epochs
   \\
		\rowcolor{mygray}	$f_1$ & MMNN1 (S2) & (388, 18, 6) &  73035 / 73035 & 
$\bm{2.05 \times 10^{-6}}$  &  $1.88 \times 10^{-2}$
   & 27.4s / 1000 epochs\\			
			% \midrule
			$f_1$ & FCNN1-2 & (120, --,  6)  & 72961 / 72961 & 
$1.73 \times 10^{-4}$  &  $1.14 \times 10^{-1}$
   & 22.3s / 1000 epochs 
   \\
			\midrule
			 \rowcolor{mygray}
			 $f_2$ & MMNN2 (S1) & (789, 36, 12) &  313630 / 637120 & 
$\bm{4.61 \times 10^{-6}}$  &  $\bm{1.55 \times 10^{-2}}$
    & 30.3s / 10 epochs \\
    			$f_2$ & FCNN2-1 & (168, --, 12) & 312985 / 312985 & 
$2.42 \times 10^{-4}$  &  $2.75 \times 10^{-1}$
   & 26.7s / 10 epochs \\
			% \midrule
		\rowcolor{mygray}	$f_2$ & MMNN2 (S2) & (789, 36, 12)  &  637120 / 637120 & 
$6.17 \times 10^{-6}$  &  $6.05 \times 10^{-2}$
   & 35.8s / 10 epochs\\			
			% \midrule
			$f_2$ & FCNN2-2 & (240, --, 12) & 637201 / 637201 & 
$3.28 \times 10^{-5}$  &  $1.39 \times 10^{-1}$
   & 29.3s / 10 epochs \\
%			Training time (s) & 690  & 922 & 984 &  1328  \\
			\bottomrule% \bottomrule[1.2pt] 
		\end{tabular} 
	}%%% \resizebox
% 	\vskip -0.1in
\end{table} 

% \begin{figure}%[htbp]%[h]%[htbp!]
%     \centering	
%      \begin{subfigure}[b]{0.40\textwidth}
%     \centering            \includegraphics[width=0.808\textwidth]{figures/vsFCNN/1Derrors_vsFCNN.pdf}
%     \subcaption{$f_1$.}
%     		% \label{subfig:vsFCNN:errors:1D}
%     \end{subfigure}
%     \hspace{16pt}
%     \begin{subfigure}[b]{0.40\textwidth}
%     \centering            \includegraphics[width=0.808\textwidth]{figures/vsFCNN/2Derrors_vsFCNN.pdf}
%     \subcaption{$f_2$.}
%     		% \label{subfig:vsFCNN:errors:1D}
%     \end{subfigure}
%     \caption{
%     The y-axes is $\log_{10}(\text{test-error-aver})$, where the ``test-error-aver" is the same as Figure~\ref{fig:S1:vs:S2:errors}. 
%     % for epoch \( k \) is calculated by averaging the errors from the preceding and following 100 epochs (if applicable). Specifically, it is the average of the errors in epochs \(\max\{1, k-100\}\) to \(\min\{k+100, \#\text{epochs}\}\).
%     }
% % \caption{The first and second rows correspond to $f_1$ and $f_2$, respectively. First column: Log plot of test errors vs. epochs. Second column: MMNN with the learning strategy proposed in Section~\ref{sec:MMNN:learning:stategy}, where $\bmA_i$'s and $\bmc_i$'s are trained and $\bmW_i$'s and $\bmb_i$'s are fixed. Third column: MMNN (all parameters trained). Fourth column: FCNN.}
%   \label{fig:vsFCNN:errors}
%     % 		\vskip 0.02in
% \end{figure}


% \begin{figure}%[htbp]%[h]%[htbp!]
%     \centering	
%     \begin{subfigure}[b]{0.2311\textwidth}
%     \centering            \includegraphics[width=0.985\textwidth]{figures/vsFCNN/1DvsFCNN_Diff_w388r18d6kind3.pdf}
%     \subcaption{MMNN1 (S1).}
%     % 		\label{subfig:vsFCNN:errors:2D}
%     \end{subfigure}
%         \hfill
%              \begin{subfigure}[b]{0.2311\textwidth}
%     \centering            \includegraphics[width=0.985\textwidth]{figures/vsFCNN/1DvsFCNN_Diff_w83r0d6kind0.pdf}
%     \subcaption{FCNN1-1.}
%     		% \label{subfig:vsFCNN:errors:1D}
%     \end{subfigure}
%     \hfill
%          \begin{subfigure}[b]{0.2311\textwidth}
%     \centering            \includegraphics[width=0.985\textwidth]{figures/vsFCNN/1DvsFCNN_Diff_w388r18d6kind2.pdf}
%     \subcaption{MMNN1 (S2).}
%     % 		\label{subfig:vsFCNN:errors:1D}
%     \end{subfigure}
%     \hfill
%     \begin{subfigure}[b]{0.2311\textwidth}
%     \centering            \includegraphics[width=0.985\textwidth]{figures/vsFCNN/1DvsFCNN_Diff_w120r0d6kind1.pdf}
%     \subcaption{FCNN1-2.}
%     % 		\label{subfig:vsFCNN:errors:2D}
%     \end{subfigure}\\[6pt]
%     \begin{subfigure}[b]{0.2311\textwidth}
%     \centering            \includegraphics[width=0.998055\textwidth]{figures/vsFCNN/2DvsFCNN_Diff_w789r36d12kind3.png}
%     \subcaption{MMNN2 (S1).}
%     % 		\label{subfig:vsFCNN:errors:2D}
%     \end{subfigure}
%         \hfill
%               \begin{subfigure}[b]{0.2311\textwidth}
%     \centering            \includegraphics[width=0.985\textwidth]{figures/vsFCNN/2DvsFCNN_Diff_w168r0d12kind0.png}
%     \subcaption{FCNN2-1.}
%     		% \label{subfig:vsFCNN:errors:1D}
%     \end{subfigure}
%     \hfill
%          \begin{subfigure}[b]{0.2311\textwidth}
%     \centering            \includegraphics[width=0.998055\textwidth]{figures/vsFCNN/2DvsFCNN_Diff_w789r36d12kind2.png}
%     \subcaption{MMNN2 (S2).}
%     % 		\label{subfig:vsFCNN:errors:1D}
%     \end{subfigure}
%     \hfill
%     \begin{subfigure}[b]{0.2311\textwidth}
%     \centering            \includegraphics[width=0.998055\textwidth]{figures/vsFCNN/2DvsFCNN_Diff_w240r0d12kind1.png}
%     \subcaption{FCNN2-2.}
%     % 		\label{subfig:vsFCNN:errors:2D}
%     \end{subfigure}
%     \caption{The differences between the learned network and $f_2$. }
%   \label{fig:vsFCNN:NN:fun:diff2}
%     % 		\vskip 0.02in
% \end{figure}

\begin{figure}%[htbp]%[h]%[htbp!]
    \centering	
     \begin{subfigure}[b]{0.442\textwidth}
    \centering            \includegraphics[width=0.9015\textwidth]{figures/vsFCNN/1Derrors_vsFCNN.pdf}
    \subcaption{$f_1$.}
    		% \label{subfig:vsFCNN:errors:1D}
    \end{subfigure}
    \hspace{14pt}
    \begin{subfigure}[b]{0.442\textwidth}
    \centering            \includegraphics[width=0.9015\textwidth]{figures/vsFCNN/2Derrors_vsFCNN.pdf}
    \subcaption{$f_2$.}
    		% \label{subfig:vsFCNN:errors:1D}
    \end{subfigure}
    \\
    \begin{subfigure}[b]{0.24311\textwidth}
    \centering            \includegraphics[width=0.985\textwidth]{figures/vsFCNN/1DvsFCNN_Diff_w388r18d6kind3.pdf}
    \subcaption{MMNN1 (S1).}
    % 		\label{subfig:vsFCNN:errors:2D}
    \end{subfigure}
        \hfill
             \begin{subfigure}[b]{0.24311\textwidth}
    \centering            \includegraphics[width=0.985\textwidth]{figures/vsFCNN/1DvsFCNN_Diff_w83r0d6kind0.pdf}
    \subcaption{FCNN1-1.}
    		% \label{subfig:vsFCNN:errors:1D}
    \end{subfigure}
    \hfill
         \begin{subfigure}[b]{0.24311\textwidth}
    \centering            \includegraphics[width=0.985\textwidth]{figures/vsFCNN/1DvsFCNN_Diff_w388r18d6kind2.pdf}
    \subcaption{MMNN1 (S2).}
    % 		\label{subfig:vsFCNN:errors:1D}
    \end{subfigure}
    \hfill
    \begin{subfigure}[b]{0.24311\textwidth}
    \centering            \includegraphics[width=0.985\textwidth]{figures/vsFCNN/1DvsFCNN_Diff_w120r0d6kind1.pdf}
    \subcaption{FCNN1-2.}
    % 		\label{subfig:vsFCNN:errors:2D}
    \end{subfigure}\hfill
 \\[6pt]
    \begin{subfigure}[b]{0.24311\textwidth}
    \centering            \includegraphics[width=0.998055\textwidth]{figures/vsFCNN/2DvsFCNN_Diff_w789r36d12kind3.png}
    \subcaption{MMNN2 (S1).}
    % 		\label{subfig:vsFCNN:errors:2D}
    \end{subfigure}
        \hfill
              \begin{subfigure}[b]{0.24311\textwidth}
    \centering            \includegraphics[width=0.985\textwidth]{figures/vsFCNN/2DvsFCNN_Diff_w168r0d12kind0.png}
    \subcaption{FCNN2-1.}
    		% \label{subfig:vsFCNN:errors:1D}
    \end{subfigure}
    \hfill
         \begin{subfigure}[b]{0.24311\textwidth}
    \centering            \includegraphics[width=0.998055\textwidth]{figures/vsFCNN/2DvsFCNN_Diff_w789r36d12kind2.png}
    \subcaption{MMNN2 (S2).}
    % 		\label{subfig:vsFCNN:errors:1D}
    \end{subfigure}
    \hfill
    \begin{subfigure}[b]{0.24311\textwidth}
    \centering            \includegraphics[width=0.998055\textwidth]{figures/vsFCNN/2DvsFCNN_Diff_w240r0d12kind1.png}
    \subcaption{FCNN2-2.}
    % 		\label{subfig:vsFCNN:errors:2D}
    \end{subfigure}
    \caption{First row: base-10 logarithm of ``test-error-aver'' vs. epoch, where ``test-error-aver'' for epoch \( k \) is calculated by averaging the errors in epochs \(\max\{1, k-100\}\) to \(\min\{k+100, \#\text{epochs}\}\).
    All errors shown on the $y$-axis are in base-10 logarithmic scale.
    Second row: differences between learned networks and $f_1$. 
    Third row: differences between learned networks and $f_2$.}
  \label{fig:vsFCNN:errors:and:NN:fun:diff}
    % 		\vskip 0.02in
\end{figure}

As illustrated in 
Table~\ref{tab:vsFCNN:error:comparison} and Figure~\ref{fig:vsFCNN:errors:and:NN:fun:diff}, MMNNs outperform FCNNs when both have the same depth and a comparable number of parameters, particularly for relatively oscillatory target functions. 
Moreover, as indicated in Table~\ref{tab:vsFCNN:error:comparison}, the training time for MMNN (S1) is similar to that of FCNN, while MMNN (S2) takes a bit more time.
We remark that the primary advantage of MMNNs lies in capturing high-frequency components.
As we can see from Figure~\ref{fig:vsFCNN:errors:and:NN:fun:diff}, the differences between network approximations and the corresponding target functions show that FCNNs approximate high-frequency parts of the target functions poorly. In contrast, the approximation errors for MMNNs, especially with the S1 learning strategy, are more evenly distributed across the entire domain, indicating their effectiveness in capturing high-frequency components. The Adam optimizer \cite{DBLP:journals/corr/KingmaB14} is applied throughout the training process.


% The differences between the learned networks and the target functions show that FCNNs struggle to capture the high-frequency components of the target functions. In contrast, MMNNs, especially with the S1 learning strategy, excel at capturing these high-frequency components.
% With a relatively small number of parameters, MMNNs show a significant performance advantage over FCNNs (See Figure~\ref{subfig:vsFCNN:errors:1D}). Even with a larger or sufficient number of parameters, MMNNs continue to perform better than FCNNs (See Figure~\ref{subfig:vsFCNN:errors:2D}).


% \begin{figure}%[htbp]%[h]%[htbp!]
%     \centering	
%      \begin{subfigure}[c]{0.431\textwidth}
%     \centering            \includegraphics[width=0.998055\textwidth]{figures/vsFCNN/1Derrors_vsFCNN.pdf}
%     \subcaption{$f_1$.}
%     		\label{subfig:vsFCNN:errors:1D}
%     \end{subfigure}
%     \hspace{25pt}
%     \begin{subfigure}[c]{0.431\textwidth}
%     \centering            \includegraphics[width=0.998055\textwidth]{figures/vsFCNN/2Derrors_vsFCNN.pdf}
%     \subcaption{$f_2$.}
%     		\label{subfig:vsFCNN:errors:2D}
%     \end{subfigure}
% \caption{Log plot of test errors vs. epoch.}
%   \label{fig:vsFCNN:errors}
%     % 		\vskip 0.02in
% \end{figure}



% \subsection{PDE Solving}
\subsection{Solving partial differential equations}
\label{sec:PDE}


Next, we compare the performance of MMNNs and FCNNs for solving partial differential equations (PDEs).  
We consider a classical example: the two-dimensional Poisson equation with zero Dirichlet boundary conditions:
\[
-\Delta u(x, y) = f(x, y), \quad (x, y) \in (-1,1)^2,\quad u|_{\partial \Omega} = 0,
\]
where the source term is given by
% \[
% f(x,y) = -112\pi^2\sin(7\pi x)\sin(8\pi y) - 108\pi^2\sin(6\pi x)\sin(9\pi y).
% \]
\[
f(x,y) = -113\pi^2\sin(7\pi x)\sin(8\pi y) - 117\pi^2\sin(6\pi x)\sin(9\pi y).
\]
It is easy to verify that the exact solution is
\[
u(x,y) = \sin(7\pi x)\sin(8\pi y) + \sin(6\pi x)\sin(9\pi y).
\]
% where the source term is given by
% \[
% f(x,y) = -72\pi^2\sin(6\pi x)\sin(6\pi y) - 50\pi^2\sin(5\pi x)\sin(5\pi y).
% \]
% It is easy to verify that the exact solution is
% \[
% u(x,y) = \sin(6\pi x)\sin(6\pi y) + \sin(5\pi x)\sin(5\pi y).
% \]


 We approximate the solution \( u(x, y) \) by a neural network \( u_\bmtheta(x, y) \), where \( \bmtheta \) denotes the network parameters and solve the Poisson equation using the Physics-Informed Neural Network (PINN) \cite{RAISSI2019686} formulation.
The PDE residual is defined as \(\mathcal{R}(x, y) = -\Delta u_\bmtheta(x, y) - f(x, y)\), where \(\Delta u_\bmtheta = \frac{\partial^2 u_\bmtheta}{\partial x^2} + \frac{\partial^2 u_\bmtheta}{\partial y^2}\), and the neural network is trained to minimize the loss function \(\mathcal{L}_{\text{total}} = \lambda_{\text{PDE}}\mathcal{L}_{\text{PDE}} + \lambda_{\text{BC}}\mathcal{L}_{\text{BC}}\), combining the PDE loss \(\mathcal{L}_{\text{PDE}} = \frac{1}{N_r} \sum_{i=1}^{N_r} \mathcal{R}(x_i, y_i)^2\) for collocation points \((x_i, y_i) \in (-1,1)^2\) and the boundary loss \(\mathcal{L}_{\text{BC}} = \frac{1}{N_b} \sum_{j=1}^{N_b} u_\bmtheta(x_j, y_j)^2\) for boundary points \((x_j, y_j) \in \partial \Omega\).


To compare FCNNs and MMNNs in terms of representation accuracy, we avoid soft boundary condition enforcement by including the penalty term \(\lambda_{\text{BC}}\mathcal{L}_{\text{BC}}\) in the loss, which requires careful tuning of \(\lambda_{\text{BC}}\). 
We adopt a hard-constraint formulation that guarantees the boundary condition is satisfied exactly. Specifically, we define the network output as
\[
u_{\bm{\theta}}(x, y) = h_{\bm{\theta}}(x, y) \cdot \cos\left(\frac{\pi x}{2}\right)\cos\left(\frac{\pi y}{2}\right),
\]
where \( h_{\bm{\theta}}(x, y) \) is modeled by either an FCNN or an MMNN. This construction ensures that \( u_{\bm{\theta}}(x, y) = 0 \) on the boundary \( \partial \Omega \) for all values of \( \bm{\theta} \).


% We select an MMNN of size (301, 16, 6), another MMNN of size (503, 20, 6), and an FCNN of size (100, --, 6), denoted as MMNN1, MMNN2, and FCNN, respectively, for simplicity. We sample \(100^2\) data points on a uniform grid in \([-1,1]^2\) for $f(x,y)$, using a mini-batch size of 2000. $\lambda_{\tn{PDE}}=0.001$.
% As for the learning rate of $\lfloor k/200\rfloor / 800$
% if $k<16000$ and 
% \(0.001 \times 0.9^{\lfloor (k-16000)/1600 \rfloor}\) if $k\ge 1600$ for epochs \(k = 1, 2, \dots, 160000\) .
We select an MMNN of size (301, 16, 6), another MMNN of size (503, 20, 6), and an FCNN of size (100, --, 6), denoted as MMNN1, MMNN2, and FCNN, respectively, for simplicity. 
Since \ReLU{} is not differentiable, we use the sine function as the activation function.
We sample \(100^2\) data points for \(f(x,y)\) on a uniform grid in \((-1,1)^2\), using a mini-batch size of 2000 and setting \(\lambda_{\text{PDE}} = 0.001\). 
We adopt the Adam optimizer, with the learning rate set to 
\(\lfloor k/200 \rfloor / 800\) for \(k < 16000\), and to 
\(0.001 \times 0.9^{\lfloor (k - 16000)/1600 \rfloor}\) for 
\(k \geq 16000\), where \(k = 1, 2, \dots, 160000\) denotes the training epoch.
We note that for \( k < 16000 \), we use an increasing learning rate to facilitate warm-up (see, e.g., \cite{NEURIPS2020_288cd256,NEURIPS2024_ca98452d}), which enhances training performance.



Our experiments reveal that initial parameters significantly impact training, particularly for FCNNs. To ensure experimental reliability, we repeated the experiments with 16 different seeds. As demonstrated in Figures~\ref{fig:PDE:error:comparison}, \ref{fig:PDE:fun:diff}, and Table~\ref{tab:PDE:error:comparison}, both MMNNs surpass the FCNN in solving PDEs with PINNs, even with comparable depth and total (or training) parameters.
We note that using the same seed may yield different outcomes across various code environments. Our tests consistently show that MMNNs always succeed, whereas FCNNs fail with high probability (increasing FCNN width may improve the likelihood of success). 


\begin{figure}%[htbp]%[htbp!]
    \centering	
            \begin{subfigure}[c]{0.3253012\textwidth}
    \centering            \includegraphics[width=0.975859598055\textwidth]{figures/PDE/PDE_error_train.pdf}
    \subcaption{Training error.}
    		% \label{subfig:f3}
    \end{subfigure}
    \hfill
    \begin{subfigure}[c]{0.3253012\textwidth}
    \centering            \includegraphics[width=0.975859598055\textwidth]{figures/PDE/PDE_error_train.pdf}
    \subcaption{Test error (MSE).}
    		% \label{subfig:f3}
    \end{subfigure}
    \hfill 
 \begin{subfigure}[c]{0.3253012\textwidth}
    \centering            \includegraphics[width=0.975859598055\textwidth]{figures/PDE/PDE_error_testmax.pdf}
    \subcaption{Test error (MAX).}
    		% \label{subfig:f3}
    \end{subfigure}
\caption{Average errors across 16 seeds versus epoch. The training error corresponds to the PDE loss, while the test error (MSE or MAX) quantifies the difference between the learned network and the true solution. All errors shown on the $y$-axis are in base-10 logarithmic scale.}
    % 		\vskip 0.02in
\label{fig:PDE:error:comparison}
\end{figure}


\begin{figure}%[htbp]%[htbp!]
    \centering	
    \begin{subfigure}[b]{0.24311\textwidth}
    \centering            \includegraphics[width=0.998055\textwidth]{figures/PDE/PDE_diff_model_true.png}
    \subcaption{Truth.}
    % 		\label{subfig:vsFCNN:errors:2D}
    \end{subfigure}
        \hfill
              \begin{subfigure}[b]{0.24311\textwidth}
    \centering            \includegraphics[width=0.985\textwidth]{figures/PDE/PDE_diff_model2.png}
    \subcaption{MMNN1.}
    		% \label{subfig:vsFCNN:errors:1D}
    \end{subfigure}
    \hfill
         \begin{subfigure}[b]{0.24311\textwidth}
    \centering            \includegraphics[width=0.998055\textwidth]{figures/PDE/PDE_diff_model1.png}
    \subcaption{MMNN2.}
    % 		\label{subfig:vsFCNN:errors:1D}
    \end{subfigure}
    \hfill
    \begin{subfigure}[b]{0.24311\textwidth}
    \centering            \includegraphics[width=0.998055\textwidth]{figures/PDE/PDE_diff_model3.png}
    \subcaption{FCNN.}
    % 		\label{subfig:vsFCNN:errors:2D}
    \end{subfigure}
\caption{Comparison of three networks: (a) true solution; (b, c, d) differences between the true solution and predictions from learned networks. For each network, we select the best trained model from 16 seeds.}
  \label{fig:PDE:fun:diff}
    % 		\vskip 0.02in
\end{figure}



\begin{table}[h]%[htbp!]
	\centering  
 \setlength{\tabcolsep}{0.868em} 
 \renewcommand{\arraystretch}{1.15}
\caption{Comparison of test errors. }
	\label{tab:PDE:error:comparison}
	\resizebox{0.95\textwidth}{!}{ 
		\begin{tabular}{ccccccccc} 
			\toprule% \toprule[1.2pt]  
            % \multicolumn{2}{c}
            % {
            % target function
            % }
            % & \multicolumn{2}{c}{$f_1:[-1,1]\to \R$}
            % & \multicolumn{2}{c}{$f_2:[-1,1]^2\to \R$}
            % % & \multicolumn{2}{c}{$f_3:[-1,1]^2\to \R$}
            % & \multicolumn{2}{c}{$f_3:[-1,1]^3\to \R$}\\        
            % % \cmidrule(lr){1-2} 
            % % \cmidrule(lr){3-4}
            % % \cmidrule(lr){5-6}
            % % \cmidrule(lr){7-8}
            % \cmidrule(lr){2-3}
            % \cmidrule(lr){4-5}
            % \cmidrule(lr){6-7}
            % % \cmidrule(lr){8-9}
         % \multicolumn{2}{c}
            {\#parameters (trained / all)}
            &\multicolumn{2}{c}{24462 / \textbf{50950}} &\multicolumn{2}{c}{\textbf{50904} / 105228} &\multicolumn{2}{c}{\textbf{50901} / \textbf{50901}}   \\
            \cmidrule(lr){2-3}
            \cmidrule(lr){4-5}
            \cmidrule(lr){6-7}
            % \cmidrule(lr){1-2} 
            % \cmidrule(lr){3-4}
            % \cmidrule(lr){5-6}
            % \cmidrule(lr){7-8}
        \rowcolor{mygray}   
        % \multicolumn{2}{c}{network} 
        % &
        &\multicolumn{2}{c}{MMNN1 of size (301, 16, 6)} &\multicolumn{2}{c}{MMNN2 of size (503, 20, 6)} &\multicolumn{2}{c}{FCNN of size (100,  --, 6)}
        
        % &\multicolumn{2}{c}{ResMMNN of size (1024,32,10)}
        \\
           %             \cmidrule(lr){3-4}
           %  \cmidrule(lr){5-6}
           %  \cmidrule(lr){7-8}
           % & &\multicolumn{2}{c}{MMNN1}
           %  &\multicolumn{2}{c}{MMNN2}
           %  &\multicolumn{2}{c}{FCNN1}
           %  \\
            % \cmidrule(lr){3-4}
            % \cmidrule(lr){5-6}
            % \cmidrule(lr){7-8}
            \cmidrule(lr){2-3}
            \cmidrule(lr){4-5}
            \cmidrule(lr){6-7}
            % \cmidrule(lr){8-9}
			 %    target  &  activation &        TE (MSE) 
    % & TE (MAX) &        TE (MSE) 
    % & TE (MAX)
    % &        TE (MSE) 
    % & TE (MAX)\\
    % \multicolumn{2}{c}
        {seed for randomness} &         MSE 
    &  MAX &    MSE 
    &  MAX &    MSE 
    &  MAX 
    % &    MSE 
    % &  MAX
    \\
  \midrule  



$0$   &  $ 2.60 \times 10^{-6} $  &  $ 6.22 \times 10^{-3} $  &  $ 1.77 \times 10^{-8} $  &  $ 3.38 \times 10^{-3} $  &  $ 4.95 \times 10^{-1} $  &  $ 1.94 \times 10^{0} $ 
 \\ 

\rowcolor{mygray}$1$   &  $ 2.18 \times 10^{-7} $  &  $ 2.08 \times 10^{-3} $  &  $ 9.32 \times 10^{-9} $  &  $ 6.76 \times 10^{-4} $  &  $ \bm{7.95 \times 10^{-7}} $  &  $ 3.35 \times 10^{-3} $ 
 \\ 

 $2$   &  $ 1.18 \times 10^{-6} $  &  $ 4.18 \times 10^{-3} $  &  $ 2.38 \times 10^{-8} $  &  $ 1.55 \times 10^{-3} $  &  $ 5.01 \times 10^{-1} $  &  $ 1.94 \times 10^{0} $ 
 \\ 

\rowcolor{mygray}$3$   &  $ 3.19 \times 10^{-6} $  &  $ 5.57 \times 10^{-3} $  &  $ 6.68 \times 10^{-9} $  &  $ 1.19 \times 10^{-3} $  &  $ 5.03 \times 10^{-1} $  &  $ 1.94 \times 10^{0} $ 
 \\ 

 $4$   &  $ 5.34 \times 10^{-7} $  &  $ 3.12 \times 10^{-3} $  &  $ 2.35 \times 10^{-8} $  &  $ 2.89 \times 10^{-3} $  &  $ 4.99 \times 10^{-1} $  &  $ 1.94 \times 10^{0} $ 
 \\ 

\rowcolor{mygray}$5$   &  $ 1.76 \times 10^{-6} $  &  $ 5.15 \times 10^{-3} $  &  $ 1.78 \times 10^{-8} $  &  $ 2.51 \times 10^{-3} $  &  $ 4.96 \times 10^{-1} $  &  $ 1.94 \times 10^{0} $ 
 \\ 

 $6$   &  $ 4.92 \times 10^{-7} $  &  $ 2.61 \times 10^{-3} $  &  $ 9.68 \times 10^{-9} $  &  $ 4.65 \times 10^{-3} $  &  $ 5.04 \times 10^{-1} $  &  $ 1.94 \times 10^{0} $ 
 \\ 

\rowcolor{mygray}$7$   &  $ 1.97 \times 10^{-6} $  &  $ 5.34 \times 10^{-3} $  &  $ 3.87 \times 10^{-8} $  &  $ 3.37 \times 10^{-3} $  &  $ 5.01 \times 10^{-1} $  &  $ 1.94 \times 10^{0} $ 
 \\ 

 $8$   &  $ 4.98 \times 10^{-7} $  &  $ 2.44 \times 10^{-3} $  &  $ 4.04 \times 10^{-8} $  &  $ 9.83 \times 10^{-4} $  &  $ 8.20 \times 10^{-7} $  &  $ \bm{3.22 \times 10^{-3}} $ 
 \\ 

\rowcolor{mygray}$9$   &  $ 2.57 \times 10^{-6} $  &  $ 4.10 \times 10^{-3} $  &  $ \bm{1.69 \times 10^{-9}} $  &  $ \bm{6.66 \times 10^{-4}} $  &  $ 5.00 \times 10^{-1} $  &  $ 1.94 \times 10^{0} $ 
 \\ 

 $10$   &  $ 2.47 \times 10^{-6} $  &  $ 6.15 \times 10^{-3} $  &  $ 2.72 \times 10^{-8} $  &  $ 3.88 \times 10^{-3} $  &  $ 5.02 \times 10^{-1} $  &  $ 1.94 \times 10^{0} $ 
 \\ 

\rowcolor{mygray}$11$   &  $ 1.03 \times 10^{-6} $  &  $ 3.86 \times 10^{-3} $  &  $ 2.02 \times 10^{-8} $  &  $ 1.31 \times 10^{-3} $  &  $ 5.03 \times 10^{-1} $  &  $ 1.94 \times 10^{0} $ 
 \\ 

 $12$   &  $ 1.61 \times 10^{-6} $  &  $ 5.13 \times 10^{-3} $  &  $ 4.38 \times 10^{-9} $  &  $ 1.16 \times 10^{-3} $  &  $ 1.60 \times 10^{-5} $  &  $ 1.12 \times 10^{-2} $ 
 \\ 

\rowcolor{mygray}$13$   &  $ \bm{9.35 \times 10^{-8}} $  &  $ \bm{1.52 \times 10^{-3}} $  &  $ 3.85 \times 10^{-8} $  &  $ 1.07 \times 10^{-2} $  &  $ 2.86 \times 10^{-2} $  &  $ 4.39 \times 10^{-1} $ 
 \\ 

 $14$   &  $ 7.42 \times 10^{-7} $  &  $ 3.41 \times 10^{-3} $  &  $ 1.51 \times 10^{-8} $  &  $ 1.77 \times 10^{-3} $  &  $ 5.01 \times 10^{-1} $  &  $ 1.95 \times 10^{0} $ 
 \\ 

\rowcolor{mygray}$15$   &  $ 6.58 \times 10^{-6} $  &  $ 6.25 \times 10^{-3} $  &  $ 2.74 \times 10^{-8} $  &  $ 1.23 \times 10^{-3} $  &  $ 6.85 \times 10^{-6} $  &  $ 7.97 \times 10^{-3} $ 
 \\ 



 \bottomrule% \bottomrule[1.2pt] 
		\end{tabular} 
	}%%% \resizebox
% 	\vskip -0.1in
\end{table} 





\section{Multi-component and multi-layer decomposition}
\label{sec:decomposition}

% Although a one-hidden-layer neural network is a low-pass filter that can not represent or learn high-frequency features effectively \cite{ZZZZ-23}, we use mathematical construction to show that MMNNs, which are composed of one-hidden-layer neural networks, can overcome this difficulty by decomposition of the complexity through components and/or depth.
It has been shown in \cite{ZZZZ-23} that a one-hidden-layer neural network acts as a low-pass filter and cannot effectively represent or learn high-frequency features. Using mathematical construction, we demonstrate that MMNNs, which are composed of one-hidden-layer neural networks, can overcome this difficulty by decomposing complexity through their components and/or depth.
We emphasize that the decomposition is highly non-unique. Our construction is ``man-made" which can be different from the one by computer through an optimization (learning) process.
%In this section, we examine how neural networks handle high-frequency components. We previously demonstrated in Section~\ref{sec:shallow:low:pass:filter} that shallow neural networks act as low-pass filters, struggling with high-frequency signals. We then investigate the proficiency of deep networks in managing high-frequency elements, analyzing their capabilities through a decomposition approach. 
Our discussion begins with one-dimensional construction in Section~\ref{sec:decomposition:1D} and later extends to higher dimensions in Section~\ref{sec:decomposition:highD}.


%\subsection{Challenges of Shallow Networks in Handling High-Frequency}
%\label{sec:shallow:low:pass:filter}

%In our previous work \cite{ZZZZ-23}, we demonstrated that shallow neural networks function as low-pass filters, indicating their difficulty in processing high-frequency signals. In this section, we explore the ability of deep networks to manage high-frequency components, analyzing this capability from a decomposition viewpoint. In Section~\ref{sec:decomposition:1D}, we initiate our discussion with the one-dimensional scenario. Subsequently, in Section~\ref{sec:decomposition:highD}, we extend our findings to cases involving higher dimensions.

% \subsection{Managing High Frequency Through Constructed Balanced Decomposition}
\subsection{One dimensional construction}
% \subsection{Univariant function}
\label{sec:decomposition:1D}

We begin with a two-component decomposition in the one-dimensional as both an illustration and an example in Section~\ref{sec:1D:two:component:decom}. Later in Section~\ref{sec:1D:general:multi:component:decom}, we introduce the general multi-component decomposition. Finally in Section~\ref{sec:1D:decom:examples}, we use concrete examples for demonstration.

% \subsubsection{One Dimensional Case}
% \label{sec:decomposition:1D}

% \subsubsection{Multi-layer decomposition}
% We use the following example to show that any Fourier mode, no matter how high the frequency is, can be decomposed and approximated well using a multi-layer network even with one component per layer. Consider the target function
% % $f(x)=\cos(2n\pi x)$.
% % The function is redefined as 
% $f(x)=\big(\cos(2n\pi x)-1\big)/2$, ensuring $f(0)=0$. The functions $f_1$ and $f_2$ are defined piecewise as follows:
% \begin{equation*}
% f_1=\begin{cases}
% \big(\cos(2n\pi x)-1\big)/2 & \text{for}\ x\in [-1,0),\\
% x & \text{for}\ x\in [0,1],
% \end{cases}
% \end{equation*}
% and
% \begin{equation*}
% f_2=\begin{cases}
% x & \text{for}\ x\in [-1,0),\\
% \big(\cos(2n\pi x)-1\big)/2 & \text{for}\ x\in [0,1].
% \end{cases}
% \end{equation*}
% Consequently, it can be shown that
% \begin{equation*}
% f(x)=f_2\circ f_1(x)\quad \text{for any $x\in [-1,1]$}.
% \end{equation*}
% The central strategy involves a divide and conquer approach. While the original target function exhibits high-frequency oscillations throughout its domain, $f_1$ and $f_2$ oscillate only within a half-space, hence making them less oscillatory. Through a few layers of such (de)composition, it becomes smooth enough to be approximated by the last layer of shallow network. 

% \subsubsection{Rank 1}
% We start from a simple example by assuming the target function is $f(x)=\cos(2n\pi  x)$.

% $f(x)=\big(\cos(2n\pi x)-1\big)/2$ with $f(0)=0$.
% \begin{equation*}
% 	f_1=\begin{cases}
% 		\big(\cos(2n\pi x)-1\big)/2 & \tn{for}\  x\in [-1,0),\\
% 		x & \tn{for}\  x\in [0,1].
% 	\end{cases}
% \end{equation*}
% \begin{equation*}
% 	f_2=\begin{cases}
% 		x & \tn{for}\  x\in [-1,0),\\
% 		\big(\cos(2n\pi x)-1\big)/2  & \tn{for}\  x\in [0,1].
% 	\end{cases}
% \end{equation*}
% Then we have 
% \begin{equation*}
% 	f(x)=f_2\circ f_1(x)\quad \tn{for any $x\in [-1,1]$}
% \end{equation*}



% \subsubsection{Multi-component decomposition}
\subsubsection{Two-component decomposition}
\label{sec:1D:two:component:decom}

We demonstrate a simple ``divide and conquer'' strategy for an example function 
$f(x)=\cos(2n\pi x)$, a high frequency Fourier mode when $n$ is large. 
Define
\begin{equation*}
     f_2:(u,v)\in [-1,1]^2 \mapsto 
    \cos\big(n\pi  (u+1)\big)+ \cos\big(n\pi  (v-1)\big)\in \R.
\end{equation*}
and 
$\bmf_1=\begin{bsmallmatrix}
    f_{1,1}\\  f_{1,2}\end{bsmallmatrix}:[-1,1]\mapsto[-1,1]^2$, where the components $f_{1,1}$ and $f_{1,2}$ are given by
\begin{equation}\label{eq:component:f11}
	f_{1,1}(x)=\ReLU(2x)-1=\begin{cases}
		-1 & \tn{for}\  x\in [-1,0),\\
		2x-1 & \tn{for}\  x\in [0,1],
	\end{cases}
\end{equation}
and
\begin{equation}\label{eq:component:f12}
	f_{1,2}(x)=-\ReLU(-2x)+1=\begin{cases}
		2x+1 & \tn{for}\  x\in [-1,0),\\
		1  & \tn{for}\  x\in [0,1].
	\end{cases}
\end{equation}
% Define
% \begin{equation*}
%     \bmf_1:[-1,1]\mapsto[-1,1]^2, \qquad   \bmf_1=\begin{bmatrix}
%     f_{1,1}\\  f_{1,2}\end{bmatrix},
% \end{equation*}
% \begin{equation*}
% 	f_{1,1}(x)=\ReLU(2x)-1=\begin{cases}
% 		-1 & \tn{for}\  x\in [-1,0),\\
% 		2x-1 & \tn{for}\  x\in [0,1],
% 	\end{cases}
% \end{equation*}
% \begin{equation*}
% 	f_{1,2}(x)=-\ReLU(-2x)+1=\begin{cases}
% 		2x+1 & \tn{for}\  x\in [-1,0),\\
% 		1  & \tn{for}\  x\in [0,1],
% 	\end{cases}
% \end{equation*}
% and
% \begin{equation*}
%      f_2:(u,v)\in [-1,1]^2 \mapsto 
%     \cos\big(n\pi  (u+1)\big)+ \cos\big(n\pi  (v-1)\big)\in \R.
% \end{equation*}
Then for any $x\in [-1,1]$ we have 
\begin{equation*}
	\begin{split}
	    f(x)&=\cos\Big(n\pi  \cdot\ReLU(2x)\Big)+\cos\Big(-n\pi  \cdot\ReLU(-2x)\Big)
     \\& =\cos\Big(n\pi  \big(f_{1,1}(x)+1\big)\Big)+\cos\Big(n\pi  \big(f_{1,2}(x)-1\big)\Big)
 =
 f_2\circ \bmf_1(x).
	\end{split}
\end{equation*}
Through this decomposition and piecewise linear transformation, which can be approximated easily by a single layer of {\ReLU} network, one only needs to approximate a function that is smoother than the original $f$: $\bm{f}_1$ is simplified, while $f_2$ is reduced to half of the frequency of the original target function $f$. 
%In the rank 1 scenario, the strategy involves distributing the challenges through depth, utilizing composite structures. On the other hand, in the rank 2 scenario, challenges can be distributed across different ranks.
% In the case of rank 1, we can exclusively distribute the difficulty across the depth (complex structures). Conversely, in the case of rank 2, we can distribute the difficulty across the ranks.

We observe that this decomposition approach is universally applicable for any function $f:[-1,1]\mapsto\mathbb{R}$. Specifically, the decomposition is defined as
\begin{equation*}
    f_2:(u,v)\in[-1,1]^2\mapsto 
    f\big(\tfrac{u+1}{2}\big)+ f\big(\tfrac{v-1}{2}\big)-f(0)\in\R.
\end{equation*}
and 
$\bmf_1=\begin{bsmallmatrix}
    f_{1,1}\\  f_{1,2}\end{bsmallmatrix}:[-1,1]\mapsto[-1,1]^2$,
where $f_{1,1}$ and $f_{1,2}$ are given in \eqref{eq:component:f11} and \eqref{eq:component:f12}.
% \begin{equation*}
% 	f_{1,1}(x)=\ReLU(2x)-1=\begin{cases}
% 		-1 & \tn{for}\  x\in [-1,0),\\
% 		2x-1 & \tn{for}\  x\in [0,1],
% 	\end{cases}
% \end{equation*}
% % and
% \begin{equation*}
% 	f_{1,2}(x)=-\ReLU(-2x)+1=\begin{cases}
% 		2x+1 & \tn{for}\  x\in [-1,0),\\
% 		1  & \tn{for}\  x\in [0,1].
% 	\end{cases}
% \end{equation*}
% \begin{equation*}
%     \bmf_1:[-1,1]\mapsto[-1,1]^2, \qquad   \bmf_1=\begin{bmatrix}
%     f_{1,1}\\  f_{1,2}\end{bmatrix},
% \end{equation*}
% where the component functions $f_{1,1}$ and $f_{1,2}$ are defined by
% \begin{equation*}
% 	f_{1,1}(x)=\ReLU(2x)-1=\begin{cases}
% 		-1 & \tn{for}\  x\in [-1,0),\\
% 		2x-1 & \tn{for}\  x\in [0,1],
% 	\end{cases}
% \end{equation*}
% and
% \begin{equation*}
% 	f_{1,2}(x)=-\ReLU(-2x)+1=\begin{cases}
% 		2x+1 & \tn{for}\  x\in [-1,0),\\
% 		1  & \tn{for}\  x\in [0,1].
% 	\end{cases}
% \end{equation*}
% Moreover, 
% \begin{equation*}
%     f_2:(u,v)\in[-1,1]^2\mapsto 
%     f\big(\tfrac{u+1}{2}\big)+ f\big(\tfrac{v-1}{2}\big)-f(0)\in\R.
% \end{equation*}
% \begin{equation*}
% \bm{f}1:[-1,1]\to[-1,1]^2, \quad \bm{f}1=\begin{pmatrix}
% f{1,1}\ f{1,2}\end{pmatrix},
% \end{equation*}
% where
% \begin{equation*}
% f_2:[-1,1]^2\to\mathbb{R},\quad f_2:(u,v)\mapsto
% f\left(\frac{u+1}{2}\right) + f\left(\frac{v-1}{2}\right) - f(0),
% \end{equation*}
% and the component functions $f_{1,1}$ and $f_{1,2}$ are defined by
% \begin{equation*}
% f_{1,1}(x)=\text{ReLU}(2x)-1=\begin{cases}
% -1 & \text{for}\ x\in [-1,0),\
% 2x-1 & \text{for}\ x\in [0,1],
% \end{cases}
% \end{equation*}
% \begin{equation*}
% f_{1,2}(x)=-\text{ReLU}(-2x)+1=\begin{cases}
% 2x+1 & \text{for}\ x\in [-1,0),\
% 1 & \text{for}\ x\in [0,1].
% \end{cases}
% \end{equation*}
Hence, for any $x\in [-1,1]$, we achieve the following reconstruction of $f(x)$:
\begin{equation*}
	\begin{split}
	    f(x)&=f\Big(\tfrac{\ReLU(2x)}{2}\Big)+ f\Big(\tfrac{-\ReLU(-2x)}{2}\Big)-f(0)
     \\& =f\Big(\tfrac{f_{1,1}(x)+1}{2}\Big)+ f\Big(\tfrac{f_{1,2}(x)-1}{2}\Big)-f(0)
 =
 f_2\circ \bmf_1(x)
	\end{split}
\end{equation*}
demonstrating a structured decomposition that allows the function to be expressed through the composition of a smoother function with a piecewise (component-wise) transformation and rescaling.
% We remark that this decomposition is valid for all $f:[-1,1]\to\R$.
% \begin{equation*}
%     \bmf_1:[-1,1]\to[-1,1]^2, \qquad   \bmf_1=\begin{bmatrix}
%     f_{1,1}\\  f_{1,2}\end{bmatrix}
% \end{equation*}
% \begin{equation*}
%     f_2:[-1,1]^2\to\R,\quad  f_2:(u,v)\mapsto 
%     f\big(\tfrac{u+1}{2}\big)+ f\big(\tfrac{v-1}{2}\big)-f(0)
% \end{equation*}
% \begin{equation*}
% 	f_{1,1}(x)=\ReLU(2x)-1=\begin{cases}
% 		-1 & \tn{for}\  x\in [-1,0),\\
% 		2x-1 & \tn{for}\  x\in [0,1].
% 	\end{cases}
% \end{equation*}
% \begin{equation*}
% 	f_{1,2}(x)=-\ReLU(-2x)+1=\begin{cases}
% 		2x+1 & \tn{for}\  x\in [-1,0),\\
% 		1  & \tn{for}\  x\in [0,1].
% 	\end{cases}
% \end{equation*}
% Then for any $x\in [-1,1]$ we have 
% \begin{equation*}
% 	\begin{split}
% 	    f(x)&=f\Big(\tfrac{\ReLU(2x)}{2}\Big)+ f\Big(\tfrac{-\ReLU(-2x)}{2}\Big)-f(0)
%      \\& =f\Big(\tfrac{f_{1,1}(x)+1}{2}\Big)+ f\Big(\tfrac{f_{1,2}(x)-1}{2}\Big)-f(0)
%  =
%  f_2\circ \bmf_1(x)
% 	\end{split}
% \end{equation*}


\subsubsection{General multi-component decomposition}
\label{sec:1D:general:multi:component:decom}

Now we extend to a general multi-component adaptive decomposition, a ``divide and conquer" strategy, that can distribute the complexity of a target function evenly to multiple components. 


Given a sequence $x_0<x_1<\cdots<x_n$ where the target function is defined on the interval $[x_0,x_n]$, we will demonstrate how our new architecture allows us to partition the complexities of the function $f$ into smaller intervals $[x_{i-1},x_i]$. By rescaling each subinterval, one only needs to deal with a much smoother function in each interval. This approach enables us to effectively approximate the target function over the entire interval $[x_0,x_n]$. %This strategy exemplifies the \textbf{divide and conquer principle}.



Let $\calL_i:[a_i,b_i]\to [x_{i-1},x_i]$ be the linear map with 
\begin{equation}\label{eq:1D:Li}
    \calL_i(a_i)=x_{i-1} \quad \tn{and}\quad \calL_i(b_i)=x_i.
\end{equation}
Define
\begin{equation}\label{eq:1D:fi}
    f_i=f\circ \calL_i: [a_i,b_i]\to \R. 
\end{equation}




To decompose the target function into smoother pieces, we define a piecewise linear transformation 
$\psi_i$ using a linear combination of two \texttt{ReLU} functions (or a simple single layer network),
\begin{equation}\label{eq:1D:psii}
    \psi_i(x)=s_i\cdot \ReLU\left( {x-x_{i-1}} \right)
    -s_i\cdot \ReLU\left( {x-x_{i}} \right)+a_i.
\end{equation}
Here $s_i=\frac{b_i-a_i}{x_i-x_{i-1}}$ is the ``slope'' of $\calL_i^{-1}$, which is a local rescaling. For example, $f_i$ becomes a smoother function than $f$ after stretching $[x_{i-1},x_i]$ to a larger domain $[a_i,b_i]$. See an illustration of $\psi_i(x)$ in Figure~\ref{fig:psi:i}. 

\begin{figure}%[htbp]%[h]%[htbp!]
    \centering
    \includegraphics[width=0.615\textwidth]{figures/Others/psii.pdf}
    \caption{An illustration of $\psi_i(x)$.}
    \label{fig:psi:i}
\end{figure}

% \red{We can implement $\psi_i$ with two active \ReLU-activated neurons, i.e.,
% \begin{equation}\label{eq:1D:psii}
%     \psi_i(x)=s_i\cdot \ReLU\left( {x-x_{i-1}} \right)
%     -s_i\cdot \ReLU\left( {x-x_{i}} \right)+a_i,
% \end{equation}
% where $s_i=\frac{b_i-a_i}{x_i-x_{i-1}}$ is the ``slope'' of $\calL_i^{-1}$.} 
% It is easy to check
% \begin{equation*}%\label{eq:1D:psi}
%     \psi_i(x)=\begin{cases}
%         b_i &\tn{if  } x>x_{i},\\
%         \calL_i^{-1}(x) &\tn{if  }x\in [x_{i-1},x_i],\\
%         a_i &\tn{if  }x<x_{i-1}.
%     \end{cases}
% \end{equation*}

\begin{figure}%[htbp]%[h]%[htbp!]
    \centering	
    \begin{subfigure}[c]{0.998\textwidth}
    \centering            \includegraphics[width=0.827998055\textwidth]{figures/Others/fDecomposition.pdf}
    \subcaption{Decomposition of target function $f=c+ \sum_{i=1}^n f_i\circ \psi_i$:
    oscillatory $f$ to smooth $f_i$'s.}
    %		\label{subfig:f3}
    \end{subfigure}\\
    \vspace*{12pt}
        \begin{subfigure}[c]{0.998\textwidth}
    \centering            \includegraphics[width=0.827998055\textwidth]{figures/Others/hDecomposition.pdf}
    \subcaption{Neural network architecture of $h=c+ \sum_{i=1}^n h_i\circ \psi_i$ by using $h_i\approx f_i$.}
    %		\label{subfig:f3}
    \end{subfigure}
    \caption{
Visual representations of the decompositions of $f$ and $h$ are provided with $c=\sum_{i=0}^{n-1}f(x_i)$ being a constant and $s_i$ being the slope. Here, the function $f$ is dissected into several simpler functions, labeled as $f_i$. Each $f_i$ represents a simplified and more manageable segment of $f$, allowing for the straightforward application of sub-network $h_i$ to closely approximate $f_i$, even with the use of shallow networks.
    }
    	\label{fig:1D:decomposition:approx}
    % 		\vskip 0.02in
\end{figure}

\begin{theorem}
\label{thm:decomposition:1D}
Given \(x_0<x_1<\cdots<x_n\),
    suppose $\calL_{i}$ and $\psi_i$ are given in Equations~\eqref{eq:1D:Li} and \eqref{eq:1D:psii}, respectively. Then the target function $f:[x_0,x_n]\to\R$ has the following (smoother) decomposition ($f_i$) with a piecewise linear transformation ($\psi_i$),
    
\begin{equation*}
    \begin{split}
        f(x)=\sum_{i=1}^n f_i\circ \psi_i(x)
    -\underbrace{ \sum_{i=1}^{n-1}f(x_i)}_{\tn{constant}}\quad \tn{for any $x\in [x_{0},x_n]$,}
    \end{split}
\end{equation*}
where $f_{i}$ is given in Equation~\eqref{eq:1D:fi}.
\end{theorem}

%This theorem implies that an oscillatory
%function can be decomposed into linear combinations and compositions of non-oscillatory functions.

% \red{Add a figure to illustrate the decomposition later.}



% \red{Consider an example of a target function $f$ that utilizes adaptive decomposition. For instance, we select the points $x_0, x_1, x_2, x_3, x_4$ as $-1, -0.7, -0.4, 0, 1$, respectively.}

% Next, let verify that for any $x\in [x_0,x_n]=\cup_{i=1}^{n} [x_{i-1},x_i]$, 
% it holds that
% \begin{equation*}
%     f(x)=\sum_{i=1}^n f_i\circ \psi_i(x)-\underbrace{\sum_{i=1}^{n-1}f(x_i)}_{\tn{constant}}.
% \end{equation*}
\begin{proof}%[Proof of Theorem~\ref{thm:decomposition:1D}]
By definition of $\psi_i$ in Equation~\eqref{eq:1D:psii}, 
it is easy to check
\begin{equation*}%\label{eq:1D:psi}
    \psi_i(x)=\begin{cases}
        b_i &\tn{if  } x>x_{i},\\
        \calL_i^{-1}(x) &\tn{if  }x\in [x_{i-1},x_i],\\
        a_i &\tn{if  }x<x_{i-1},
    \end{cases}
    \quad \Longrightarrow \quad 
        \psi_i(x)=\begin{cases}
        b_i &\tn{if  }i\le j-1,\\
        \calL_j^{-1}(x) &\tn{if  }i= j ,\\
        a_i &\tn{if  }i\ge j+1,
    \end{cases}
\end{equation*}
for a fixed $j\in \{1,2,\cdots,n\}$ and any $x\in [x_{j-1}, x_j]$.
% Fixing $j$, for any $x\in [x_{j-1}, x_j]$, we have
% \begin{equation*}
%     \psi_i(x)=\begin{cases}
%         b_i &\tn{if  }i\le j-1,\\
%         \calL_j^{-1}(x) &\tn{if  }i= j ,\\
%         a_i &\tn{if  }i\ge j+1,
%     \end{cases}
% \end{equation*}
% Then $\psi_i(x)=x_i$ if $i\le j-1$
% and $\psi_i(x)=x_{i-1}$ if $i\ge j+1$.
% Then we have
% from which we deduce
% \begin{equation*}
%     \begin{split}
%         \sum_{i=1}^n f_i\circ \psi_i(x)=\sum_{i=1}^n f\circ \calL_i\circ \psi_i(x)
%     &=\sum_{i=1}^{j-1} f\circ \calL_i\circ \psi_i(x)
%     +f\circ \calL_j\circ \psi_j(x)
%     +\sum_{i=j+1}^n f\circ \calL_i\circ \psi_i(x)
%     \\& = \sum_{i=1}^{j-1} f\circ \calL_i(b_i)
%     +f\circ \calL_j\circ\calL_j^{-1}(x) 
%     +\sum_{i=j+1}^n f\circ \calL_i(a_i)
%     \\& = \sum_{i=1}^{j-1} f (x_i)
%     +f(x)
%     +\sum_{i=j+1}^n f (x_{i-1})
%     =f(x)+\underbrace{\sum_{i=1}^{n-1}f(x_i)}_{\tn{constant}}.
%     \end{split}
% \end{equation*}
It follows that
\begin{equation*}
    \begin{split}
        % & \phantom{=\;\;}
        \sum_{i=1}^n f_i\circ \psi_i(x)
        &=\sum_{i=1}^n f\circ \calL_i\circ \psi_i(x)
    \\&=\sum_{i=1}^{j-1} f\circ \calL_i\circ \psi_i(x)
    +f\circ \calL_j\circ \psi_j(x)
    +\sum_{i=j+1}^n f\circ \calL_i\circ \psi_i(x)
    \\& = \sum_{i=1}^{j-1} f\circ \calL_i(b_i)
    +f\circ \calL_j\circ\calL_j^{-1}(x) 
    +\sum_{i=j+1}^n f\circ \calL_i(a_i)
    \\& = \sum_{i=1}^{j-1} f (x_i)
    +f(x)
    +\sum_{i=j+1}^n f (x_{i-1})
    =f(x)+\underbrace{\sum_{i=1}^{n-1}f(x_i)}_{\tn{constant}}.
    \end{split}
\end{equation*}
It follows that
\begin{equation*}
    \begin{split}
        f(x)=\sum_{i=1}^n f_i\circ \psi_i(x)
    -\underbrace{\sum_{i=1}^{n-1}f(x_i)}_{\tn{constant}}\quad \tn{for any $x\in [x_{j-1},x_j]$.}
    \end{split}
\end{equation*}
Since $j$ is arbitrary, the above equation holds for all
$x=\cup_{j=1}^n[x_{j-1}, x_j]= [x_0,x_n]$. 
\end{proof}

For each smoother $f_i$, one can use a shallow network component $\phi_i$, a linear combination of random basis functions, to approximate $f_i$ well on $[a_i,b_i]$. Then
\begin{equation*}
    f(x)=\sum_{i=1}^n f_i\circ \psi_i(x)-\underbrace{\sum_{i=1}^{n-1}f(x_i)}_{\tn{constant}}
    \approx \sum_{i=1}^n \phi_i\circ \psi_i(x)-\underbrace{\sum_{i=1}^{n-1}f(x_i)}_{\tn{constant}}
   \eqqcolon h(x),
\end{equation*}
$h(x)$ is a one-hidden-layer neural network approximation of the target function $f(x)$ that can approximate a complex function better than a single layer.
See Figure~\ref{fig:1D:decomposition:approx} for an illustration. In practice, one can choose repeated decomposition using a multi-component and multi-layer network structure which is the motivation for MMNN. It is well-known that neural networks can approximate smooth functions well. For
localized rapid change/oscillation, our construction shows that a small network in terms of the width as well as the number of components and layers can achieve adaptive decomposition and deal with it rather easily. Hence, MMNN is effective in approximating a function with localized fine features. This is an important advantage in dealing with low-dimensional structures embedded in high dimensions. The most difficult situation is approximating global highly oscillatory functions, especially with diverse frequency modes, for which wider networks with more components and layers are needed to deal with both the complexity and curse of dimensions.


\subsubsection{Examples}
\label{sec:1D:decom:examples}


% To provide an intuitive demonstration of the decomposition described earlier, we present several examples. Initially, we consider the Runge function defined by $f(x) = \frac{1}{25x^2+1}$, and modify it to $f(x) = \frac{1}{1000x^2+1}$. For simplicity, we select the points $x_0, x_1, x_2, x_3, x_4$ as $-1, -0.2, 0, 0.2, 1$, and set $a_i=1$ and $b_i=1$ for all i. Although in practice the choices of $x_i$, $a_i$, and $b_i$ can be adaptive, here we use a straightforward selection. Consequently, $f(x)$ can be decomposed into a series of smooth components, expressed as:

Here we use two examples to demonstrate the complexity decomposition strategy presented in the previous section. We start with the Runge function $f(x) = \frac{1}{25x^2+1}$ and modify it to $f(x) = \frac{1}{1000x^2+1}$, which has a localized rapid change near $0$. As an example, we use four components $n=4$, choose points $x_0, x_1, x_2, x_3, x_4$ at $-1, -0.2, 0, 0.2, 1$, and let $a_i = -1$ and $b_i = 1$ for all 
 $i$. In practice, each component is approximated by a single-layer network - a linear combination of basis functions, and trained by an optimization method, e.g., Adam. Our examples here are just a proof of concept for the decomposition of a target function into smoother components using MMNN structure in the form
\[
f(x) = \sum_{i=1}^4 f_i\circ \psi_i(x) - \underbrace{\sum_{i=1}^{3} f(x_i)}_{\text{constant}},
\]
where $f_i$ and $\psi_i$ (piecewise tranformation/rescaling) are defined as in~\eqref{eq:1D:fi} and \eqref{eq:1D:psii}, respectively. These components are illustrated in Figure~\ref{fig:Runge:F:decom}.
% \[
% f_i = f \circ \mathcal{L}_i: [a_i, b_i] \to \mathbb{R}, \qquad
% \psi_i(x) = \begin{cases} 
% b_i & \text{if } x > x_i, \\
% \mathcal{L}_i^{-1}(x) & \text{if } x \in [x_{i-1}, x_i], \\
% a_i & \text{if } x < x_{i-1}.
% \end{cases}
% \]
Each component is relatively smooth, making it easier for approximation and learning through shallow networks. This approach essentially utilizes a divide-and-conquer principle.

% \clearpage
% \begin{figure}%[htbp]%[h]%[htbp!]
%     \centering	
%     \begin{subfigure}[c]{0.48\textwidth}
%     \centering            \includegraphics[width=0.998055\textwidth]{figures/DecomEg1D/1DdecomRungeF1.pdf}
%     % \subcaption{Decompostion}
%     %		\label{subfig:f3}
%     \end{subfigure}\hfill
%         \begin{subfigure}[c]{0.48\textwidth}
%     \centering            \includegraphics[width=0.998055\textwidth]{figures/DecomEg1D/1DdecomRungeF2.pdf}
%     % \subcaption{Neural .}
%     %		\label{subfig:f3}
%     \end{subfigure}\\
%         \begin{subfigure}[c]{0.248\textwidth}
%     \centering            \includegraphics[width=0.998055\textwidth]{figures/DecomEg1D/1DdecomRungeFi1.pdf}
%     % \subcaption{Decompostion}
%     %		\label{subfig:f3}
%     \end{subfigure}\hfill
%         \begin{subfigure}[c]{0.248\textwidth}
%     \centering            \includegraphics[width=0.998055\textwidth]{figures/DecomEg1D/1DdecomRungeFi2.pdf}
%     % \subcaption{Neural .}
%     %		\label{subfig:f3}
%     \end{subfigure}\hfill
%     \begin{subfigure}[c]{0.248\textwidth}
%     \centering            \includegraphics[width=0.998055\textwidth]{figures/DecomEg1D/1DdecomRungeFi3.pdf}
%     % \subcaption{Decompostion}
%     %		\label{subfig:f3}
%     \end{subfigure}\hfill
%         \begin{subfigure}[c]{0.248\textwidth}
%     \centering            \includegraphics[width=0.998055\textwidth]{figures/DecomEg1D/1DdecomRungeFi4.pdf}
%     % \subcaption{Neural .}
%     %		\label{subfig:f3}
%     \end{subfigure}\\
%         \begin{subfigure}[c]{0.248\textwidth}
%     \centering            \includegraphics[width=0.998055\textwidth]{figures/DecomEg1D/1DdecomRungePsii1.pdf}
%     % \subcaption{Decompostion}
%     %		\label{subfig:f3}
%     \end{subfigure}\hfill
%         \begin{subfigure}[c]{0.248\textwidth}
%     \centering            \includegraphics[width=0.998055\textwidth]{figures/DecomEg1D/1DdecomRungePsii2.pdf}
%     % \subcaption{Neural .}
%     %		\label{subfig:f3}
%     \end{subfigure}\hfill
%     \begin{subfigure}[c]{0.248\textwidth}
%     \centering            \includegraphics[width=0.998055\textwidth]{figures/DecomEg1D/1DdecomRungePsii3.pdf}
%     % \subcaption{Decompostion}
%     %		\label{subfig:f3}
%     \end{subfigure}\hfill
%         \begin{subfigure}[c]{0.248\textwidth}
%     \centering            \includegraphics[width=0.998055\textwidth]{figures/DecomEg1D/1DdecomRungePsii4.pdf}
%     % \subcaption{Neural .}
%     %		\label{subfig:f3}
%     \end{subfigure}
%     \caption{
% Illustration of $f(x)=\frac{1}{1000x^2+1}$ and its multi-component decomposition through $f_i$ and $\psi_i$ $i=1,2,3,4$, where
% $f(x) = \sum_{i=1}^4 f_i\circ \psi_i(x) -  {\sum_{i=1}^{3} f(x_i)}$.
% % \\
% % \red{Choose points $x_0, x_1, x_2, x_3, x_4$ at $-1, -0.2, 0, 0.2, 1$}
% % \\ 
% % \red{how to choose $x_i$: length of curve;  $\int |f'(t)|ds$,  $ds= \sqrt{1+|f'(t)|^2}dt$}
%     }
%     	\label{fig:Runge:F:decom}
%     % 		\vskip 0.02in
% \end{figure}

\begin{figure}%[htbp]%[h]%[htbp!]
    \centering	
    \begin{subfigure}[c]{0.1948\textwidth}
    \centering            \includegraphics[width=0.998055\textwidth]{figures/DecomEg1D/1DdecomRungeF1.pdf}
    % \subcaption{Decompostion}
    %		\label{subfig:f3}
    \end{subfigure}\hfill
    \begin{subfigure}[c]{0.1948\textwidth}
    \centering            \includegraphics[width=0.998055\textwidth]{figures/DecomEg1D/1DdecomRungeFi1.pdf}
    % \subcaption{Decompostion}
    %		\label{subfig:f3}
    \end{subfigure}\hfill
        \begin{subfigure}[c]{0.1948\textwidth}
    \centering            \includegraphics[width=0.998055\textwidth]{figures/DecomEg1D/1DdecomRungeFi2.pdf}
    % \subcaption{Neural .}
    %		\label{subfig:f3}
    \end{subfigure}\hfill
    \begin{subfigure}[c]{0.1948\textwidth}
    \centering            \includegraphics[width=0.998055\textwidth]{figures/DecomEg1D/1DdecomRungeFi3.pdf}
    % \subcaption{Decompostion}
    %		\label{subfig:f3}
    \end{subfigure}\hfill
        \begin{subfigure}[c]{0.1948\textwidth}
    \centering            \includegraphics[width=0.998055\textwidth]{figures/DecomEg1D/1DdecomRungeFi4.pdf}
    % \subcaption{Neural .}
    %		\label{subfig:f3}
    \end{subfigure}\\
    \begin{subfigure}[c]{0.1948\textwidth}
    \centering            \includegraphics[width=0.998055\textwidth]{figures/DecomEg1D/1DdecomRungeF2.pdf}
    % \subcaption{Neural .}
    %		\label{subfig:f3}
    \end{subfigure}
        \begin{subfigure}[c]{0.1948\textwidth}
    \centering            \includegraphics[width=0.998055\textwidth]{figures/DecomEg1D/1DdecomRungePsii1.pdf}
    % \subcaption{Decompostion}
    %		\label{subfig:f3}
    \end{subfigure}\hfill
        \begin{subfigure}[c]{0.1948\textwidth}
    \centering            \includegraphics[width=0.998055\textwidth]{figures/DecomEg1D/1DdecomRungePsii2.pdf}
    % \subcaption{Neural .}
    %		\label{subfig:f3}
    \end{subfigure}\hfill
    \begin{subfigure}[c]{0.1948\textwidth}
    \centering            \includegraphics[width=0.998055\textwidth]{figures/DecomEg1D/1DdecomRungePsii3.pdf}
    % \subcaption{Decompostion}
    %		\label{subfig:f3}
    \end{subfigure}\hfill
        \begin{subfigure}[c]{0.1948\textwidth}
    \centering            \includegraphics[width=0.998055\textwidth]{figures/DecomEg1D/1DdecomRungePsii4.pdf}
    % \subcaption{Neural .}
    %		\label{subfig:f3}
    \end{subfigure}
    \caption{
Illustrations of $f(x)=\frac{1}{1000x^2+1}$ and its multi-component decomposition through $f_i$ and $\psi_i$ $i=1,2,3,4$, where
$f(x) = \sum_{i=1}^4 f_i\circ \psi_i(x) -  {\sum_{i=1}^{3} f(x_i)}$.
% \\
% \red{Choose points $x_0, x_1, x_2, x_3, x_4$ at $-1, -0.2, 0, 0.2, 1$}
% \\ 
% \red{how to choose $x_i$: length of curve;  $\int |f'(t)|ds$,  $ds= \sqrt{1+|f'(t)|^2}dt$}
    }
    	\label{fig:Runge:F:decom}
    % 		\vskip 0.02in
\end{figure}

% \begin{figure}%[htbp]%[h]%[htbp!]
%     \centering	
%     \begin{subfigure}[c]{0.48\textwidth}
%     \centering            \includegraphics[width=0.998055\textwidth]{figures/DecomEg1D/1DdecomRunge2F1.pdf}
%     % \subcaption{Decompostion}
%     %		\label{subfig:f3}
%     \end{subfigure}\hfill
%         \begin{subfigure}[c]{0.48\textwidth}
%     \centering            \includegraphics[width=0.998055\textwidth]{figures/DecomEg1D/1DdecomRunge2F2.pdf}
%     % \subcaption{Neural .}
%     %		\label{subfig:f3}
%     \end{subfigure}\\
%         \begin{subfigure}[c]{0.248\textwidth}
%     \centering            \includegraphics[width=0.998055\textwidth]{figures/DecomEg1D/1DdecomRunge2Fi1.pdf}
%     % \subcaption{Decompostion}
%     %		\label{subfig:f3}
%     \end{subfigure}\hfill
%         \begin{subfigure}[c]{0.248\textwidth}
%     \centering            \includegraphics[width=0.998055\textwidth]{figures/DecomEg1D/1DdecomRunge2Fi2.pdf}
%     % \subcaption{Neural .}
%     %		\label{subfig:f3}
%     \end{subfigure}\hfill
%     \begin{subfigure}[c]{0.248\textwidth}
%     \centering            \includegraphics[width=0.998055\textwidth]{figures/DecomEg1D/1DdecomRunge2Fi3.pdf}
%     % \subcaption{Decompostion}
%     %		\label{subfig:f3}
%     \end{subfigure}\hfill
%         \begin{subfigure}[c]{0.248\textwidth}
%     \centering            \includegraphics[width=0.998055\textwidth]{figures/DecomEg1D/1DdecomRunge2Fi4.pdf}
%     % \subcaption{Neural .}
%     %		\label{subfig:f3}
%     \end{subfigure}\\
%         \begin{subfigure}[c]{0.248\textwidth}
%     \centering            \includegraphics[width=0.998055\textwidth]{figures/DecomEg1D/1DdecomRunge2Psii1.pdf}
%     % \subcaption{Decompostion}
%     %		\label{subfig:f3}
%     \end{subfigure}\hfill
%         \begin{subfigure}[c]{0.248\textwidth}
%     \centering            \includegraphics[width=0.998055\textwidth]{figures/DecomEg1D/1DdecomRunge2Psii2.pdf}
%     % \subcaption{Neural .}
%     %		\label{subfig:f3}
%     \end{subfigure}\hfill
%     \begin{subfigure}[c]{0.248\textwidth}
%     \centering            \includegraphics[width=0.998055\textwidth]{figures/DecomEg1D/1DdecomRunge2Psii3.pdf}
%     % \subcaption{Decompostion}
%     %		\label{subfig:f3}
%     \end{subfigure}\hfill
%         \begin{subfigure}[c]{0.248\textwidth}
%     \centering            \includegraphics[width=0.998055\textwidth]{figures/DecomEg1D/1DdecomRunge2Psii4.pdf}
%     % \subcaption{Neural .}
%     %		\label{subfig:f3}
%     \end{subfigure}
%     \caption{
% Illustrations of $f(x)=\frac{1}{1000x^2+1}$ and its decomposition components $f_i$ and $\psi_i$. Note that $f$ can decomposed as
% $f(x) = \sum_{i=1}^4 f_i\circ \psi_i(x) -  {\sum_{i=1}^{3} f(x_i)}$. \\ \red{Equally partition the total variation;\quad\\ choose points $x_0, x_1, x_2, x_3, x_4$ at $-1, -1/\sqrt{1000}, 0, 1/\sqrt{1000}, 1$}
%     }
%     	\label{fig:Runge:F:decom}
%     % 		\vskip 0.02in
% \end{figure}



% Note that the example discussed previously exhibits only local high frequency. Below, we provide an example with global high frequency:
% \[
% f(x) = \cos^2(6\pi x) + \sin(10\pi x^2).
% \]
% For simplicity, we choose points $x_0, x_1, x_2, x_3, x_4$ at $-1, -0.2, 0, 0.2, 1$, and assign $a_i = 1$ and $b_i = -1$ for all 
%  $i$.
% As illustrated in Figure~\ref{fig:CosSin:F:decom}, even though $f(x)$ features global high frequency, it can still be decomposed into a series of components that are relatively smooth. This smoothness makes them amenable to approximation and learning via shallow networks, effectively applying the divide and conquer principle.

The second example is a globally oscillatory function of the form
\[
f(x) = \cos^2(6\pi x) + \sin(10\pi x^2).
\]
Again we illustrate using four components $n=4$, selecting points $x_0, x_1, x_2, x_3, x_4$ at $-1, -0.7, 0, 0.7, 1$, and setting $a_i = -1$ and $b_i = 1$ for all $i$. As shown in Figure~\ref{fig:CosSin:F:decom}, the target function $f(x)$ is decomposed into components that are less oscillatory again facilitating their approximation and learning through shallow networks.


% \begin{figure}%[htbp]%[h]%[htbp!]
%     \centering	
%     \begin{subfigure}[c]{0.48\textwidth}
%     \centering            \includegraphics[width=0.998055\textwidth]{figures/DecomEg1D/1DdecomCosSinF1.pdf}
%     % \subcaption{Decompostion}
%     %		\label{subfig:f3}
%     \end{subfigure}\hfill
%         \begin{subfigure}[c]{0.48\textwidth}
%     \centering            \includegraphics[width=0.998055\textwidth]{figures/DecomEg1D/1DdecomCosSinF2.pdf}
%     % \subcaption{Neural .}
%     %		\label{subfig:f3}
%     \end{subfigure}\\
%         \begin{subfigure}[c]{0.248\textwidth}
%     \centering            \includegraphics[width=0.998055\textwidth]{figures/DecomEg1D/1DdecomCosSinFi1.pdf}
%     % \subcaption{Decompostion}
%     %		\label{subfig:f3}
%     \end{subfigure}\hfill
%         \begin{subfigure}[c]{0.248\textwidth}
%     \centering            \includegraphics[width=0.998055\textwidth]{figures/DecomEg1D/1DdecomCosSinFi2.pdf}
%     % \subcaption{Neural .}
%     %		\label{subfig:f3}
%     \end{subfigure}\hfill
%     \begin{subfigure}[c]{0.248\textwidth}
%     \centering            \includegraphics[width=0.998055\textwidth]{figures/DecomEg1D/1DdecomCosSinFi3.pdf}
%     % \subcaption{Decompostion}
%     %		\label{subfig:f3}
%     \end{subfigure}\hfill
%         \begin{subfigure}[c]{0.248\textwidth}
%     \centering            \includegraphics[width=0.998055\textwidth]{figures/DecomEg1D/1DdecomCosSinFi4.pdf}
%     % \subcaption{Neural .}
%     %		\label{subfig:f3}
%     \end{subfigure}\\
%         \begin{subfigure}[c]{0.248\textwidth}
%     \centering            \includegraphics[width=0.998055\textwidth]{figures/DecomEg1D/1DdecomCosSinPsii1.pdf}
%     % \subcaption{Decompostion}
%     %		\label{subfig:f3}
%     \end{subfigure}\hfill
%         \begin{subfigure}[c]{0.248\textwidth}
%     \centering            \includegraphics[width=0.998055\textwidth]{figures/DecomEg1D/1DdecomCosSinPsii2.pdf}
%     % \subcaption{Neural .}
%     %		\label{subfig:f3}
%     \end{subfigure}\hfill
%     \begin{subfigure}[c]{0.248\textwidth}
%     \centering            \includegraphics[width=0.998055\textwidth]{figures/DecomEg1D/1DdecomCosSinPsii3.pdf}
%     % \subcaption{Decompostion}
%     %		\label{subfig:f3}
%     \end{subfigure}\hfill
%         \begin{subfigure}[c]{0.248\textwidth}
%     \centering            \includegraphics[width=0.998055\textwidth]{figures/DecomEg1D/1DdecomCosSinPsii4.pdf}
%     % \subcaption{Neural .}
%     %		\label{subfig:f3}
%     \end{subfigure}
%     \caption{
% Illustration of $f(x)=\cos^2(6\pi x)+\sin(10\pi x^2)$ and its decomposition components $f_i$ and $\psi_i$ such that
% $f(x) = \sum_{i=1}^4 f_i\circ \psi_i(x) -  {\sum_{i=1}^{3} f(x_i)}$.
%     }
%     	\label{fig:CosSin:F:decom}
%     % 		\vskip 0.02in
% \end{figure}


\begin{figure}%[htbp]%[h]%[htbp!]
    \centering	
    \begin{subfigure}[c]{0.1948\textwidth}
    \centering            \includegraphics[width=0.998055\textwidth]{figures/DecomEg1D/1DdecomCosSinF1.pdf}
    % \subcaption{Decompostion}
    %		\label{subfig:f3}
    \end{subfigure}\hfill
        \begin{subfigure}[c]{0.1948\textwidth}
    \centering            \includegraphics[width=0.998055\textwidth]{figures/DecomEg1D/1DdecomCosSinFi1.pdf}
    % \subcaption{Decompostion}
    %		\label{subfig:f3}
    \end{subfigure}\hfill
        \begin{subfigure}[c]{0.1948\textwidth}
    \centering            \includegraphics[width=0.998055\textwidth]{figures/DecomEg1D/1DdecomCosSinFi2.pdf}
    % \subcaption{Neural .}
    %		\label{subfig:f3}
    \end{subfigure}\hfill
    \begin{subfigure}[c]{0.1948\textwidth}
    \centering            \includegraphics[width=0.998055\textwidth]{figures/DecomEg1D/1DdecomCosSinFi3.pdf}
    % \subcaption{Decompostion}
    %		\label{subfig:f3}
    \end{subfigure}\hfill
        \begin{subfigure}[c]{0.1948\textwidth}
    \centering            \includegraphics[width=0.998055\textwidth]{figures/DecomEg1D/1DdecomCosSinFi4.pdf}
    % \subcaption{Neural .}
    %		\label{subfig:f3}
    \end{subfigure}\\
            \begin{subfigure}[c]{0.1948\textwidth}
    \centering            \includegraphics[width=0.998055\textwidth]{figures/DecomEg1D/1DdecomCosSinF2.pdf}
    % \subcaption{Neural .}
    %		\label{subfig:f3}
    \end{subfigure}\hfill
        \begin{subfigure}[c]{0.1948\textwidth}
    \centering            \includegraphics[width=0.998055\textwidth]{figures/DecomEg1D/1DdecomCosSinPsii1.pdf}
    % \subcaption{Decompostion}
    %		\label{subfig:f3}
    \end{subfigure}\hfill
        \begin{subfigure}[c]{0.1948\textwidth}
    \centering            \includegraphics[width=0.998055\textwidth]{figures/DecomEg1D/1DdecomCosSinPsii2.pdf}
    % \subcaption{Neural .}
    %		\label{subfig:f3}
    \end{subfigure}\hfill
    \begin{subfigure}[c]{0.1948\textwidth}
    \centering            \includegraphics[width=0.998055\textwidth]{figures/DecomEg1D/1DdecomCosSinPsii3.pdf}
    % \subcaption{Decompostion}
    %		\label{subfig:f3}
    \end{subfigure}\hfill
        \begin{subfigure}[c]{0.1948\textwidth}
    \centering            \includegraphics[width=0.998055\textwidth]{figures/DecomEg1D/1DdecomCosSinPsii4.pdf}
    % \subcaption{Neural .}
    %		\label{subfig:f3}
    \end{subfigure}
    \caption{
Illustrations of $f(x)=\cos^2(6\pi x)+\sin(10\pi x^2)$ and its decomposition components $f_i$ and $\psi_i$ such that
$f(x) = \sum_{i=1}^4 f_i\circ \psi_i(x) -  {\sum_{i=1}^{3} f(x_i)}$.
    }
    	\label{fig:CosSin:F:decom}
    		% \vskip 0.2in
\end{figure}



\subsection{High dimensional cases}
\label{sec:decomposition:highD}

% Let us now consider the extension to multi-dimension using two dimensions as an example since the simple dimension-by-dimension strategy applies to any dimension. 
Let us now consider the extension to multiple dimensions, using the case of two dimensions as an example, since the straightforward dimension-by-dimension strategy can be applied to any number of dimensions.

Given \(x_0<x_1<\cdots<x_n\) and \(y_0<y_1<\cdots<y_m\), dividing the domain of the function \(f(x,y)\) into small Cartesian rectangles \([x_{i-1},x_i]\times [y_{j-1}, y_j]\).
%simplifies the problem by focusing on smaller segments. This approach, which applies the \textbf{divide and conquer principle}, effectively manages the function across the full domain \([x_0,x_n]\times [y_0,y_m]\).
Let $\calL_{1,i}:[a_i,b_i]\to [x_{i-1},x_i]$ and $\calL_{2,j}:[c_i,d_i]\to [y_{j-1},y_j]$ be the linear maps with 
\begin{equation}
\label{eq:L:1i:2j}
\begin{cases}
        \calL_{1,i}(a_i)=x_{i-1},\\
    \calL_{1,i}(b_i)=x_i
\end{cases}
     \quad \tn{and}\quad\,\   
     \begin{cases}
        \calL_{2,j}(c_i)=y_{j-1},\\
    \calL_{2,j}(d_i)=y_j.
\end{cases}
\end{equation}
For $i=1,2,\cdots,n$ and $j=1,2,\cdots,m$, we define
\begin{equation}
\label{eq:fij}
\left\{\begin{array}{l}
    f_{i,0}(x,y)\coloneqq  f\Big(\calL_{1,i}(x),\, y\Big),
\\
    f_{0,j}(x,y)\coloneqq  f\Big(x,\, \calL_{2,j}(y)\Big),
    \\
% \end{equation}
% \begin{equation}
% \label{eq:fij}
    f_{i,j}(x,y)\coloneqq  f\Big(\calL_{1,i}(x),\, \calL_{2,j}(y)\Big)=f_{0,j}\Big(\calL_{1,i}(x),\, y\Big)=f_{i,0}\Big(x,\, \calL_{2,j}(y)\Big).
\end{array}\right.
\end{equation}
It is evident that with appropriate transformation and rescaling, \(f_{i,0}(x,y)\) is smooth in \(x\) when \(y\) is held constant, \(f_{0,j}(x,y)\) is smooth in \(y\) when \(x\) is fixed, and \(f_{i,j}(x,y)\) is smooth in both \(x\) and \(y\). 
Define
\begin{equation}\label{eq:psii:phij}
    \psi_i(x)=\begin{cases}
        b_i &\tn{if  } x>x_{i},\\
        \calL_{1,i}^{-1}(x) &\tn{if  }x\in [x_{i-1},x_i],\\
        a_i &\tn{if  }x<x_{i-1}
    \end{cases}
    \quad \tn{and}\quad 
        \phi_j(y)=\begin{cases}
        d_j &\tn{if  } y>y_{j},\\
        \calL_{2,j}^{-1}(y) &\tn{if  }y\in [y_{j-1},y_j],\\
        c_j &\tn{if  }y<y_{j-1}.
    \end{cases}
\end{equation}
The theorem below provides a decomposition of $f$ that fits into the MMNN structure.
% \red{
% The proof of this theorem can be found in the appendix.
% }
% We claim 
% \begin{equation}
% \label{eq:highD:decomposition}
%     \begin{split}          
%         f( x, \, y)
%         & =\sum_{i=1}^n \sum_{j=1}^m f_{i,j}\Big( \psi_i(x), \, \phi_j(y)\Big)-
%         \sum_{i=1}^n\sum_{j=1}^{m-1} f_{i,0}\Big(\psi_i(x),\, y_{j}\Big)
%         \\ & \quad\   -  \sum_{i=1}^{n-1} \sum_{j=1}^m f_{0,j}\Big( x_i, \, \phi_j(y)\Big)+
%         \sum_{i=1}^{n-1}\sum_{j=1}^{m-1} f (x_i,\, y_{j})
%     \end{split}
% \end{equation}
% for all $(x,y)\in [x_0,x_n]\times [y_0,y_m]$. This implies that 
% \(f(x, y)\) can be expressed as a series of smooth function components.

\begin{theorem}
\label{thm:decomposition:2D}
Given \(x_0<x_1<\cdots<x_n\) and \(y_0<y_1<\cdots<y_m\),
    suppose $\calL_{1,i}, \calL_{2,j}$ and $\psi_i,\phi_j$ are given in Equations~\eqref{eq:L:1i:2j} and \eqref{eq:psii:phij}, respectively. Then the function $f:[x_0,x_n]\times [y_0, y_m]\to\R$ can be expressed as 
\begin{equation}
\label{eq:highD:decomposition}
    \begin{split}          
        f( x, \, y)
        & =\sum_{i=1}^n \sum_{j=1}^m f_{i,j}\Big( \psi_i(x), \, \phi_j(y)\Big)-
        \sum_{i=1}^n\sum_{j=1}^{m-1} f_{i,0}\Big(\psi_i(x),\, y_{j}\Big)
        \\ & \quad\   -  \sum_{i=1}^{n-1} \sum_{j=1}^m f_{0,j}\Big( x_i, \, \phi_j(y)\Big)+
        \sum_{i=1}^{n-1}\sum_{j=1}^{m-1} f (x_i,\, y_{j})
    \end{split}
\end{equation}
for all $(x,y)\in [x_0,x_n]\times [y_0,y_m]$, where $f_{i,j}$ are given in 
Equation~\eqref{eq:fij}.
% Equations~\eqref{eq:fi0:f0j} and \eqref{eq:fij}.
\end{theorem}



%This theorem implies that \(f(x, y)\) can be expressed as a series of smooth function components. Below let us give the detailed proof of Theorem~\ref{thm:decomposition:2D}.

\begin{proof}%[Proof of Theorem~\ref{thm:decomposition:2D}]
% of Equation~\eqref{eq:highD:decomposition}.
Fixing $(k,j)$, for any $(x,y)\in [x_{k-1}, x_k]\times [y_{\ell-1}, y_\ell]$, we have
\begin{equation*}
    \psi_i(x)=\begin{cases}
        b_i &\tn{if  }i\le k-1,\\
        \calL_{1,k}^{-1}(x) &\tn{if  }i= k ,\\
        a_i &\tn{if  }i\ge k+1 
    \end{cases}
        \quad \tn{and}\quad 
        \phi_j(y)=\begin{cases}
        d_j &\tn{if  } j \le \ell-1,\\
        \calL_{2,\ell}^{-1}(y) &\tn{if  } j=\ell,\\
        c_j &\tn{if  } j\ge \ell+1.
    \end{cases}
\end{equation*}
It follows that 
\begin{equation*}
    \begin{split}
          &\phantom{=\;}\sum_{i=1}^n f_{i,0}\Big( \psi_i(x), \, y\Big)
        = \sum_{i=1}^n f\Big( \calL_{1,i}\circ \psi_i(x), \, y\Big)
        \\
        &=
        \sum_{i=1}^{k-1} f\Big( \calL_{1,i}\circ \psi_i(x), \, y\Big)
        +f\Big( \calL_{1,k}\circ \psi_k(x), \, y\Big)
        +
        \sum_{i=k+1}^n f\Big( \calL_{1,i}\circ \psi_i(x), \, y \Big)
        \\
        &=
        \sum_{i=1}^{k-1} f\Big( \calL_{1,i}(b_i), \, y\Big)
        +f\Big( \calL_{1,k}\circ \calL_{1,k}^{-1}(x), \, y\Big)
        +
        \sum_{i=k+1}^n f\Big( \calL_{1,i}(a_i), \, y\Big)
        \\
        &=
        \sum_{i=1}^{k-1} f( x_i, \, y)
        +f( x, \, y)
        +
        \sum_{i=k+1}^n f( x_{i-1}, \, y)
=
        f( x, \, y)
        +\sum_{i=1}^{n-1} f( x_i, \, y),
    \end{split}
\end{equation*}
implying 
\begin{equation*}
%    \begin{split}          
        f( x, \, y)=\sum_{i=1}^n f_{i,0}\Big( \psi_i(x), \, y\Big)-  \sum_{i=1}^{n-1} f( x_i, \, y).
%    \end{split}
\end{equation*}
For each \(i\), using the one-dimensional decomposition technique described in Section~\ref{sec:decomposition:1D}, we find the decompositions for 
\(f_{i,0}\big( \psi_i(x), \, y\big)\) and \(f( x_i, \, y)\).
We have 
% \begin{equation*}
%     \begin{split}
%           &\phantom{=\;}\sum_{j=1}^m f_{i,j}\Big( \psi_i(x), \, \phi_j(y)\Big)
%         = \sum_{j=1}^m f_{i,0}\Big(\psi_i(x),\, \calL_{2,j}\circ \phi_j(y)\Big)\\
%         &=
%         \sum_{j=1}^{\ell-1} f_{i,0}\Big(\psi_i(x),\, \calL_{2,j}\circ \phi_j(y)\Big)
%         +f_{i,0}\Big(\psi_i(x),\, \calL_{2,\ell }\circ \phi_\ell(y)\Big)
%         +
%         \sum_{j=\ell+1}^m f_{i,0}\Big(\psi_i(x),\, \calL_{2,j}(\circ \phi_j(y))\Big)
%         \\
%         &=
%         \sum_{j=1}^{\ell-1} f_{i,0}\Big(\psi_i(x),\, \calL_{2,j}(d_j)\Big)
%         +f_{i,0}\Big(\psi_i(x),\, \calL_{2,\ell }\circ \calL_{2,\ell }^{-1}(y)\Big)
%         +
%         \sum_{j=\ell +1}^m f_{i,0}\Big(\psi_i(x),\, \calL_{2,j}(c_j)\Big)
%         \\
%         &=
%         \sum_{j=1}^{\ell-1} f_{i,0}\Big(\psi_i(x),\, y_j\Big)
%         +f_{i,0}\Big(\psi_i(x),\, y\Big)
%         +
%          \sum_{j=\ell+1}^m f_{i,0}\Big(\psi_i(x),\, y_{j-1}\Big)\\
%         &=f_{i,0}\Big(\psi_i(x),\, y\Big)
%         +
%         \sum_{j=1}^{m-1} f_{i,0}\Big(\psi_i(x),\, y_{j-1}\Big),
%     \end{split}
% \end{equation*}
% {
% \fontsize{8.7pt}{10pt}\selectfont 
% \footnotesize
{\small\begin{equation*}
    \begin{split}
          &\phantom{=\;}\sum_{j=1}^m f_{i,j}\Big( \psi_i(x), \, \phi_j(y)\Big)
        = \sum_{j=1}^m f_{i,0}\Big(\psi_i(x),\, \calL_{2,j}\circ \phi_j(y)\Big)\\
        &=
        \sum_{j=1}^{\ell-1} f_{i,0}\Big(\psi_i(x),\, \calL_{2,j}\circ \phi_j(y)\Big)
        +f_{i,0}\Big(\psi_i(x),\, \calL_{2,\ell }\circ \phi_\ell(y)\Big)
        +      \sum_{j=\ell+1}^m f_{i,0}\Big(\psi_i(x),\, \calL_{2,j}\circ \phi_j(y)\Big)
        \\
        &=
        \sum_{j=1}^{\ell-1} f_{i,0}\Big(\psi_i(x),\, \calL_{2,j}(d_j)\Big)
        +f_{i,0}\Big(\psi_i(x),\, \calL_{2,\ell }\circ \calL_{2,\ell }^{-1}(y)\Big)
        +
        \sum_{j=\ell +1}^m f_{i,0}\Big(\psi_i(x),\, \calL_{2,j}(c_j)\Big)
        \\
        &=
        \sum_{j=1}^{\ell-1} f_{i,0}\Big(\psi_i(x),\, y_j\Big)
        +f_{i,0}\Big(\psi_i(x),\, y\Big)
        +
         \sum_{j=\ell+1}^m f_{i,0}\Big(\psi_i(x),\, y_{j-1}\Big)
         \\&=f_{i,0}\Big(\psi_i(x),\, y\Big)
        +
        \sum_{j=1}^{m-1} f_{i,0}\Big(\psi_i(x),\, y_{j-1}\Big),
    \end{split}
\end{equation*}}implying 
\begin{equation}\label{eq:fi0:psiixy}
    \begin{split}
          f_{i,0}\Big(\psi_i(x),\, y\Big)
        =\sum_{j=1}^m f_{i,j}\Big( \psi_i(x), \, \phi_j(y)\Big)-
        \sum_{j=1}^{m-1} f_{i,0}\Big(\psi_i(x),\, y_{j}\Big).
    \end{split}
\end{equation}
Moreover, 
 \begin{equation*}
    \begin{split}
          &\phantom{=\;}\sum_{j=1}^m f_{0,j}\Big( x_i, \, \phi_j(y)\Big)
        = \sum_{j=1}^m f \Big(x_i,\, \calL_{2,j}\circ \phi_j(y)\Big)\\
        &=
        \sum_{j=1}^{\ell-1} f \Big(x_i,\, \calL_{2,j}\circ \phi_j(y)\Big)
        +f \Big(x_i,\, \calL_{2,\ell}\circ \phi_\ell(y)\Big)
        +
         \sum_{j=\ell+1}^mf \Big(x_i,\, \calL_{2,j}\circ \phi_j(y)\Big)
        \\
        &=
        \sum_{j=1}^{\ell-1} f \Big(x_i,\, \calL_{2,j}(d_j)\Big)
        +f \Big(x_i,\, \calL_{2,\ell }\circ \calL_{2,\ell }^{-1}(y)\Big)
        +
        \sum_{j=\ell+1}^m f \Big(x_i,\, \calL_{2,j}(c_j)\Big)
        \\
        &=
        \sum_{j=1}^{\ell-1} f (x_i,\, y_j)
        +f (x_i,\, y)
        +
        \sum_{j=\ell+1}^m f (x_i,\, y_{j-1})        =f (x_i,\, y)
        +
        \sum_{j=1}^{m-1} f (x_i,\, y_{j}),
    \end{split}
\end{equation*}
implying 
\begin{equation}\label{eq:fxiy}
    \begin{split}
          f (x_i,\, y)
        =\sum_{j=1}^m f_{0,j}\Big( x_i, \, \phi_j(y)\Big)-
        \sum_{j=1}^{m-1} f (x_i,\, y_{j}).
    \end{split}
\end{equation}
Therefore, for any $(x,y)\in [x_{k-1}, x_k]\times [y_{\ell-1}, y_\ell]$, by \eqref{eq:fi0:psiixy}  and \eqref{eq:fxiy}, we have
% \begin{equation*}
%     \begin{split}          
%         f( x, \, y)
%         &=\sum_{i=1}^n f_{i,0}\Big( \psi_i(x), \, y\Big)-  \sum_{i=1}^{n-1} f( x_i, \, y)
%         \\& =\sum_{i=1}^n 
%         \Bigg(\sum_{j=1}^m f_{i,j}\Big( \psi_i(x), \, \phi_j(y)\Big)-
%         \sum_{j=1}^{m-1} f_{i,0}\Big(\psi_i(x),\, y_{j}\Big)\Bigg)
%         \\ & \quad\   -  \sum_{i=1}^{n-1}\Bigg(\sum_{j=1}^m f_{0,j}\Big( x_i, \, \phi_j(y)\Big)-
%         \sum_{j=1}^{m-1} f (x_i,\, y_{j})\Bigg)
%         \\& =\sum_{i=1}^n \sum_{j=1}^m f_{i,j}\Big( \psi_i(x), \, \phi_j(y)\Big)-
%         \sum_{i=1}^n\sum_{j=1}^{m-1} f_{i,0}\Big(\psi_i(x),\, y_{j}\Big)
%         \\ & \quad\   -  \sum_{i=1}^{n-1} \sum_{j=1}^m f_{0,j}\Big( x_i, \, \phi_j(y)\Big)+
%         \sum_{i=1}^{n-1}\sum_{j=1}^{m-1} f (x_i,\, y_{j}).
%     \end{split}
% \end{equation*}
\begin{equation*}
    \begin{split}          
        % &\phantom{=\;\;}
        f( x, \, y)
        &=\sum_{i=1}^n f_{i,0}\Big( \psi_i(x), \, y\Big)-  \sum_{i=1}^{n-1} f( x_i, \, y)
        % \\& =\sum_{i=1}^n 
        % \Bigg(\sum_{j=1}^m f_{i,j}\Big( \psi_i(x), \, \phi_j(y)\Big)-
        % \sum_{j=1}^{m-1} f_{i,0}\Big(\psi_i(x),\, y_{j}\Big)\Bigg)
        %   -  \sum_{i=1}^{n-1}\Bigg(\sum_{j=1}^m f_{0,j}\Big( x_i, \, \phi_j(y)\Big)-
        % \sum_{j=1}^{m-1} f (x_i,\, y_{j})\Bigg)
        \\& =\sum_{i=1}^n \sum_{j=1}^m f_{i,j}\Big( \psi_i(x), \, \phi_j(y)\Big)-
        \sum_{i=1}^n\sum_{j=1}^{m-1} f_{i,0}\Big(\psi_i(x),\, y_{j}\Big)
           \\& \quad\  -  \sum_{i=1}^{n-1} \sum_{j=1}^m f_{0,j}\Big( x_i, \, \phi_j(y)\Big)+
        \sum_{i=1}^{n-1}\sum_{j=1}^{m-1} f (x_i,\, y_{j}).
    \end{split}
\end{equation*}
Since $k$ and $\ell$ are arbitrary, the above equation holds for all
$(x,y)=\cup_{k=1}^n\cup_{\ell=1}^m[x_{k-1}, x_k]\times [y_{\ell-1}, y_\ell]= [x_0,x_n]\times [y_0,y_m]$.
\end{proof}

% to decompose  $f(x,y)$ as  $f(x,y)=\sum_{i=1}^n f_i(\phi_i(x),y)$


% {\small
% \begin{equation*}
%     \begin{split}
%         &\quad \sum_{j=1}^m\sum_{i=1}^n f_{i,j}\Big( \psi_i(x), \phi_j(y)\Big)
%         =\sum_{j=1}^m\sum_{i=1}^n f\Big( \calL_{1,i}\circ \psi_i(x), \calL_{2,j}\circ \phi_j(y)\Big)\\
%         &=\sum_{j=1}^m
%         \Bigg(
%         \sum_{i=1}^{k-1} f\Big( \calL_{1,i}\circ \psi_i(x), \calL_{2,j}\circ \phi_j(y)\Big)
%         +f\Big( \calL_{1,k}\circ \psi_k(x), \calL_{2,j}\circ \phi_j(y)\Big)
%         +
%         \sum_{i=k+1}^n f\Big( \calL_{1,i}\circ \psi_i(x), \calL_{2,j}\circ \phi_j(y)\Big)
%         \Bigg)\\
%         &=\sum_{j=1}^m
%         \Bigg(
%         \sum_{i=1}^{k-1} f\Big( \calL_{1,i}(b_i), \calL_{2,j}\circ \phi_j(y)\Big)
%         +f\Big( \calL_{1,k}\circ \calL_{1,k}^{-1}(x), \calL_{2,j}\circ \phi_j(y)\Big)
%         +
%         \sum_{i=k+1}^n f\Big( \calL_{1,i}(a_i), \calL_{2,j}\circ \phi_j(y)\Big)
%         \Bigg)\\
%         &=\sum_{j=1}^m
%         \Bigg(
%         \sum_{i=1}^{k-1} f\Big( x_i, \calL_{2,j}\circ \phi_j(y)\Big)
%         +f\Big( x, \calL_{2,j}\circ \phi_j(y)\Big)
%         +
%         \sum_{i=k+1}^n f\Big( x_{i-1}, \calL_{2,j}\circ \phi_j(y)\Big)
%         \Bigg)\\
%         &=\sum_{j=1}^m
%         \Bigg(
%         f\Big( x, \calL_{2,j}\circ \phi_j(y)\Big)
%         +\sum_{i=1}^{n-1} f\Big( x_i, \calL_{2,j}\circ \phi_j(y)\Big)
%         \Bigg)\\
%         &=\sum_{j=1}^{\ell-1}
%         \Bigg(
%         f\Big( x, \calL_{2,j}\circ \phi_j(y)\Big)
%         +\sum_{i=1}^{n-1} f\Big( x_i, \calL_{2,j}\circ \phi_j(y)\Big)
%         \Bigg)
%         +
%         \Bigg(f\Big( x, \calL_{2,\ell}\circ \phi_\ell(y)\Big)
%         +\sum_{i=1}^{n-1} f\Big( x_i, \calL_{2,\ell}\circ \phi_\ell(y)\Big)\Bigg)
%         +\sum_{\ell=\ell-1}^m
%         \Bigg(
%         f\Big( x, \calL_{2,j}\circ \phi_j(y)\Big)
%         +\sum_{i=1}^{n-1} f\Big( x_i, \calL_{2,j}\circ \phi_j(y)\Big)
%         \Bigg)
%     %     \sum_{i=1}^n f\circ \calL_i\circ \psi_i(x)
%     % &=\sum_{i=1}^{j-1} f\circ \calL_i\circ \psi_i(x)
%     % +f\circ \calL_j\circ \psi_j(x)
%     % +\sum_{i=j+1}^n f\circ \calL_i\circ \psi_i(x)
%     % \\& = \sum_{i=1}^{j-1} f\circ \calL_i(b)
%     % +f\circ \calL_j\circ\calL_j^{-1}(x) 
%     % +\sum_{i=j+1}^n f\circ \calL_i(a)
%     % \\& = \sum_{i=1}^{j-1} f (x_i)
%     % +f(x)
%     % +\sum_{i=j+1}^n f (x_{i-1})
%     % =f(x)+\underbrace{\sum_{i=1}^{n-1}f(x_i)}_{\tn{constant}}.
%     \end{split}
% \end{equation*}
% }


% \begin{equation*}
%     f(x,y)=\sum_{j=1}^m\sum_{i=1}^n f_{i,j}\Big( \psi_i(x), \phi_j(y)\Big)-\underbrace{\red{\sum_{j=1}^{m-1}\sum_{i=1}^{n-1}f(x_{i,j})}}_{\tn{constant}}.
% \end{equation*}

% \red{Add a figure to illustrate the decomposition for 2D later.}

% The decomposition here can be generalized to high dimensional cases. For example, $f:[-1,1]^2\to\R$
% \begin{equation*}
%     \bmf_1:[-1,1]^2\to[-1,1]^2, \qquad   \bmf_1=\begin{bmatrix}
%     f_{1,1}\\  f_{1,2}\end{bmatrix}
% \end{equation*}
% \begin{equation*}
%     f_2:[-1,1]^2\to\R,\quad  f_2:(u,v)\mapsto 
%     f\big(\tfrac{u+1}{2}\big)+ f\big(\tfrac{v-1}{2}\big)-f(0)
% \end{equation*}
% \begin{equation*}
% 	f_{1,1}(x,y)=\Big(\ReLU(2x)-1, y\Big)
%  % =\begin{cases}
% 	% 	-1 & \tn{for}\  x\in [-1,0),\\
% 	% 	2x-1 & \tn{for}\  x\in [0,1].
% 	% \end{cases}
% \end{equation*}
% \begin{equation*}
% 	f_{1,2}(x,y)=\Big(-\ReLU(-2x)+1,y\Big)
%  % =\begin{cases}
% 	% 	2x+1 & \tn{for}\  x\in [-1,0),\\
% 	% 	1  & \tn{for}\  x\in [0,1].
% 	% \end{cases}
% \end{equation*}
% Then for any $x\in [-1,1]$ we have 
% \begin{equation*}
% 	\begin{split}
% 	    f(x,y)&=f\Big(\tfrac{\ReLU(2x)}{2},y\Big)+ f\Big(\tfrac{-\ReLU(-2x)}{2},y\Big)-f(0,y)
%      \\& =f\Big(\tfrac{f_{1,1}(x)+1}{2},y\Big)+ f\Big(\tfrac{f_{1,2}(x)-1}{2},y\Big)-f(0,y)
%  =
%  f_2\circ \bmf_1(x,y)
% 	\end{split}
% \end{equation*}


% \subsection{Connection to Kolmogorov-Arnold superposition theorem}

% The Kolmogorov-Arnold superposition theorem (KST) \cite{Kol1957} is typically utilized to simplify multidimensional problems into one-dimensional ones. Essentially, the compositional representation in KST can be viewed as a unique form of neural network with complex activation functions that are determined by the target function, rendering KST impractical for real-world computations.

% \begin{theorem}[KST \cite{Kol1957}]
% 	\label{thm:kst}
% 	There exist continuous functions $h_{i,j}\in C([0,1])$ for $i=0,1,\cdots,2d$ and $j=1,2,\cdots,d$ such that any continuous function $f\in C([0,1]^d)$ can be represented as 
% 	\begin{equation*}
% 		f(\bmx)=\sum_{i=0}^{2d}  g_i \Big(\sum_{j=1}^d h_{i,j}(x_j)\Big)\quad \tn{for any $\bmx=[x_1,x_2,\cdots,x_d]^\ts\in [0,1]^d$,}
% 	\end{equation*}
% 	where $g_i:\R\to\R$ is a continuous function for each $i\in \{0,1,\cdots,2d\}$.
% \end{theorem}


% KST allows us to represent multidimensional functions using univariate functions, although these may be complex. By employing the 1D decomposition techniques described in Section~\ref{sec:decomposition:1D}, we can approximate each univariate function with smooth functions. By integrating all these smooth functions, we represent a continuous high-dimensional function smoothly. However, it is important to note that these are theoretical results. In practice, the univariate functions used in KST to represent continuous high-dimensional functions exhibit ``wild behavior," which limits the practical applicability of this 




% \section{Additional discussions}


% \subsection{A balanced decomposition through components and depth}

% \subsection{Characterization of difficulties of learning various functions}

% \red{Eg: hard to use the infinity norm or some integral norm of the target function to characterize the learning difficulty}
% \begin{equation*}
%     f_1(x)= \one_{\{|x|<0.02\}}\cdot\sin(50\pi x) 
%     \quad \tn{and}\quad 
%     f_2(x)=\frac{1}{50}\sin(\pi x)
% \end{equation*}
% \begin{equation*}   \|f_1\|_{L^1([-1,1])}
% =\int_{-1}^{1}|f_1(x)|dx=\int_{-0.02}^{0.02}|\sin(50\pi  x)|\frac{d(50x)}{50}
% =
% \int_{-1}^{1}|\sin(\pi  y)|\frac{dy}{50}
% =\|f_2\|_{L^1([-1,1])}
% \end{equation*}

% \red{
% We can describe the complexity of learning using the expression:
% \begin{equation*}
% \|f'\|_{L^\infty([-1,1])}^\alpha \cdot V^\beta(f)\quad \text{with appropriate $\alpha,\beta>0$},
% \end{equation*}
% where $V(f)$ represents the total variation of $f$. The values of $\alpha$ and $\beta$ should exceed 1, and their determination depends on the network architecture and learning hyperparameters. It is important to note that this function in $f$ is generally not a norm.
% }

% \subsubsection{Total variation}
% We may assume the target function $f\in C^1([0,1]^d; [0,1])$

% % Total variant 
% % \begin{equation*}
% %     V(f, \Omega) = \int_{\Omega} \|\nabla f(\bmx)\|_2 \, \mathrm{d}\bmx
% % \end{equation*}


% \begin{definition}
%  Let \( \Omega \) be an open subset of \( \mathbb{R}^n \). Given a function \( f \) belonging to \( L^1(\Omega) \), the total variation of \( f \) in \( \Omega \) is defined as
% \[
% V(f, \Omega) \coloneqq \sup \left\{ \int_{\Omega} f(\bmx) \operatorname{div} \phi (\bmx) \, \mathrm{d}\bmx : \phi \in C_c^{1}(\Omega, \mathbb{R}^n), \ \|\phi\|_{L^{\infty}(\Omega)} \leq 1 \right\},
% \]
% where
% \begin{itemize}
%   \item \( C_c^{1}(\Omega, \mathbb{R}^n) \) is the set of continuously differentiable vector functions of compact support contained in \( \Omega \),
%   \item \( \|\cdot\|_{L^{\infty}(\Omega)} \) is the essential supremum norm, and
%   \item \( \operatorname{div} \) is the divergence operator.
% \end{itemize}
% This definition does not require that the domain \( \Omega \subseteq \mathbb{R}^n \) of the given function be a bounded set.
% \end{definition}


% \begin{theorem}
%  Given a \( C^{1}(\overline{\Omega}) \) function \( f \) defined on a bounded open set \( \Omega \subseteq \mathbb{R}^{n} \), with \( \partial \Omega \) of class \( C^{1} \), the total variation of \( f \) has the following expression:
% \[
% V(f, \Omega) = \int_{\Omega} \|\nabla f(\bmx)\|_2 \, \mathrm{d}\bmx .
% \]    
% \end{theorem}

% \subsubsection{Learning difficulty characterization}



% Given a point $\bmx\in\Omega$, learning difficulty of $f$ at $\bmx$ is characterized by 
% \begin{equation*}
%     \int_{\Omega\cup B(\bmx,r)} \|\nabla f(\bmx)\|_2 \, \mathrm{d}\bmx,
% \end{equation*}
% where $B(\bmx,r)$ is a ball with center $\bmx$ and radius $r$ with $r^{4}\sim $ learning rate.


% % \newpage
% Write $V(f,\Omega)=V(f)$ in the case $\Omega=[0,1]^d$.





% \red{I believe that we can employ $V(f)$ or \blue{$V_\alpha (f, \Omega) = \int_{\Omega} \|\nabla f(\bmx)\|_2^\alpha \, \mathrm{d}\bmx$ with proper $\alpha>0$} to represent the complexity involved in approximating the function $f$, belonging to $C^1([0,1]^d; [0,1])$, utilizing $\bmh$ with the specified (above) architecture. For shallow networks $\bmh$, a comparable result is not available.}
% % \red{Use $V(f)$ to describe the complexity involved in approximating the function $f\in  C^1([0,1]^d; [0,1])$ using $\bmh$ with the architecture above. In the case of shallow netowrks $\bmh$, we do not have a similar result.}









% % \section{Theoretical analysis}

% % \section{Further discussions}

% % \section{Additional Discussions}

% % \subsection{}



% \subsection{Low-pass filter nature of shallow networks}

% \subsection{Overcome low-pass filter nature of shallow networks by compositions}
% A few bumps

% \subsection{Efficient deep network architectures learning complicated functions}

% {\color{teal}{
% We explain the initial stage of training. The composition structure of the network
% \begin{equation}
%     \bmh = \bmh_m\circ \cdots \circ \bmh_2\circ \bmh_1 
% \end{equation}
% Let $\bmh_1 = \bmh_1(\bmx; \theta)$, then the gradient of the loss function $\mathcal L (f, \bmh) = \int_{\Omega} |f - \bmh|^2 d\bmx$ is (not rigorous)
% \begin{equation*}
%     \dfrac{\partial \mathcal L}{\partial\theta} =  \nabla \bmh_m \nabla \bmh_{m-1} \cdots \frac{\partial}{\partial \theta}\bmh_{1} 
% \end{equation*}

% The gradient will be zero immediately once the evaluation of $\bmh_k$ is negative for $1\le k\le m$. For large rank and width, the statistical behavior is dominating. Each product truncates away roughly half of the possibility of being nonzero of the gradient. Having a smaller rank provides randomness if $m$ is not large.

% To fulfill the feature-by-feature training dynamics, the localized components (on both sides) in layer $\bmh_1$ must have nonzero gradients. Thus when training a highly oscillatory function of a frequency bandwidth $\nu$, the width of $\bmh_1$ should be at the order of $2^m \nu$? A significantly smaller width could suffer from vanishing gradients.

% \begin{remark}
%     It seems the rank can leverage the vanishing gradient effect, while a large rank would fall into the statistical regime where the coefficients are less useful. Ideally, there is competition between the two: vanishing gradient vs. slow speed.
% \end{remark}
% }}



% \subsection{Importance of rank}
% In this section, we demonstrate the importance of the rank. The simplest construction of the network is using $\texttt{rank}$-1 structure. The three-layer $\texttt{rank}$-1 network can be viewed as function composition. Once the composition consists of ``smoother'' functions, then it may be possible to boost the training speed over shallow network.
% \begin{problem}
% Given a smooth function $h: [0, 1]\to \mathbb{R}$, find ``smoother'' functions $f: I\to \mathbb{R}$ and $g: [0, 1]\to I$ such that $h(x) = f(g(x))$.
% \end{problem}
% Here the terminology ``smoother'' is still a vague description. Let's narrow down a possible definition for it. Define the critical set of $h$ by  $\mathcal{C}_h = \{z\in [0,1]\mid h'(z) = 0\}$, then using chain rule that $h'(z) = f'(g(z)) g'(z)$, we introduce the subsets 
% \begin{equation}
% \begin{aligned}
%     \mathcal{C}_f &=  \{z\in [0,1]\mid f'(g(z)) = 0\}\subset \mathcal C_h, \\
%     \mathcal{C}_g &=  \{z\in [0,1]\mid g'(z) = 0\}\subset \mathcal C_h,
% \end{aligned}
% \end{equation}
% then we define the smoothness index for $h$,  $f$ and $g$ by 
% \begin{equation}
% \begin{aligned}
%     \kappa_h &= \min_{z_1\neq z_2\in \mathcal{C}_h} |z_1 - z_2|, \\
%     \kappa_f &= \min_{z_1\neq z_2\in \mathcal{C}_f} |z_1 - z_2|, \\
%     \kappa_g &= \min_{z_1\neq z_2\in \mathcal{C}_g} |z_1 - z_2|,
% \end{aligned}
% \end{equation}
% and ``smoother'' means $\kappa_h < \kappa_f$ and $\kappa_h < \kappa_g$. Clearly the smoothness index only cares about the fastest oscillation.
% \subsection{Existence of composition}
% Clearly, when $h$ is a monotone function, such a composition exists since $\mathcal{C}_h = \emptyset$. Next, we consider $h$ as an oscillatory function such that all critical points are not degenerate, which means $h''(z) \neq 0$ for all $z\in \mathcal{C}_h$. 

% \begin{lemma}
%     $\mathcal{C}_f \cap \mathcal {C}_g = \emptyset$.
% \end{lemma}
% \begin{proof}
%     If $z\in\mathcal{C}_f \cap \mathcal {C}_g $, then $h''(z) = f''(g(z)) |g'(z)|^2 + f'(g(z)) g''(z) = 0$. It is a contradiction of the assumption.
% \end{proof}
% \begin{theorem}
%     If $\mathrm{card}(\mathcal{C}_h) \ge 4$ and h has distinct values on $\mathcal{C}_h$, then it is impossible to find ``smoother'' $f$ and $g$.
% \end{theorem}
% \begin{proof}
%     We prove by construction. Let $z_0, z_1, z_2, z_3$ be four consecutive critical points of $h$,  without loss of generality, we assume $z_0 < z_1 < z_2 < z_3$ and $h(z_i)$ are distinct, $i=0,1,2,3$. If $\kappa_h > \kappa_f$ and $\kappa_h > \kappa_g$, then there are two possible configurations. 
%     \begin{enumerate}
%         \item $z_0, z_2\in \mathcal{C}_f$, $z_1, z_3\in \mathcal{C}_g$. We focus on the interval $(z_0, z_2)$, since $f'(g(z))$ has a fixed sign on $(z_0, z_2)$ and $g'(z) \neq 0$ at $z = z_0$ or $z = z_2$, it implies the signs of $g'(z)$ on $(z_0, z_1)$ and $(z_1, z_2)$ must differ. If $g' > 0$ on $(z_0, z_1)$,  we let $I_1 = [g(z_0), g(z_1)]$ and $I_2 = [g(z_2), g(z_1)]$, otherwise we let $I_1 = [g(z_1), g(z_0)]$ and $I_2 = [g(z_1), g(z_2)]$.  

%         Note that we must have $I_1\subset I_2$, otherwise there exists a point $z_3\in (z_1, z_2)$ that $g(z_3) = g(z_0)$, which causes $f'(g(z_3)) = 0$, which is a contradiction. Similarly, we must have $I_2\subset I_1$, which implies we must have $I_1 = I_2$, that is $g(z_2) = g(z_0)$ hence $h(z_0) = h(z_2)$, which is a contradiction with our assumption.

%         \item $z_0, z_2\in \mathcal{C}_g$, $z_1, z_3\in \mathcal{C}_f$. We simply reuse the above argument on the interval $(z_1, z_3)$ instead.
%     \end{enumerate}
%     Therefore, once $h$ has at least four critical points with distinct values, it is impossible to construct the composition.  
% \end{proof}

% The above theorem shows that if $h$ is a general oscillatory function (contains at least four critical points with distinct evaluations) then in general the set $\mathcal C_f$ and $\mathcal C_g$ must contain consecutive critical points from $\mathcal C_h$. Therefore having a deeper network with $\texttt{rank}$-1 network structure cannot distribute the fast oscillation into slower ones to different layers. That is why we need at least rank two. The following example shows it is possible to find $f$ and $g$ when the ``smoother'' requirement is lifted.



% \newpage
% \red{
% \subsection{Learning}
% SGD with momentum is much better than SGD without momentum


% \subsection{Example}



% Use Fourier series to approximate the target function. 

% Smooth,  coefficients small for large $n$, $\cos(n\pi x)$ or $\sin(n\pi x)$.


% \subsection{To do}
% \textbf{local example}

% FCNN, low-rank, high-rank
% }


% \red{
% \subsection{
% Dual-phase learning approach}

% \begin{equation*}
%     f(x)=(1-x^2)\sin(20\pi x)
% \end{equation*}
% Learning such a target function
% $f$ is challenging, as previously mentioned. To bridge this gap, we suggest a dual-phase learning strategy, grounded in the observation that if a neural network (NN) can learn a local feature effectively, it can then rapidly acquire other features.

% Our method involves two stages:

% \begin{itemize}
%     \item Initial Stage: This focuses on identifying a basic feature by implementing a mask ($\one_{\{|x|<0.1\}}$).
% \item
% Subsequent Stage: In this phase, we proceed with standard training, this time without the mask.
% \end{itemize}

% See GIFs for illustrations (200 epochs in total; our dual-phase method using mask for first 60 epochs): 
% \blue{output\_idx0.gif} for the normal method; 
% \blue{output\_idx1.gif} for our dual-phase method.

% \bigskip
% We want to note that although the mask $\exp(-100x^2)$ is functional, it does not perform as efficiently as $\one_{\{|x|<0.1\}}$.

% To test the effectiveness of the new method on complex functions, we experiment with 
% $f(x)=(1-x^2)\sin(40\pi x)$. Initially, we use a mask for the first 100 epochs, followed by 400 epochs without a mask ($\one_{\{|x|<0.05\}}$). See \blue{output\_idx2.gif} for visual illustrations.
% }


% \clearpage
% \section{Understanding How Narrow yet Deep Networks Grasp High-Frequency Features}
% % \section{Why narrow but deep network can capture high frequency}

% % \section{Understanding the Ability of Narrow and Deep Networks to Capture High-Frequency Components}
% width 10, depth 6, rank 5, can learn one period of $\sin(50\pi x)$ well


% \section{Networks with Two Hidden Layers Have a Higher Probability of Effective Learning}
% Rashomon set

% two hidden layer easier, 
% first hidden layer width 2, second hidden layer width 1000

% \section{High dim eg; Name of architecture}

% data samples focus on high frequency

% optimization




% \clearpage
% \section{Experiments}
%  comparison: 1) w=1; b=0; trivial w and b;
%  2) random w and b; 3) fix all w, b, a, use least square

%  connecting to our previous paper; composition of many ``shallow networks"

% \begin{itemize}
% \item New architecture, focus on high frequency functions
%     \item   adaptive architecture, target function
%     \item  fix w, b; train a; easy and efficient
%     \item  Start from boundary, 
% \end{itemize}

% \begin{equation*}
%     \bmphi
% \end{equation*}
% \begin{enumerate}[(i)]
%     \item  Deep network of small width can capture local bump; 1) monotone, arctan(100x)
%     2) 2D3D
%     \item  Our architecture can learn global high frequency, e.g., $\sin(50\pi x)$, 2D $\sum_{i=1}^4 \cos \sin$; 2D star;
%      3D
%      \item data local sampling
%      \item Learn from boundary, 
% \end{enumerate}


%%%%%%%%%%%%%%%%%%%%%%%%%%%%%%%
\subsection{Related work}
\label{sec:related:work}
Several lines of research are closely related to this work, including approximation theory, low-rank methods, random feature models, and architectures inspired by the Kolmogorov–Arnold representation.
\paragraph{Approximation}
Extensive research has examined the approximation capabilities of neural networks, focusing on various architectures to approximate diverse target functions. Early studies concentrated on the universal approximation power of single-hidden-layer networks~\cite{Cybenko1989ApproximationBS,HORNIK1991251,HORNIK1989359}, which demonstrated that sufficiently large neural networks could approximate specific functions with arbitrary precision mathematically, without quantifying the error relative to network size. Subsequent research, such as~\cite{yarotsky18a,yarotsky2017,doi:10.1137/18M118709X,ZHOU2019,10.3389/fams.2018.00014,2019arXiv190501208G,2019arXiv190207896G,MO,shijun:NonlineArpprox,shijun:Characterized:by:Numer:Neurons,shijun:smooth:functions,shijun:arbitrary:error:with:fixed:size,shijun:thesis,shijun:intrinsic:parameters,shijun:RCNet,JMLR:v25:23-0912,shijun:net:arc:beyond:width:depth,shijun:optimal:rate:in:width:and:depth}, analyzed the approximation error for different networks in terms of size characterized by width, depth, or the number of parameters. 
%However, these studies focused primarily on mathematical construction and theory, with little attention to numerical errors and effective learning methods. This work introduces a new network architecture based on the ``low pass filter" nature of shallow networks in terms of both representation and learning dynamics studied in~\cite{ZZZZ-23}. The balanced structure of MMNN and the corresponding learning strategy leads to a smooth decomposition which is conducive to an effective training process.
% \zhou{Those studies focused primarily on the mathematical constructions that lead to the existence theory for such neural networks. Little attention has been given to obtaining the parameters in the neural networks, and the numerical errors, especially those caused by finite precision in computer simulations. This motivated us to conduct the present investigation considering attainable training processes and numerical errors. More precisely, the balanced structure of MMNN, the selection of training parameters, and the corresponding learning strategy considered here are designed to create a smooth partition of the function which is conducive to an effective training process.} 
Those studies have primarily concentrated on the mathematical theory that supports the existence theory for such neural networks. However, there has been limited focus on determining the parameters within these networks computationally and the numerical errors, particularly those arising from finite precision in computer simulations. This gap motivated our current investigation, which considers practical training processes and numerical errors. Specifically, the balanced structure of MMNN, the choice of training parameters, and the associated learning strategy discussed here are intended to facilitate a smooth decomposition of the function, thereby promoting an efficient training process.

% \paragraph{Low-rank methods }
% Low-rank structures in $\bmW$ of a fully connected neural network have been investigated by a variety of groups.
% % which is equivalent to pursuing a linear dimension reduction between layers.
% \zhou{For example, the methods proposed in \cite{2022arXiv220913569R,9157223,6638949} focus on accelerating training and reducing memory requirements while maintaining final performance. The low-rank idea is generalized to tensor train decomposition in~\cite{novikov2015tensorizing}.}
% %However this may not be effective in a nonlinear representation using a deep neural network. For example, the low-rank methods proposed in \cite{2022arXiv220913569R,9157223,6638949} focus on accelerating training and reducing memory requirements while maintaining final performance. 
% \zhou{This is different from the MMNN proposed here in two aspects. First there are two matrices in each layer, $\bmA$ outside and $\bmW$ inside the activation functions. Each row of $\bmA$ is the weights to a linear combination of a set of random basis functions which forms a component in each layer. The number of rows of $\bmA$, which is the number of components, are selected depending on the complexity of the function and is usually much smaller than the number of columns, which corresponds to the number of basis functions. Each row of $(\bmW, \bmb)$ corresponds to a random parametrization of a basis function. The number of rows of $\bmW$ corresponds to the number of basis functions which is usually much larger than the number of columns of $\bmW$ which is the input dimension. Secondly, in our MMNN, only $\bmA$ is trained while $\bmW$ is fixed with randomly initialized values. The theoretical study and numerical experiments demonstrate that the architecture of MMNN, combined with the learning strategy, is effective in approximating complicate functions. 
% %Our key contribution is the design of a balanced multi-component and multi-layer structure to decompose the complexity of the target function that is conducive to approximation and learning. We theoretically and numerically demonstrate that our architecture, combined with our learning strategy, is effective in approximating complex functions. 
% }

\paragraph{Low-rank methods}
Low-rank structures in the weight matrix $\bmW$ of a fully connected neural network have been investigated by various groups. For example, the methods proposed in \cite{2022arXiv220913569R,9157223,6638949} focus on accelerating training and reducing memory requirements while maintaining final performance. The concept of low-rank structures is further extended to tensor train decomposition in~\cite{novikov2015tensorizing}.
The MMNN proposed here differs in two key aspects. First, each layer contains two matrices: $\bmA$ outside and $\bmW$ inside the activation functions. Each row of $\bmA$ represents the weights for a linear combination of a set of random basis functions, forming a component in each layer. The number of rows in $\bmA$, which equals the number of components, is selected based on the complexity of the function and is typically much smaller than the number of columns, corresponding to the number of basis functions. Each row of $(\bmW, \bmb)$ represents a random parameterization of a basis function, with the number of rows in $\bmW$ corresponding to the number of basis functions, usually much larger than the number of columns in $\bmW$, which is the input dimension.
Secondly, in our MMNN, only $\bmA$ is trained while $\bmW$ remains fixed with randomly initialized values. Theoretical studies and numerical experiments demonstrate that the architecture of MMNN, combined with the learning strategy, is effective in approximating complex functions.


% \paragraph{Random features} \zhou{Fixing $(\bmW, \bmb)$ in the MMNN is inspired by an earlier approach called Random features \cite{NIPS2007_013a006f,NIPS2016_e7061188,9495136} (please add citations here). Different from the existing random feature methods, in which only the linear combination parameters at the output layer is learned, we train the matrices $\bmA$ in each layer as well. On one hand, this increases the computational cost over the random feature methods, although it is still less expensive than training a network with $\bmA$ and $(\bmW, \bmb)$ as learnable parameters. On the other hand, the approximation power is improved over the random feature methods, yet the training process is not as difficult as training a standard fully connected networks of the same size. Our experiments seem to indicate that this change provides a reasonable balance between approximation accuracy and the training cost. 
% }

\paragraph{Random features} 
Fixing $(\bmW, \bmb)$ of each layer and use of random basis functions in the MMNNs is inspired by a previous approach known as random features \cite{NIPS2007_013a006f,NIPS2016_e7061188,9495136,NIPS2017_61b1fb3f,peng2021random}. In typical random feature methods, only the linear combination parameters at the output layer are trained which also leads to the issue of ill-conditioning of the representation. While in MMNNS matrix $\bmA$ and vector $\bmc$ of each layer are trained. 
% Although this approach increases the computational cost compared to random feature methods, it is still less expensive than training a network where both $(\bmA,\bmc)$ and $(\bmW, \bmb)$ are learnable parameters. Additionally,
Our MMNN employs a composition architecture and learning mechanism that enhances the approximation capabilities compared to random feature methods while achieving a more effective training process than a standard fully connected network of equivalent size. Extensive experiments demonstrate that our approach can strike a satisfactory balance between approximation accuracy and training cost.


% \paragraph{Komogolrov-Arnold (KA) representation:} \zhou{KA representation states a multivariate continuous function on a bounded domain can be written as a finite composition of continuous functions of a single variable and the binary operation of addition. Despite of the possibility of encountering non-smooth or even fractal univariate functions in the KA representation, it provides an appealing framework to build neural network structures to approximate functions especially in higher dimensions. This idea has been explored in a number of studies \cite{2024arXiv240419756L} (please add citations). For example, a recent proposed network called KA network (KAN) uses spline functions to approximate the univariate functions in the KA representation. Aimed at approximating highly oscillatory functions, the proposed MMNN is based on the idea of ``divide and conquer'', which uses different network architectures, activation functions, and training cost than those KA representation based methods.}
%Recently, a network structure inspired by the Kolmogorov-Arnold (KA) representation, called KAN, is proposed. KA representation states a multivariate continuous function on a bounded domain can be written as a finite composition of continuous functions of a single variable and the binary operation of addition. The key idea of KAN is to use spline functions to approximate those univariate functions. However, those univariate functions may be highly non-smooth or even fractal which makes splines difficult to approximate well. Hence a network that is wider and deeper than the compact form of KA representation is needed in practice.  Moreover, the coordinate system may be artificial and simple additions of different univariate functions may not be effective in approximating a generic multi-variate function. Also the training process seems relative expensive. 

\paragraph{Komogolrov-Arnold (KA) representation} 
The KA representation theorem \cite{kolmogorov1957} states that any multivariate continuous function on a hypercube can be expressed as a finite composition of continuous univariate functions and the binary operation of addition. However, this elegant mathematical representation may result in compositions of non-smooth or even fractal univariate functions in general, a computational challenge one has to address in practice. 
%Although this representation may result in non-smooth or even fractal univariate functions, it offers a compelling framework for constructing neural network structures to approximate functions, particularly in higher dimensions. 
KA representation has been explored in several studies \cite{2024arXiv240419756L,MAIOROV199981,shijun:arbitrary:error:with:fixed:size,ISMAYILOVA2024106333}. A recently proposed network known as the KA network (KAN) utilizes spline functions to approximate the univariate functions in the KA representation. The proposed MMNN is motivated by a multi-component and multi-layer smooth decomposition, or a ``divide and conquer" approach, employing distinct network architectures, activation functions, and training strategies.



%%%%%%%%%%%%%%%%%%%%%%%%%%%%%%%%%%%%%%
\section{Numerical experiments}
\label{sec:experiments}


We perform extensive experiments to validate our analysis and demonstrate the effectiveness of MMNNs through multi-component and multi-layer decomposition studied in Section~\ref{sec:decomposition}. In particular, our tests show its ability in (1) adaptively capturing localized high-frequency features in Section~\ref{sec:local}, (2) approximating highly oscillatory functions in Section~\ref{sec:oscillatory}, 
(3) approximating discontinuous functions with porous structures in Section~\ref{sec:discontinuous:func:porous}, (4) some interesting learning dynamics in Section~\ref{sec:dynamics}, and 
(5) solving problems in three and higher dimensions in Section~\ref{sec:higher}.
All our experiments involve target functions that include high-frequency components in various ways and are difficult to handle by shallow networks (no matter how wide) as shown in our previous work \cite{ZZZZ-23}. Moreover, our experience on these tests shows that using a fully connected deep neural network would require many more parameters and is much harder (if possible) to train to get a comparable result. This is mainly due to a balanced and structured network design of MMNN in terms of (1) the network width $w$, which is the number of hidden neurons or random basis functions in each component, (2) the rank $r$, which is the number of components in each layer, and (3) the network depth $l$, which is the number of layers in the network. The use of a controllable number of collective components (through $\bmA$) in each layer instead of a large number of individual neurons and the use of fixed and randomly chosen weights $(\bmW,\bmb)$ make the training process more effective.

In all tests, (1) data are sampled enough to resolve fine features in the target function, (2) the Adam optimizer is used in training, (3) the mean squared error (MSE) is the loss function, (4) the default activation function used is \ReLU, (5)
all computation and training use single precision in PyTorch,
(6) 
all parameters are initialized according to the PyTorch default initialization (Section~\ref{sec:MMNN:learning:stategy}) unless otherwise specified,
(7) $\bmW$'s and $\bmb$'s (the parameters inside the activation functions, see Section~\ref{sec:MMNN:architecture}) are fixed and only $\bmA$'s and $\bmc$'s (the parameters outside the activation functions) are trained,
% $\bmW$'s and $\bmb$'s (the parameters inside the activation functions) are chosen according to Xavier initialization and fixed (see Section~\ref{sec:MMNN:learning:stategy}) and only $\bmA$'s and $\bmc$'s (the parameters outside the activation functions) are trained, 
% 5) computations are implemented on NVIDIA RTX 3500 Ada Generation Laptop GPU.
(8)
computations are conducted on a \textit{NVIDIA RTX 3500 Ada Generation Laptop GPU (power cap 130W)}, with most experiments concluding within a range from a few dozen to several thousand seconds.
% with most experiments concluding within a range from several minutes to tens of minutes.
All our MMNN setups are specified by three parameters $(w,r,l)$ which depend on the function complexity. Another tuning parameter is the learning rate which is guided by the following criteria: (1) not too large initially due to stability, (2) a decreasing rate with iterations such that the learning rate becomes small near the equilibrium to achieve a good accuracy while not decreasing too fast (especially during a long training process for more difficult target functions) so that the training is stalled. 

\subsection{Localized rapid changes}\label{sec:local}
% Deep network of small width can capture local bump; 1) monotone, arctan(100x)
%     2) 2D3D

We begin with two one-dimensional examples. The first is $f(x) = \arctan(100x+20)$, which is smooth but features a rapid transition at zero. While demonstrated in our previous work \cite{ZZZZ-23}, a shallow network struggles to capture such a simple local fast transition which contains high-frequencies, we show that this function can be approximated easily by a composition of a smooth function on top of a (repeated) spatial decomposition and local rescaling using MMNN structure in Section~\ref{sec:MMNN:architecture}.  Our test indeed verifies that our new architecture can effectively capture a localized fast transition rather easily using a very small network of size $(16,4,3)$ as shown in Figure~\ref{fig:arctan:100x}. 
For this test, a total of 1000 data points are uniformly sampled in the range $[-1, 1]$, with a mini-batch size of 100, a learning rate of $10^{-3}$, and the number of epochs set to 2000. Figure~\ref{fig:arctan:100x:error} gives the error plot.
% For this test: a total of 1000 data are uniformly sampled in $[-1,1]$ with a mini-batch size of 100, a learning rate at $10^{-3}$, and a number of epoch at 1000.
% \begin{figure}%[htbp]%[h]%[htbp!]
%     \centering	
%     \foreach \iiidx in {100,200,300,1000}
%     {
%     \begin{subfigure}[c]{0.231\textwidth}
%     \centering            \includegraphics[width=0.998055\textwidth]{figures/Local1D/s1/epoch\iiidx_idx0_w10r3d5.pdf}
%     \subcaption{Epoch \iiidx.}
%     		% \label{subfig:f3}
%     \end{subfigure}
%     }
%     \caption{Illustrations of the training process.}
%     	\label{fig:arctan:100x}
%     % 		\vskip 0.02in
% \end{figure}
% \begin{figure}%[htbp]%[h]
%     \centering
%     \begin{subfigure}[c]{0.2431\textwidth}
%         \centering
%         \includegraphics[width=0.998055\textwidth]{figures/Local1D/s1/epoch60_idx0_w12r3d5.pdf}
%         \subcaption{Epoch 60.}
%     \end{subfigure}
%     \hfill
%     \begin{subfigure}[c]{0.2431\textwidth}
%         \centering
%         \includegraphics[width=0.998055\textwidth]{figures/Local1D/s1/epoch120_idx0_w12r3d5.pdf}
%         \subcaption{Epoch 120.}
%     \end{subfigure}
%     \hfill
%     \begin{subfigure}[c]{0.2431\textwidth}
%         \centering
%         \includegraphics[width=0.998055\textwidth]{figures/Local1D/s1/epoch180_idx0_w12r3d5.pdf}
%         \subcaption{Epoch 180.}
%     \end{subfigure}
%     \hfill
%     \begin{subfigure}[c]{0.2431\textwidth}
%         \centering
%         \includegraphics[width=0.998055\textwidth]{figures/Local1D/s1/epoch2000_idx0_w12r3d5.pdf}
%         \subcaption{Epoch 2000.}
%     \end{subfigure}
%     \caption{Illustrations of the training process.}
%     \label{fig:arctan:100x}
% \end{figure}
% \begin{figure}%[htbp]%[h]
%     \centering
%     \begin{subfigure}[c]{0.2431\textwidth}
%         \centering
%         \includegraphics[width=0.998055\textwidth]{figures/Local1D/s1/epoch100_idx0_w16r4d3.pdf}
%         \subcaption{Epoch 100.}
%     \end{subfigure}
%     \hfill
%     \begin{subfigure}[c]{0.2431\textwidth}
%         \centering
%         \includegraphics[width=0.998055\textwidth]{figures/Local1D/s1/epoch200_idx0_w16r4d3.pdf}
%         \subcaption{Epoch 200.}
%     \end{subfigure}
%     \hfill
%     \begin{subfigure}[c]{0.2431\textwidth}
%         \centering
%         \includegraphics[width=0.998055\textwidth]{figures/Local1D/s1/epoch500_idx0_w16r4d3.pdf}
%         \subcaption{Epoch 500.}
%     \end{subfigure}
%     \hfill
%     \begin{subfigure}[c]{0.2431\textwidth}
%         \centering
%         \includegraphics[width=0.998055\textwidth]{figures/Local1D/s1/epoch2000_idx0_w16r4d3.pdf}
%         \subcaption{Epoch 2000.}
%     \end{subfigure}
%     \caption{Illustrations of the training process.}
%     \label{fig:arctan:100x}
% \end{figure}

\begin{figure}%[htbp]%[h]
\begin{minipage}[b]{0.70\linewidth}
    \centering
    \begin{subfigure}[c]{0.32431\textwidth}
        \centering
        \includegraphics[width=0.998055\textwidth]{figures/Local1D/arctan/epoch100_idx0_w16r4d3.pdf}
        \subcaption{Epoch 100.}
    \end{subfigure}
    \hfill
    \begin{subfigure}[c]{0.32431\textwidth}
        \centering
        \includegraphics[width=0.998055\textwidth]{figures/Local1D/arctan/epoch200_idx0_w16r4d3.pdf}
        \subcaption{Epoch 200.}
    \end{subfigure}
    \hfill
    \begin{subfigure}[c]{0.32431\textwidth}
        \centering
        \includegraphics[width=0.998055\textwidth]{figures/Local1D/arctan/epoch2000_idx0_w16r4d3.pdf}
        \subcaption{Epoch 2000.}
    \end{subfigure}
    \caption{Illustrations of the training process.}
    \label{fig:arctan:100x}
    \end{minipage}
    \hfill
    \begin{minipage}[b]{0.27\linewidth}
   \centering       \includegraphics[width=0.998055\textwidth]{figures/Local1D/arctan/errors1D_w16r4d3arctan.pdf}
    \caption{Training and test errors (MSE).}
    \label{fig:arctan:100x:error} 
    \end{minipage}
\end{figure}


% Next, we consider a more complicated target function, $f(x) = \one_{\{|x+0.2| < 0.02\}} \cdot \sin(50\pi x)$, which represents a localized fast oscillation.
% %\[
% %f(x) = \one_{\{|x| < 0.02\}} \cdot \sin(50\pi x).
% %\]
% %    	\label{fig:sin:50pi:one:period:sample1}
% %    	\label{fig:sin:50pi:one:period:sample1}
%  For this example, we will conduct two tests. The first one is to show the flexibility of MMNN to automatically adapt to local features. The network has a small size as above $(16, 4, 3)$. Each layer has a network width of 16. In other words, each component is a linear combination of 16 \texttt{ReLU} functions which has no way to approximate such a target function well. However, with a multi-layer and multi-component decomposition with parameters appropriately trained by Adam, MMNN can adapt to the behavior of the target function as shown in Figure~\ref{fig:sin:50pi:one:period:sample1}. Figure~\ref{fig:sin:50pi:one:period:error} gives the error plot. Also, the test shows that this example is more difficult to train.  For this test, there are a total of 1000 uniformly sampled points in $[-1,1]$ with a mini-batch size of 100 and a learning rate of $0.002\times 0.95^{\lfloor k/1000 \rfloor}$, where $\lfloor \cdot \rfloor$ denotes floor operation and $k=1,2,\cdots,20000$ is the epoch number. 
% It should be noted that in this test, we initialize the biases $\bmb$'s to $\bm{0}$ and use the PyTorch default initialization method for the weights $\bmW$. This approach, inspired by Xavier initialization, is chosen because the target function is locally oscillatory 
% % near $0$
% and the MMNN size is quite small, necessitating a setup adaptive to the target function to facilitate the training. For other experiments, both the biases and weights use the PyTorch default initialization.
% We then compare with least square approximation using uniform finite element method (FEM) basis with the same degrees of freedom. As shown in Figure~\ref{fig:sin:50pi:FEM:vs:low:rank:NN}, MMNN renders a better approximation due to automatic adaptation through the training process. 
% We would like to remark that when training an extremely compact MMNN which does not have much flexibility and makes the training more subtle, the training hyperparameters, such as learning rate, min-batch size, and etc., need to be more carefully tuned. However, when there is some redundancy in MMNN, i.e., an MMNN with a slightly larger size, MMNN becomes more flexible and the training process becomes easier. On the other hand, when the network becomes too large, then training a large number of parameters and over-redundancy will lead to potential difficulties for optimization. This also shows that there is a trade-off between representation and optimization one needs to balance in practice.

% %The second test is on non-uniform sampling. 
% %or this function, we will conduct two comparisons: 1) We will compare the finite element method and our architecture, which has ``insufficient" degrees of freedom (or parameters), using uniform grid samples; 2) We will evaluate the performance of training our network with uniform grid samples against training under conditions of localized concentration.
% % A hybrid approach where 10\% of the samples are concentrated within $[-0.02, 0.02]$, consisting of uniform grid sampling with 900 points across the domain $[-1,1]$ and 100 uniform grid samples within $[-0.02, 0.02]$; 3) Modifying the second method by increasing the central concentration to 70\% of the samples.


% % \begin{figure}%[htbp]%[h]%[htbp!]
% %     \centering	
% %     \foreach \iiidx in {3000,3500,4000,10000}
% %     {
% %     \begin{subfigure}[c]{0.231\textwidth}
% %     \centering            \includegraphics[width=0.998055\textwidth]{figures/Local1D/s1/epoch\iiidx_idx1_w10r3d5.pdf}
% %     \subcaption{Epoch \iiidx.}
% %     		% \label{subfig:f3}
% %     \end{subfigure}
% %     }
% % \caption{Illustrations of the training process.}
% %     	\label{fig:sin:50pi:one:period:sample1}
% %     % 		\vskip 0.02in
% % \end{figure}
% \begin{figure}%[htbp]%[h]
%     \centering
%     \begin{subfigure}[c]{0.231\textwidth}
%         \centering
%         \includegraphics[width=0.998055\textwidth]{figures/Local1D/sin1P/epoch500_idx0_w16r4d3.pdf}
%         \subcaption{Epoch 500.}
%     \end{subfigure}
%     \hfill
%     \begin{subfigure}[c]{0.231\textwidth}
%         \centering
%         \includegraphics[width=0.998055\textwidth]{figures/Local1D/sin1P/epoch1000_idx0_w16r4d3.pdf}
%         \subcaption{Epoch 1000.}
%     \end{subfigure}
%     \hfill
%     \begin{subfigure}[c]{0.231\textwidth}
%         \centering
%         \includegraphics[width=0.998055\textwidth]{figures/Local1D/sin1P/epoch2000_idx0_w16r4d3.pdf}
%         \subcaption{Epoch 2000.}
%     \end{subfigure}
%     \hfill
%     \begin{subfigure}[c]{0.231\textwidth}
%         \centering
%         \includegraphics[width=0.998055\textwidth]{figures/Local1D/sin1P/epoch20000_idx0_w16r4d3.pdf}
%         \subcaption{Epoch 20000.}
%     \end{subfigure}
%     \caption{Illustrations of the training process.}
%     \label{fig:sin:50pi:one:period:sample1}
% \end{figure}


% \begin{figure}%[htbp]%[h]%[htbp!]
%     \begin{minipage}[b]{0.32\linewidth}
%    \centering       \includegraphics[width=0.886\textwidth]{figures/Local1D/sin1P/errors1D_w16r4d3sin1P.pdf}
%     \caption{Training and test errors measured in MSE vs. epoch.}    \label{fig:sin:50pi:one:period:error} 
%     \end{minipage}
%     \hfill
% \begin{minipage}[b]{0.64\linewidth}
%     \centering	
%             \begin{subfigure}[c]{0.458\textwidth}
%     \centering            \includegraphics[width=0.8998055\textwidth]{figures/Local1D/sin1P/fem153.pdf}
%     \subcaption{FEM.}
%     		% \label{subfig:f3}
%     \end{subfigure}
%     \hspace{5pt}
%     \begin{subfigure}[c]{0.458\textwidth}
%     \centering            \includegraphics[width=0.8998055\textwidth]{figures/Local1D/sin1P/epoch20000_idx0_w16r4d3.pdf}
%     \subcaption{MMNN.}
%     		% \label{subfig:f3}
%     \end{subfigure}
% \caption{(a) Least square using equally spaced 153 FEM bases. (b) MMNN with $(16+1) \times 4 \times (3-1) + (16+1) = 153$ free parameters.} 	\label{fig:sin:50pi:FEM:vs:low:rank:NN}
%     % 		\vskip 0.02in
%     \end{minipage}
% \end{figure}

%As demonstrated in Figure~\ref{fig:sin:50pi:one:period:sample1}, our compact yet deep network effectively learns local high frequencies using uniform grid sampling. Compared to the Finite Element Method (FEM) with the same degrees of freedom, our network exhibits superior performance, as illustrated in Figure~\ref{fig:sin:50pi:FEM:vs:low:rank:NN}. This highlights the \textbf{adaptability} of our network architecture, even with limited degrees of freedom.


%As mentioned in our prior publication \cite{ZZZZ-23}, functions featuring a localized bump (peak) present a significant challenge for shallow networks. Nonetheless, the figures above demonstrate that these types of functions are readily handled by our newly proposed architecture, which employs balanced decompositions. This, in a way, showcases the effectiveness of compositional strategies in deep neural networks.

Next, we consider a more complicated target function, $f(x) = \one_{\{|x+0.2| < 0.02\}} \cdot \sin(50\pi x)$, which represents a localized fast oscillation.
%\[
%f(x) = \one_{\{|x| < 0.02\}} \cdot \sin(50\pi x).
%\]
%    	\label{fig:sin:50pi:one:period:sample1}
%    	\label{fig:sin:50pi:one:period:sample1}
 For this example, we will conduct two tests. The first one is to show the flexibility of MMNN to automatically adapt to local features. The network has a small size as above $(16, 4, 3)$. Each layer has a network width of 16. In other words, each component is a linear combination of 16 \texttt{ReLU} functions which has no way to approximate such a target function well. However, with a multi-layer and multi-component decomposition with parameters appropriately trained by Adam, MMNN can adapt to the behavior of the target function as shown in Figure~\ref{fig:sin:50pi:one:period:sample1}. Figure~\ref{fig:sin:50pi:one:period:error} gives the error plot. Also, the test shows that this example is more difficult to train.  For this test, there are a total of 1000 uniformly sampled points in $[-1,1]$ with a mini-batch size of 100 and a learning rate of $0.002\times 0.95^{\lfloor k/1000 \rfloor}$, where $\lfloor \cdot \rfloor$ denotes floor operation and $k=1,2,\cdots,20000$ is the epoch number. 
It should be noted that in this test, we initialize the biases $\bmb$'s to $\bm{0}$ and use the PyTorch default initialization method for the weights $\bmW$. This approach, inspired by Xavier initialization, is chosen because the target function is locally oscillatory 
% near $0$
and the MMNN size is quite small, necessitating a setup adaptive to the target function to facilitate the training. For other experiments, both the biases and weights use the PyTorch default initialization.
We then compare with least square approximation using uniform finite element method (FEM) basis with the same degrees of freedom. As shown in Figure~\ref{fig:sin:50pi:FEM:vs:low:rank:NN}, MMNN renders a better approximation due to automatic adaptation through the training process. 
We remark that when training an MMNN with an extreme compact size with respect to the target function complexity, due to the lack of flexibility/redundancy, the training may become more subtle and need more careful calibration, such as initialization, learning rate, min-batch size, and etc. 
% However, when there is some redundancy in an MMNN, i.e., an MMNN with a slightly larger size, MMNN becomes more flexible and the training process becomes easier.
However, introducing slight redundancy into an MMNN, such as by increasing its size marginally, enhances its flexibility and makes training more tractable.
On the other hand, when the network becomes too large, then training a large number of parameters and over-redundancy will lead to potential difficulties for optimization. This also shows that there is a trade-off between representation and optimization one needs to balance in practice.

%The second test is on non-uniform sampling. 
%or this function, we will conduct two comparisons: 1) We will compare the finite element method and our architecture, which has ``insufficient" degrees of freedom (or parameters), using uniform grid samples; 2) We will evaluate the performance of training our network with uniform grid samples against training under conditions of localized concentration.
% A hybrid approach where 10\% of the samples are concentrated within $[-0.02, 0.02]$, consisting of uniform grid sampling with 900 points across the domain $[-1,1]$ and 100 uniform grid samples within $[-0.02, 0.02]$; 3) Modifying the second method by increasing the central concentration to 70\% of the samples.


% \begin{figure}%[htbp]%[h]%[htbp!]
%     \centering	
%     \foreach \iiidx in {3000,3500,4000,10000}
%     {
%     \begin{subfigure}[c]{0.231\textwidth}
%     \centering            \includegraphics[width=0.998055\textwidth]{figures/Local1D/s1/epoch\iiidx_idx1_w10r3d5.pdf}
%     \subcaption{Epoch \iiidx.}
%     		% \label{subfig:f3}
%     \end{subfigure}
%     }
% \caption{Illustrations of the training process.}
%     	\label{fig:sin:50pi:one:period:sample1}
%     % 		\vskip 0.02in
% \end{figure}
\begin{figure}%[htbp]%[h]
    \centering
    \begin{subfigure}[c]{0.231\textwidth}
        \centering
        \includegraphics[width=0.998055\textwidth]{figures/Local1D/sin1P/epoch500_idx0_w16r4d3.pdf}
        \subcaption{Epoch 500.}
    \end{subfigure}
    \hfill
    \begin{subfigure}[c]{0.231\textwidth}
        \centering
        \includegraphics[width=0.998055\textwidth]{figures/Local1D/sin1P/epoch1000_idx0_w16r4d3.pdf}
        \subcaption{Epoch 1000.}
    \end{subfigure}
    \hfill
    \begin{subfigure}[c]{0.231\textwidth}
        \centering
        \includegraphics[width=0.998055\textwidth]{figures/Local1D/sin1P/epoch2000_idx0_w16r4d3.pdf}
        \subcaption{Epoch 2000.}
    \end{subfigure}
    \hfill
    \begin{subfigure}[c]{0.231\textwidth}
        \centering
        \includegraphics[width=0.998055\textwidth]{figures/Local1D/sin1P/epoch20000_idx0_w16r4d3.pdf}
        \subcaption{Epoch 20000.}
    \end{subfigure}
    \caption{Illustrations of the training process.}
    \label{fig:sin:50pi:one:period:sample1}
\end{figure}


% \begin{figure}%[htbp]%[h]%[htbp!]
%     \begin{minipage}[b]{0.32\linewidth}
%    \centering       \includegraphics[width=0.886\textwidth]{figures/Local1D/sin1P/errors1D_w16r4d3sin1P.pdf}
%     \caption{Training and test errors measured in MSE vs. epoch.}    \label{fig:sin:50pi:one:period:error} 
%     \end{minipage}
%     \hfill
% \begin{minipage}[b]{0.64\linewidth}
%     \centering	
%             \begin{subfigure}[c]{0.458\textwidth}
%     \centering            \includegraphics[width=0.8998055\textwidth]{figures/Local1D/sin1P/fem153.pdf}
%     \subcaption{FEM.}
%     		% \label{subfig:f3}
%     \end{subfigure}
%     \hspace{5pt}
%     \begin{subfigure}[c]{0.458\textwidth}
%     \centering            \includegraphics[width=0.8998055\textwidth]{figures/Local1D/sin1P/epoch20000_idx0_w16r4d3.pdf}
%     \subcaption{MMNN.}
%     		% \label{subfig:f3}
%     \end{subfigure}
% \caption{(a) Least square using equally spaced 153 FEM bases. (b) MMNN with $(16+1) \times 4 \times (3-1) + (16+1) = 153$ free parameters.} 	\label{fig:sin:50pi:FEM:vs:low:rank:NN}
%     % 		\vskip 0.02in
%     \end{minipage}
% \end{figure}

\begin{figure}%[htbp]%[h]%[htbp!]
    \begin{minipage}[b]{0.32\linewidth}
   \centering       \includegraphics[width=0.75\textwidth]{figures/Local1D/sin1P/errors1D_w16r4d3sin1P.pdf}
    \caption{Training and test errors (in MSE) vs. epoch.}    \label{fig:sin:50pi:one:period:error} 
    \end{minipage}
    \hfill
\begin{minipage}[b]{0.64\linewidth}
    \centering	
            \begin{subfigure}[c]{0.458\textwidth}
    \centering            \includegraphics[width=0.8998055\textwidth]{figures/Local1D/sin1P/fem153.pdf}
    % \subcaption{FEM.}
    		% \label{subfig:f3}
    \end{subfigure}
    \hspace{5pt}
    \begin{subfigure}[c]{0.458\textwidth}
    \centering            \includegraphics[width=0.8998055\textwidth]{figures/Local1D/sin1P/epoch20000_idx0_w16r4d3.pdf}
    % \subcaption{MMNN.}
    		% \label{subfig:f3}
    \end{subfigure}
% \caption{Left: Least square using equally spaced 153 FEM bases. Right: MMNN with $(16+1) \times 4 \times (3-1) + (16+1) = 153$ free parameters.} 
\caption{Left: Least square using equally spaced 153 FEM bases. Right: MMNN with $153$ free parameters.} 	\label{fig:sin:50pi:FEM:vs:low:rank:NN}
    % 		\vskip 0.02in
    \end{minipage}
\end{figure}


Finally, we show an example in the two-dimensional  case shown in~Figure~\ref{fig:f3D3:local:HF} and defined in polar coordinates by
\begin{equation*}
    f(r,\theta)=\begin{cases}
        0 & \tn{if  } 0.5+25\rho-25r\le 0,\\
        1 & \tn{if  } 0.5+25\rho-25r\ge 1,\\
        0.5+25\rho-25r & \tn{otherwise},\\
    \end{cases}
    \quad \tn{where}\ \rho=0.1+0.02\cos(8\pi  \theta).
\end{equation*}
% \red{1. Wavelet and FEM, to compare NN\\
% 2. width decrease, 100, 80, 60\\
% 3. target function with two fast transition, arctan(100x)-arctan(100(x-0.3))}
Again a rather compact MMNN of size $(100,10,6)$ can produce a good approximation. 
% Figure~\ref{fig:2D-2} shows the error during the training process and 
Figure~\ref{fig:2D-3} shows the $\log$ plot of training and testing errors in MSE. For this test there are a total of $400^2$ uniformly sampled points in $[-1,1]^2$ with mini-batch size of $1000$ and a learning rate of $10^{-3}\times 0.9^{\lfloor k/25 \rfloor}$, where $k=1,2,\cdots,1000$ is the epoch number. We compare the result with piecewise linear interpolation and least square approximation using FEM basis on a uniform grid with the same number of degrees of freedom in Figure~\ref{fig:comparison}. As observed before, MMNN renders the best result due to its adaptivity through training.
% \begin{Comment}
    When adaptive finite element is applicable, it is hard to beat. However, adaptive FEMs are more humanly involved, while MMNNs based on training are more automatic. The key message here is to demonstrate that MMNN has adaptive features through training/optimization, which will be useful when an adaptive finite element method becomes difficult or impossible in applications. 
% \end{Comment}


\begin{figure}%[htbp]%[h]%[htbp!]
\centering
\begin{minipage}[b]{0.5427161\linewidth}
   \centering	
    \begin{subfigure}[c]{0.48538\textwidth}
    \centering            \includegraphics[height=0.75\textwidth]{figures/Local2D/LHF2D3.png}
    % \subcaption{}
    %		\label{subfig:f3}
    \end{subfigure}
    \hfill
    \begin{subfigure}[c]{0.48538\textwidth}
    \centering            \includegraphics[height=0.75\textwidth]{figures/Local2D/LHFf3D3.pdf}
    % \subcaption{}
    %		\label{subfig:f3}
    \end{subfigure}
    \caption{Target function.}
    	\label{fig:f3D3:local:HF}
    % 		\vskip 0.02in
\end{minipage}
% \hfill
\hspace{0.1\linewidth}
\begin{minipage}[b]{0.2963\linewidth}
    \centering	    \includegraphics[width=0.95\textwidth]{figures/Local2D/errorsLHF2D3_w100r10d6.pdf}
    % \subcaption{}
    %		\label{subfig:f3}
    \caption{Errors (in MSE).}
    	\label{fig:2D-3}
    % 		\vskip 0.02in
\end{minipage}
\end{figure}


%\begin{figure}%[htbp]%[h]%[htbp!]
%    \centering	
%    \foreach \iiidx in {10,50,100,1000}
%    {
%    \begin{subfigure}[c]{0.23\textwidth}
%    \centering            \includegraphics[width=0.998055\textwidth]{figures/Local2D/epoch\iiidx_idx1_w100r10d6_diff.png}
%    \subcaption{Epoch \iiidx.}
%    %		\label{subfig:f3}
%    \end{subfigure}
%    }
%    \caption{Training process}
%    	\label{fig:2D-1}
%    % 		\vskip 0.02in
%\end{figure}

% \begin{figure}%[htbp]%[h]%[htbp!]
%     \centering	
%     \foreach \iiidx in {10,50,100,1000}
%     {
%     \begin{subfigure}[c]{0.23\textwidth}
%     \centering            \includegraphics[width=0.998055\textwidth]{figures/Local2D/epoch\iiidx_idx1_w100r10d6_f.png}
%     \subcaption{Epoch \iiidx.}
%     %\label{subfig:2D-2}
%     \end{subfigure}
%     }
%     \caption{The error during the training process.}
%     	\label{fig:2D-2}
%     % 		\vskip 0.02in
% \end{figure}

% \begin{figure}%[htbp]%[h]
%     \centering
%     \begin{subfigure}[c]{0.23\textwidth}
%         \centering
%         \includegraphics[width=0.998055\textwidth]{figures/Local2D/epoch10_idx1_w100r10d6_f.png}
%         \subcaption{Epoch 10.}
%     \end{subfigure}
%     \hfill
%     \begin{subfigure}[c]{0.23\textwidth}
%         \centering
%         \includegraphics[width=0.998055\textwidth]{figures/Local2D/epoch50_idx1_w100r10d6_f.png}
%         \subcaption{Epoch 50.}
%     \end{subfigure}
%     \hfill
%     \begin{subfigure}[c]{0.23\textwidth}
%         \centering
%         \includegraphics[width=0.998055\textwidth]{figures/Local2D/epoch100_idx1_w100r10d6_f.png}
%         \subcaption{Epoch 100.}
%     \end{subfigure}
%     \hfill
%     \begin{subfigure}[c]{0.23\textwidth}
%         \centering
%         \includegraphics[width=0.998055\textwidth]{figures/Local2D/epoch1000_idx1_w100r10d6_f.png}
%         \subcaption{Epoch 1000.}
%     \end{subfigure}
%     \caption{The error during the training process.}
%     \label{fig:2D-2}
% \end{figure}


% Target function 
% \[
% f(x)=\one_{\{|x|<0.05\}} \cdot\sin(
%         50\pi  x)
% \]
% \begin{figure}%[htbp]%[h]%[htbp!]
%     \centering	
%     \foreach \iiidx in {1,6,...,36}
%     {
%     \begin{subfigure}[c]{0.31\textwidth}
%     \centering            \includegraphics[width=0.998055\textwidth]{figures/LB2D/1/epoch\iiidx_idx1.png}
%     \subcaption{Epoch \iiidx.}
%     %		\label{subfig:f3}
%     \end{subfigure}
%     }
%     \caption{Illustrations of learned NN and the target function.}
%     	% \label{fig:}
%     % 		\vskip 0.02in
% \end{figure}



\begin{figure}%[htbp]%[h]%[htbp!]
    \centering	
            \begin{subfigure}[c]{0.312\textwidth}
    \centering            \includegraphics[width=0.75859598055\textwidth]{figures/Local2D/epoch1000_idx1_w100r10d6_diff.png}
    \subcaption{Network.}
    		% \label{subfig:f3}
    \end{subfigure}
    \hfill
    \begin{subfigure}[c]{0.312\textwidth}
    \centering            \includegraphics[width=0.75859598055\textwidth]{figures/Local2D/LHF2D4interp.png}
    \subcaption{Interpolation.}
    		% \label{subfig:f3}
    \end{subfigure}
    \hfill 
 \begin{subfigure}[c]{0.312\textwidth}
    \centering            \includegraphics[width=0.75859598055\textwidth]{figures/Local2D/LHF2D5fem.png}
    \subcaption{FEM.}
    		% \label{subfig:f3}
    \end{subfigure}\hfill
    \begin{subfigure}[c]{0.32\textwidth}
    \centering            \includegraphics[width=0.75859598055\textwidth]{figures/Local2D/epoch1000_idx1_w100r10d6_f.png}
    \subcaption{Network (error).}
    		% \label{subfig:f3}
    \end{subfigure}                
    \hfill
    \begin{subfigure}[c]{0.32\textwidth}
    \centering            \includegraphics[width=0.75859598055\textwidth]{figures/Local2D/LHF2D4interpDiff.png}
    \subcaption{Interpolation (error).}
    
    		% \label{subfig:f3}
    \end{subfigure}     
    \hfill
    \begin{subfigure}[c]{0.32\textwidth}
    \centering            \includegraphics[width=0.75859598055\textwidth]{figures/Local2D/LHF2D5femDiff.png}
    \subcaption{FEM (error).}    
    		% \label{subfig:f3}
    \end{subfigure}
    \caption{Comparison among different approximations using MMNN, interpolation, and least square FEM. The interpolation and FEM are all based on a $72\times 72=5184$ uniform grid. MMNN has $(100+1) \times 10 \times (6-1) + (100+1) = 5151$ free parameters. The maximum error is approximately $0.05$ for MMNN, $0.31$ for interpolation, and $0.38$ for FEM. The corresponding MSE errors are $0.85 \times 10^{-6}$, $1.95 \times 10^{-4}$, and $1.45 \times 10^{-4}$, respectively.}
% \caption{Comparison among our network, interpolation and FEM method. Interpolation employs $72\times 72=5184$ uniform grid points and FEM method use that many bases. Our network architecture is designed with $(100+1) \times 10 \times (6-1) + (100+1) = 5151$ degrees of freedom (parameters). The maximum error for neural networks is approximately $0.05$, for interpolation it is around $0.31$, for FEM is  $0.38$. For MSE, $0.85\times 10^{-6}$ vs 1.95e-04 vs 1.45e-04}    	\label{fig:interpolation:vs:low:rank:NN}
    % 		\vskip 0.02in
\label{fig:comparison}
\end{figure}

% \clearpage
% \red{
% Note that all settings for FMMNN remain identical to those of ReLU MMNN, except that the activation function is replaced from ReLU to sine.

% FMMNN vs (ReLU MMNN, Interpolation, and least square FEM):

% \begin{figure}%[htbp]%[h]%[ht]%[htbp!]
%     \centering	
%     \includegraphics[width=0.5595\textwidth]{figures/FMMNNerrorsLHF2D3_w100r10d6.pdf}
%     % \subcaption{}
%     %		\label{subfig:f3}
%     \caption{FMMNN: Errors (in MSE) vs. epoch.}
%     	% \label{fig:errors:LD:fun2D}
%     % 		\vskip 0.02in
% \end{figure}

% \begin{figure}%[htbp]%[h]%[ht]%[htbp!]
%     \centering	
%     \includegraphics[width=0.5595\textwidth]{figures/FMMNNvsReLUMMNNerrorsLHF2D3_w100r10d6.pdf}
%     % \subcaption{}
%     %		\label{subfig:f3}
%     \caption{FMMNN vs ReLU: test errors (in MSE) vs. epoch.}
%     	% \label{fig:errors:LD:fun2D}
%     % 		\vskip 0.02in
% \end{figure}

% \begin{figure}%[htbp]%[h]%[ht]%[htbp!]
%     \centering	
%             \begin{subfigure}[c]{0.430312\textwidth}
%     \centering            \includegraphics[width=0.75859598055\textwidth]{figures/FMMNNepoch1000_idx1_w100r10d6_diff.png}
%     \subcaption{FMMNN.}
%     		% \label{subfig:f3}
%     \end{subfigure}
%     \hfill
%     \begin{subfigure}[c]{0.430312\textwidth}
%     \centering            \includegraphics[width=0.75859598055\textwidth]{figures/FMMNNepoch1000_idx1_w100r10d6_f.png}
%     \subcaption{FMMNN (error).}
%     		% \label{subfig:f3}
%     \end{subfigure}
%     \caption{Max error $0.036$; MSE: $1.53 \times 10^{-6}$.}
% % \caption{Comparison among our network, interpolation and FEM method. Interpolation employs $72\times 72=5184$ uniform grid points and FEM method use that many bases. Our network architecture is designed with $(100+1) \times 10 \times (6-1) + (100+1) = 5151$ degrees of freedom (parameters). The maximum error for neural networks is approximately $0.05$, for interpolation it is around $0.31$, for FEM is  $0.38$. For MSE, $0.85\times 10^{-6}$ vs 1.95e-04 vs 1.45e-04}    	\label{fig:interpolation:vs:low:rank:NN}
%     % 		\vskip 0.02in
% % \label{fig:comparison}
% \end{figure}
% In this local oscillatory example, the performance of sine MMNN and ReLU MMNN is quite similar.

% }




\subsection{Highly oscillatory functions}\label{sec:oscillatory}

Globally oscillatory functions with significant high-frequency components can not be approximated well by a shallow network when a global bounded activation function of the form $\sigma(\bmW\cdot \bmx-\bmb)$, such as \texttt{ReLU}, is used. Due to almost orthogonality or high decorrelation (in terms of the inner product) between $\sigma(\bmW\cdot \bmx-\bmb)$ and oscillatory functions with high likelihood (in terms of a random choice of $(\bmW, \bmb)$), the set of parameters that can render a good approximation, namely the \emph{Rashomon set}~\cite{semenova2022existence}, becomes smaller and smaller (in terms of relative measure) and hence harder and harder to find as the target function becomes more and more oscillatory (see~\cite{ZZZZ-23}). Although this difficulty can be alleviated by complexity decomposition using MMNN as shown in Section~\ref{sec:MMNN}, it still requires a larger network in terms of width, rank, and layers and more training. Here we limit our tests to oscillatory functions in one-dimensional and two-dimensional cases due to the dramatic increase of complexity with dimensions, or the curse of dimensions, in general. 

We again start with a one-dimensional example, $f(x)=\sin(50\pi x), x\in[-1,1]$. An MMNN of size $(800,40,15)$ produces a good approximation of this highly oscillatory function, as illustrated by the error plot in Figure~\ref{fig:50pi-error}, with a smaller learning rate and a longer training process compared to previous examples with localized fine features. Due to the significant depth, we consider using ResMMNN as discussed in Section~\ref{sec:MMNN:architecture}. For this test, a total of $1000$ uniformly sampled points in $[-1,1]$ are used with a mini-batch size of $100$ and a learning rate of $10^{-4}\times 0.9^{\lfloor k/800 \rfloor}$, where $k=1,2,\cdots,40000$ is the epoch number. Also, an interesting learning dynamics for Adam is observed from Figure~\ref{fig:sin:50pi:NN}. In the beginning, nothing seems to happen until about epoch 3600 when learning starts from the boundary. Then more and more features are captured from the boundary to the inside gradually. Eventually, all features are captured and then fine-tuned together to improve the overall approximation. 

%Let's now focus on functions exhibiting global high frequency. We will limit our analysis to one-dimensional and 2D cases because high-dimensional scenarios require substantial computational power if the target function displays global high frequency in such spaces, which is affected by the curse of dimensionality.


% First, 
%We begin with the one-dimensional case, choosing the function $\sin(50\pi x)$ on the interval $[-1,1]$ as our target. The diagram in Figure~\ref{fig:sin:50pi:NN} clearly shows that learning starts at the boundary, consistent with our earlier theoretical analysis.

% \begin{figure}%[htbp]%[h]%[htbp!]
% \begin{minipage}[b]{0.65\linewidth}
%     \centering	
%     % \foreach \iiidx in {3500,4000,4500,11000,14000,40000}
%     % {\begin{subfigure}[c]{0.31\textwidth}
%     % \centering            \includegraphics[width=0.99\textwidth]{figures/Global1D/epoch\iiidx_idx0_w800r40d15.png}
%     % \subcaption{Epoch \iiidx.}
%     % %		\label{subfig:f3}
%     % \end{subfigure}
%     % }
%     \begin{subfigure}[c]{0.325\textwidth}
%         \centering       \includegraphics[width=0.99\textwidth]{figures/Global1D/epoch4000_idx0_w800r40d15.pdf}
%         \subcaption{Epoch 4000.}
%     \end{subfigure}
%     \hfill
%     \begin{subfigure}[c]{0.325\textwidth}
%         \centering
%         \includegraphics[width=0.99\textwidth]{figures/Global1D/epoch5000_idx0_w800r40d15.pdf}
%         \subcaption{Epoch 5000.}
%     \end{subfigure}
%     \hfill
%     \begin{subfigure}[c]{0.325\textwidth}
%         \centering
%         \includegraphics[width=0.99\textwidth]{figures/Global1D/epoch6000_idx0_w800r40d15.pdf}
%         \subcaption{Epoch 6000.}
%     \end{subfigure}
% \\  
%     \begin{subfigure}[c]{0.325\textwidth}
%         \centering
%         \includegraphics[width=0.99\textwidth]{figures/Global1D/epoch7000_idx0_w800r40d15.pdf}
%         \subcaption{Epoch 7000.}
%     \end{subfigure}
%     \hfill
%     \begin{subfigure}[c]{0.325\textwidth}
%         \centering
%         \includegraphics[width=0.99\textwidth]{figures/Global1D/epoch14000_idx0_w800r40d15.pdf}
%         \subcaption{Epoch 14000.}
%     \end{subfigure}
%     \hfill
%     \begin{subfigure}[c]{0.325\textwidth}
%         \centering
%         \includegraphics[width=0.99\textwidth]{figures/Global1D/epoch40000_idx0_w800r40d15.pdf}
%         \subcaption{Epoch 40000.}
%     \end{subfigure}
%     \caption{Illustrations of the training process.}
%     	\label{fig:sin:50pi:NN}
%     % 		\vskip 0.02in
%     \end{minipage}
%     \hfill
%     \begin{minipage}[b]{0.33\linewidth}
%     \centering	
%     \includegraphics[width=0.999\textwidth]{figures/Global1D/errors1D_w800r40d15.pdf}
%     % \subcaption{}
%     %		\label{subfig:f3}
%     \caption{Illustrations of training and test errors measured in MSE.}
%     	\label{fig:50pi-error}
%     % 		\vskip 0.02in
%     \end{minipage}
% \end{figure}

\begin{figure}%[htbp]%[h]%[htbp!]
\begin{minipage}[b]{0.59\linewidth}
    \centering	
    \begin{subfigure}[c]{0.32\textwidth}
        \centering       \includegraphics[width=0.99\textwidth]{figures/Global1D/epoch3600_idx0_w800r40d15.pdf}
        \subcaption{Epoch 3600.}
    \end{subfigure}
    \hfill
    \begin{subfigure}[c]{0.32\textwidth}
        \centering
        \includegraphics[width=0.99\textwidth]{figures/Global1D/epoch3800_idx0_w800r40d15.pdf}
        \subcaption{Epoch 3800.}
    \end{subfigure}
    \hfill
    \begin{subfigure}[c]{0.32\textwidth}
        \centering
        \includegraphics[width=0.99\textwidth]{figures/Global1D/epoch4200_idx0_w800r40d15.pdf}
        \subcaption{Epoch 4200.}
    \end{subfigure}
\\  
    \begin{subfigure}[c]{0.32\textwidth}
        \centering
        \includegraphics[width=0.99\textwidth]{figures/Global1D/epoch5000_idx0_w800r40d15.pdf}
        \subcaption{Epoch 5000.}
    \end{subfigure}
    \hfill
    \begin{subfigure}[c]{0.32\textwidth}
        \centering
        \includegraphics[width=0.99\textwidth]{figures/Global1D/epoch10000_idx0_w800r40d15.pdf}
        \subcaption{Epoch 10000.}
    \end{subfigure}
    \hfill
    \begin{subfigure}[c]{0.32\textwidth}
        \centering
        \includegraphics[width=0.99\textwidth]{figures/Global1D/epoch40000_idx0_w800r40d15.pdf}
        \subcaption{Epoch 40000.}
    \end{subfigure}
    \caption{Illustrations of the training process.}
    	\label{fig:sin:50pi:NN}
    % 		\vskip 0.02in
    \end{minipage}
    \hfill
    \begin{minipage}[b]{0.36\linewidth}
    \centering	
    \includegraphics[width=0.999\textwidth]{figures/Global1D/errors1D_w800r40d15.pdf}
    % \subcaption{}
    %		\label{subfig:f3}
    \caption{Training and test errors (in MSE) vs. epoch.}
    	\label{fig:50pi-error}
    % 		\vskip 0.02in
    \end{minipage}
\end{figure}

% \begin{figure}%[htbp]%[h]%[htbp!]
%     \centering	
%     \foreach \iiidx in {3500,4000,4500,5000,8000,11000,14000,40000}
%     {
%     \begin{subfigure}[c]{0.231\textwidth}
%     \centering            \includegraphics[width=0.9598055\textwidth]{figures/Global1D/epoch\iiidx_idx0_w800r40d15.png}
%     \subcaption{Epoch \iiidx.}
%     %		\label{subfig:f3}
%     \end{subfigure}
%     }
%     \caption{Illustrations of the training process.}
%     	\label{fig:sin:50pi:NN}
%     % 		\vskip 0.02in
% \end{figure}

% \begin{figure}%[htbp]%[h]%[htbp!]
%     \centering	
%     \includegraphics[width=0.5\textwidth]{figures/Global1D/errors1D_w800r40d15.pdf}
%     % \subcaption{}
%     %		\label{subfig:f3}
%     \caption{Illustrations of training and test errors measured in MSE.}
%     	\label{fig:50pi-error}
%     % 		\vskip 0.02in
% \end{figure}

Next, we consider a two-dimensional target function of the following form:
% Let $\phi_\bmtheta$ be the network-generated function parameterized with $\bmtheta$. Given a smooth target function $f$ on $[0,1]^d$, our goal is to identify 
% \begin{equation*}
% 	\bmtheta^* =\argmin_{\bmtheta} 
% 	%	\Big(\Big)
% 	\|\phi_\bmtheta-f\|_{L^2}^2
% 	%+\lambda_1\|\bmtheta\|_{L^2}^2
% \end{equation*}
% such that the learned function $\phi_{\bmtheta^*}$ approximate the target function $f$ well in the $L^\infty$-norm.
%\begin{equation*}
%	f_1(x_1,x_2)= \sum_{i=1}^2\sum_{j=1}^2 a_{ij} \sin(b_i x_i+c_{ij} x_i x_j) \cos (b_j x_j + d_{ij} x_i^2)
%\end{equation*}
%and
\begin{equation*}
	f_{s}(x_1,x_2)= \sum_{i=1}^2\sum_{j=1}^2 a_{ij} \sin(s b_i x_i+s c_{ij} x_i x_j) \cos (s b_j x_j + s d_{ij} x_i^2),
\end{equation*}
 where
\begin{equation*}
	(a_{i,j})=	\begin{bmatrix*}
		0.3 & 0.2 \\ 0.2 & 0.3
	\end{bmatrix*}, \qquad 
	(b_{i})=	\begin{bmatrix*}
		2\pi \\ 4\pi
	\end{bmatrix*},
	\qquad
	(c_{i,j})=	\begin{bmatrix*}
		2\pi & 4\pi \\ 8\pi & 4\pi
	\end{bmatrix*},  \quad 
	(d_{i,j})=	\begin{bmatrix*}
		4\pi & 6\pi \\ 8\pi & 6\pi
	\end{bmatrix*}.
\end{equation*}
In our test, we choose $s=3$ to ensure the function exhibits significant oscillations and contains diverse Fourier modes as illustrated by Figure~\ref{fig:2D-oscillation}. Given the complexity of the function, we employ an MMNN with size $(600, 30, 15)$. Again, ResMMNN is used due to the depth. For this test, a total of $400^2$ data are sampled on a uniform grid in $[-1,1]^2$ with a mini-batch size of $1000$ and a learning rate of $10^{-3}\times 0.9^{\lfloor k/40 \rfloor}$, where $k=1,2,\cdots,2000$ is the epoch number. The training process is illustrated by Figure~\ref{fig:2D-oscillation-difference:and:NN}. Figure~\ref{fig:2D-oscillation-error} shows $\log$-error plot.


\begin{figure}%[htbp]%[h]%[htbp!]
\begin{minipage}[b]{0.63597166\linewidth}
    \centering	
    \begin{subfigure}[c]{0.48736\textwidth}
    \centering            \includegraphics[height=0.7850\textwidth]{figures/Global2D/HForiginal3.png}
    % \subcaption{}
    %		\label{subfig:f3}
    \end{subfigure}
    \hfill
    \begin{subfigure}[c]{0.48736\textwidth}
    \centering            \includegraphics[height=0.7850\textwidth]{figures/Global2D/f3D.pdf}
    % \subcaption{}
    %		\label{subfig:f3}
    \end{subfigure}
    \caption{Illustrations of the target function.}
    	\label{fig:2D-oscillation}
    % 		\vskip 0.02in
\end{minipage}
\hfill
\begin{minipage}[b]{0.318\linewidth}
    \centering	
    \includegraphics[width=0.9985\textwidth]{figures/Global2D/errors2D_w600r30d15.pdf}
    % \subcaption{}
    %		\label{subfig:f3}
    \caption{Training and test errors  (in MSE) vs. epoch.}
    	\label{fig:2D-oscillation-error}
    % 		\vskip 0.02in
\end{minipage}
\end{figure}


% \begin{figure}%[htbp]%[h]%[htbp!]
%     \centering	
%     \foreach \iiidx in {10,100,500,1000}
%     {
%     \begin{subfigure}[c]{0.234\textwidth}
%     \centering            \includegraphics[width=0.998055\textwidth]{figures/Global2D/epoch\iiidx_idx1_w600r30d15_diff.png}
%     % \subcaption{Epoch \iiidx.}
%     %\label{subfig:2D}
%     \end{subfigure}
%     }
% %     \caption{}
% %     	\label{fig:2D-oscillation-training}
% %     % 		\vskip 0.02in
% % \end{figure}

% % \begin{figure}%[htbp]%[h]%[htbp!]
% %     \centering	
%     \foreach \iiidx in {25,50,1000,2000}
%     {
%     \begin{subfigure}[c]{0.234\textwidth}
%     \centering            \includegraphics[width=0.998055\textwidth]{figures/Global2D/epoch\iiidx_idx1_w600r30d15_f.png}
%     \subcaption{Epoch \iiidx.}
%     %		\label{subfig:f3}
%     \end{subfigure}
%     }
% \caption{The top row displays the learned neural network at epochs 25, 50, 100, and 2000, while the bottom row illustrates the differences between the learned neural network and the target function.}
%     	\label{fig:2D-oscillation-difference:and:NN}
%     % 		\vskip 0.02in
% \end{figure}

\begin{figure}%[htbp]%[h]
    \centering
    \begin{subfigure}[c]{0.234\textwidth}
        \centering
        \includegraphics[width=0.998055\textwidth]{figures/Global2D/epoch25_idx1_w600r30d15_diff.png}
        \subcaption{Epoch 25.}
    \end{subfigure}
    \hfill
    \begin{subfigure}[c]{0.234\textwidth}
        \centering
        \includegraphics[width=0.998055\textwidth]{figures/Global2D/epoch50_idx1_w600r30d15_diff.png}
        \subcaption{Epoch 50.}
    \end{subfigure}
    \hfill
    \begin{subfigure}[c]{0.234\textwidth}
        \centering
        \includegraphics[width=0.998055\textwidth]{figures/Global2D/epoch1000_idx1_w600r30d15_diff.png}
        \subcaption{Epoch 1000.}
    \end{subfigure}
    \hfill
    \begin{subfigure}[c]{0.234\textwidth}
        \centering
        \includegraphics[width=0.998055\textwidth]{figures/Global2D/epoch2000_idx1_w600r30d15_diff.png}
        \subcaption{Epoch 2000.}
    \end{subfigure}
\\
    \begin{subfigure}[c]{0.234\textwidth}
        \centering
        \includegraphics[width=0.998055\textwidth]{figures/Global2D/epoch25_idx1_w600r30d15_f.png}
        \subcaption{Epoch 25.}
    \end{subfigure}
    \hfill
    \begin{subfigure}[c]{0.234\textwidth}
        \centering
        \includegraphics[width=0.998055\textwidth]{figures/Global2D/epoch50_idx1_w600r30d15_f.png}
        \subcaption{Epoch 50.}
    \end{subfigure}
    \hfill
    \begin{subfigure}[c]{0.234\textwidth}
        \centering
        \includegraphics[width=0.998055\textwidth]{figures/Global2D/epoch1000_idx1_w600r30d15_f.png}
        \subcaption{Epoch 1000.}
    \end{subfigure}
    \hfill
    \begin{subfigure}[c]{0.234\textwidth}
        \centering
        \includegraphics[width=0.998055\textwidth]{figures/Global2D/epoch2000_idx1_w600r30d15_f.png}
        \subcaption{Epoch 2000.}
    \end{subfigure}
    
    \caption{The top row: the learned neural network;  the bottom row: error.
    % the differences between the learned neural network and the target function.
    }
    \label{fig:2D-oscillation-difference:and:NN}
\end{figure}


% \begin{figure}%[htbp]%[h]%[htbp!]
%     \centering	
%     \includegraphics[width=0.5\textwidth]{figures/Global2D/errors2D_w600r30d15.pdf}
%     % \subcaption{}
%     %		\label{subfig:f3}
%     \caption{Illustrations of training and test errors in difference epochs, measured in MSE.}
%     	\label{fig:2D-oscillation-error}
%     % 		\vskip 0.02in
% \end{figure}

We trained the same function using identical network settings, except we limited the domain of interest to a unit disc. We sampled $452^2$ data points uniformly distributed over the $[-1,1]^2$ area, then filtered to retain only those points that fall within the unit disk, totaling approximately $159692$ $(\approx 400^2)$ samples. As illustrated in Figure~\ref{fig:2D-oscillation-unitBall}, our network successfully learned the target function in the disc with no adjustments or modifications. This test highlights the network's flexibility for domain geometry, an advantage over traditional mesh or grid-based methods, especially in higher dimensions.

 \begin{figure}%[htbp]%[h]%[htbp!]
    \centering	
    \begin{subfigure}[c]{0.31246\textwidth}
    \centering            \includegraphics[height=0.6790\textwidth]{figures/Global2D/UnitBallepoch2000_idx1_w600r30d15_f.png}
    \subcaption{Approximation.}
    %		\label{subfig:f3}
    \end{subfigure}
    \hfill
    \begin{subfigure}[c]{0.31246\textwidth}
    \centering            \includegraphics[height=0.6790\textwidth]{figures/Global2D/UnitBallepoch2000_idx1_w600r30d15_diff.png}
    \subcaption{Error.}
    %		\label{subfig:f3}
    \end{subfigure}
    \hfill
    \begin{subfigure}[c]{0.31246\textwidth}
    \centering            \includegraphics[height=0.6675\textwidth]{figures/Global2D/UnitBallerrors2D_w600r30d15.pdf}
    \subcaption{Errors (in MSE).
    % vs. epoch
    }
    %		\label{subfig:f3}
    \end{subfigure}
\caption{MMNN approximation in a unit disk.}
    	\label{fig:2D-oscillation-unitBall}
    % 		\vskip 0.02in
\end{figure}

% \subsection{Tests in three dimension and higher}\label{sec:higher}
% In this section, we test a few examples in three and four dimensions. Even sampling an interesting function becomes challenging as the dimension becomes higher. Although our examples are limited by our computation power using a laptop, our tests show that MMNN performs well and is more effective than a fully connected network. 

% The first example is a 3D function a level set of which is shown in Figure~\ref{fig:f3D:spikyBall}. Using polar coordinates \((r, \theta, \phi)\), $\theta \in [0,\pi]$, $\phi\in[0,2\pi)$,
% %\[ \begin{cases}
% %    x = r\sin(\theta)\cos(\phi), \\
% %    y = r\sin(\theta)\sin(\phi), \\
% %    z = r\cos(\theta),
% %\end{cases} 
% %\quad \theta \in [0,\pi), \phi\in[0,2\pi)].
% %\]
% %where \(\theta\) ranges from \([0, \pi)\) and \(\phi\) from \([0, 2\pi)\).
% the target function \( f(x, y, z) \) is defined as:
% \begin{equation*}
%     f(r, \theta, \phi) = \begin{cases}
%         0 & \text{if } 0.5 + 5\rho - 5r \leq 0, \\
%         1 & \text{if } 0.5 + 5\rho - 5r \geq 1, \\
%         0.5 + 5\rho - 5r & \text{otherwise},
%     \end{cases}
% \end{equation*}
% % \begin{equation*}
% %     f(r, \theta, \phi) = \begin{cases}
% %         0 & \text{if } 0.5 + 10\rho - 10r \leq 0, \\
% %         1 & \text{if } 0.5 + 10\rho - 10r \geq 1, \\
% %         0.5 + 10\rho - 10r & \text{otherwise},
% %     \end{cases}
% % \end{equation*}
% % where \( \rho \) denotes a parameter related to the density or distribution within the spiky ball structure.
% where
% \begin{equation*}
%     \rho=\rho(\theta,\phi)=0.5+0.2\sin(6  \theta)\cos(6 \phi)\sin^2(\theta).
% \end{equation*}
% % \red{For examples, $f(r,\theta,\phi)=0.5$ corresponds to $r=r(\theta,\phi)=0.5+0.2\sin(6  \theta)\cos(6 \phi)\sin^2(\theta)$, i.e.,
% % \begin{equation*}
% %     \tn{$x = r\sin(\theta)\cos(\phi)$,\quad
% %    $y = r\sin(\theta)\sin(\phi)$,\quad
% %    $z = r\cos(\theta)$\qquad for $\theta \in [0,\pi]$,\quad $\phi\in[0,2\pi)$.
% %    }
% % \end{equation*}
% % % and $f(r,\theta,\phi)=0.1$.
% % % and $f(r,\theta,\phi)=0.9$ correspond to $r=0.54s+0.2\sin(6  \theta)\cos(6 \phi)\sin^2(\theta)$ and $r=0.46+0.2\sin(6  \theta)\cos(6 \phi)\sin^2(\theta)$, respectively, similar to that of $f(r,\theta,\phi)=0.5$.
% % It is easy to verify that the surface $f(r,\theta,\phi)=v$ for $v\in (0,1)$ is also a spiky ball, similar to that of $f(r,\theta,\phi)=0.5$.}


% \begin{figure}%[htbp]%[h]
%     \begin{minipage}[b]{0.30\linewidth}
%     \centering    \includegraphics[width=0.9585\textwidth]{figures/HighDim/f3DspikyBall.pdf}
%     \caption{Surface plot of the levelset $f(r,\theta,\phi)=0.5$.}
%     \label{fig:f3D:spikyBall}
% \end{minipage}
% \hfill
%     \begin{minipage}[b]{0.30\linewidth}
%     \centering    \includegraphics[width=0.9585\textwidth]{figures/HighDim/f3DspikyBall.pdf}
%     \caption{Surface plot of the levelset $h(r,\theta,\phi)=0.5$.}
%     \label{fig:h3D:spikyBall}
% \end{minipage}
% \hfill
% \begin{minipage}[b]{0.33\linewidth}
%     \centering
%     \includegraphics[width=0.969985\textwidth]{figures/HighDim/logErrorVsEpochw600r20d8fun2.pdf}
%  \caption{Training and test errors (MSE) vs. epoch.}
%     \label{fig:3DlogErrorVsEpoch1}
%     \end{minipage}
% \end{figure}

% Our MMNN is of a compact size $(600, 20, 8)$. For this test, a total of $111^3$ data are sampled on a uniform grid in $[-1,1]^3$ with a mini-batch size of $999$ and a learning rate of $0.0005 \times 0.9^{\lfloor k/6 \rfloor}$ for epochs $k=1,2,\cdots,300$.
% Figure~\ref{fig:3DlogErrorVsEpoch1} gives the error plot. 
% As shown in Figures~\ref{fig:f3D:spikyBall} and \ref{fig:h3D:spikyBall}, the levelsets corresponding to the target function $f$ and the learned MMNN approximation $h$ are nearly identical.
% To visually demonstrate the quality of the approximation and complex structure of the 3D function, we present several slices of the target function and the MMNN approximation by fixing either $x$, $y$, or $z$ in Figure~\ref{fig:3Dslices}.
% %Specifically, we fix the value of $z$, such as $z = 0, 0.2, 0.4, 0.6$, and consider $h(x,y,z)$ and $f(x,y,z)$ as functions of $x$ and $y$. By examining multiple values of $z$, we can clearly observe how well $h(x,y,z)$ approximates the target function $f(x,y,z)$. 
% % \begin{figure}%[htbp]%[h]
% %     \centering
% %     \includegraphics[width=0.5\textwidth]{figures/HighDim/logErrorVsEpochw600r20d8fun2.pdf}
% %  \caption{Log of training and test errors (MSE) versus epoch.}
% %     \label{fig:3DlogErrorVsEpoch1}
% % \end{figure}


% % \begin{figure}%[htbp]%[h]%[htbp!]
% %     \centering	
% %     \begin{subfigure}[c]{0.48\textwidth}
% %     \centering            \includegraphics[width=0.9598055\textwidth]{figures/HighDim/3DsliceX1.png}
% %     \subcaption{$x=0.2$.}
% %     		% \label{subfig:f3}
% %     \end{subfigure}
% %     \hfill
% %                 \begin{subfigure}[c]{0.48\textwidth}
% %     \centering            \includegraphics[width=0.9598055\textwidth]{figures/HighDim/3DsliceX2.png}
% %     \subcaption{$x=0.4$.}
% %     		% \label{subfig:f3}
% %     \end{subfigure}\\
% %     \begin{subfigure}[c]{0.48\textwidth}
% %     \centering            \includegraphics[width=0.9598055\textwidth]{figures/HighDim/3DsliceY1.png}
% %     \subcaption{$y=0.2$.}
% %     		% \label{subfig:f3}
% %     \end{subfigure}
% %     \hfill
% %                 \begin{subfigure}[c]{0.48\textwidth}
% %     \centering            \includegraphics[width=0.9598055\textwidth]{figures/HighDim/3DsliceY2.png}
% %     \subcaption{$y=0.4$.}
% %     		% \label{subfig:f3}
% %     \end{subfigure}
% %    \\
% %     \begin{subfigure}[c]{0.48\textwidth}
% %     \centering            \includegraphics[width=0.9598055\textwidth]{figures/HighDim/3DsliceZ1.png}
% %     \subcaption{$z=0.2$.}
% %     		% \label{subfig:f3}
% %     \end{subfigure}
% %     \hfill
% %                 \begin{subfigure}[c]{0.48\textwidth}
% %     \centering            \includegraphics[width=0.9598055\textwidth]{figures/HighDim/3DsliceZ2.png}
% %     \subcaption{$z=0.4$.}
% %     		% \label{subfig:f3}
% %     \end{subfigure}
% %     \caption{Slices of the true function $f(x, y, z)$ vs. those of the MMNN approximation $h(x, y, z)$.}   \label{fig:3Dslices}
% %     % 	
% % % \caption{Comparison among our network, interpolation and FEM method. Interpolation employs $72\times 72=5184$ uniform grid points and FEM method use that many bases. Our network architecture is designed with $(100+1) \times 10 \times (6-1) + (100+1) = 5151$ degrees of freedom (parameters). The maximum error for neural networks is approximately $0.05$, for interpolation it is around $0.31$, for FEM is  $0.38$. For MSE, $0.85\times 10^{-6}$ vs 1.95e-04 vs 1.45e-04}    		\vskip 0.02in
% % \end{figure}

% \begin{figure}%[htbp]%[h]%[htbp!]
%     \centering	
%        \begin{subfigure}[c]{0.3248\textwidth}
%     \centering            \includegraphics[width=0.9869598055\textwidth]{figures/HighDim/3DsliceZ0.png}
%     \subcaption{$z=0$.}
%     		% \label{subfig:f3}
%     \end{subfigure}
%     \hfill
%     \begin{subfigure}[c]{0.3248\textwidth}
%     \centering            \includegraphics[width=0.9869598055\textwidth]{figures/HighDim/3DsliceZ1.png}
%     \subcaption{$z=0.1$.}
%     		% \label{subfig:f3}
%     \end{subfigure}\hfill             
%     \begin{subfigure}[c]{0.3248\textwidth}
%     \centering            \includegraphics[width=0.9869598055\textwidth]{figures/HighDim/3DsliceY0.png}
%     \subcaption{$y=0.2$.}
%     		% \label{subfig:f3}
%     \end{subfigure}\\
%        \begin{subfigure}[c]{0.3248\textwidth}
%     \centering            \includegraphics[width=0.9869598055\textwidth]{figures/HighDim/3DsliceY1.png}
%     \subcaption{$y=0.3$.}
%     		% \label{subfig:f3}
%     \end{subfigure}
%     \hfill
%         \begin{subfigure}[c]{0.3248\textwidth}
%     \centering            \includegraphics[width=0.9869598055\textwidth]{figures/HighDim/3DsliceX1.png}
%     \subcaption{$x=0.4$.}
%     		% \label{subfig:f3}
%     \end{subfigure}
%     \hfill
%                 \begin{subfigure}[c]{0.3248\textwidth}
%     \centering            \includegraphics[width=0.9869598055\textwidth]{figures/HighDim/3DsliceX0.png}
%     \subcaption{$x=0.5$.}
%     		% \label{subfig:f3}
%     \end{subfigure}
%     \caption{Slices of the true function $f(x, y, z)$ vs. those of the MMNN approximation $h(x, y, z)$.}   \label{fig:3Dslices}
% \end{figure}


% % \sz{Include illustrative graphs in subsequent revisions, e.g.,   surface $$g(x,y,z)=f(r,\theta, \phi)=0.3$$
% % $$g(x,y,z)=f(r,\theta, \phi)=0.7$$}




% % We first consider the pdf of gaussion (normal) distribution.  

% % The multivariate normal distribution is said to be ``non-degenerate" when the symmetric covariance matrix \(\Sigma\) is positive definite. In this case the distribution has \textit{density}
% % \[
% % f(\bmx) = f(x_1, \dots, x_k) = \frac{\exp\left(-\frac{1}{2} (x-\mu)^\ts  \Sigma^{-1} (x-\mu)\right)}{\sqrt{(2\pi)^k \det(\Sigma)}}
% % \]

% % Let the target function be the pdf of gaussion (normal) distribution.  
% % For simplicity we mannuly give the inverse of  symmetric covariance matrix \(\Sigma\), which could avoid computing inverse of a matrix. That is, 
% % set $\bmmu=\bmzero$ 
% % and $\Sigam^{-1}=$ 3d example.
% Next, we consider the probability density function (PDF) of a Gaussian (normal) distribution in 4D,
% \[
% f(\bm{x}) = f(x_1, \dots, x_4) = \frac{\exp\left(-\frac{1}{2} (\bm{x}-\boldsymbol{\mu})^\ts  \bmSigma^{-1} (\bm{x}-\bm{\mu})\right)}{\sqrt{(2\pi)^k \det(\bmSigma)}}
% \]
% where \(\bmSigma\) is the covariance matrix. 
% % We begin by examining the probability density function (pdf) of the Gaussian (normal) distribution. The multivariate normal distribution is classified as ``non-degenerate" when the symmetric covariance matrix \(\Sigma\) is positive definite. Under this condition, the distribution possesses the following \textit{density}:
% % \[
% % f(\bm{x}) = f(x_1, \dots, x_k) = \frac{\exp\left(-\frac{1}{2} (\bm{x}-\boldsymbol{\mu})^\ts  \Sigma^{-1} (\bm{x}-\boldsymbol{\mu})\right)}{\sqrt{(2\pi)^k \det(\Sigma)}}
% % \]
% We set \(\bm{\mu} = \bm{0}\)  and
% $    \bmSigma^{-1}=20\left[\begin{smallmatrix}
%     1.0 & 0.9 & 0.8 & 0.7 \\
%     0.9 & 2.0 & 1.9 & 1.8 \\
%     0.8 & 1.9 & 3.0 & 2.9 \\
%     0.7 & 1.8 & 2.9 & 4.0 \\
% \end{smallmatrix}\right].$
% % \begin{equation*}
% %     \bmSigma^{-1}=20\begin{bmatrix}
% %     1.0 & 0.9 & 0.8 & 0.7 \\
% %     0.9 & 2.0 & 1.9 & 1.8 \\
% %     0.8 & 1.9 & 3.0 & 2.9 \\
% %     0.7 & 1.8 & 2.9 & 4.0 \\
% % \end{bmatrix}.
% % \end{equation*}
% We remark that the eigenvalues of 
% $\bmSigma^{-1}$ are $6.82,  9.93, 25.28, 158.05$ which means that the distribution is quite anisotropic and concentrated near the center.
% % follows:
% % \begin{itemize}
% %     \item For a 3-dimensional (3D) case:
% % \[
% % \bmSigma^{-1} = 0.05\begin{bmatrix}
% % 1.0 & 0.9 & 0.8 \\
% % 0.9 & 2.0 & 1.9 \\
% % 0.8 & 1.9 & 3.0 \\
% % \end{bmatrix}
% % \]
% % \item
% % For a 5-dimensional (5D) case:
% % \[
% % \bmSigma^{-1} = 0.05\begin{bmatrix}
% % 1.0 & 0.9 & 0.8 & 0.7 & 0.6 \\
% % 0.9 & 2.0 & 1.9 & 1.8 & 1.7 \\
% % 0.8 & 1.9 & 3.0 & 2.9 & 2.8 \\
% % 0.7 & 1.8 & 2.9 & 4.0 & 3.9 \\
% % 0.6 & 1.7 & 2.8 & 3.9 & 5.0 \\
% % \end{bmatrix}
% % \]
% % \end{itemize}

% A compact MMNN with size of $(500,12,6)$ produces a good approximation as shown in the error plot Figure~\ref{fig:logErrorVsEpoch2}. Figure~\ref{fig:4D:f:vs:h} compares the true function $f(x, y, z, u)$ and  the MMNN approximation $h(x, y, z, u)$ with $z=u=0.2$. For this test a total of $35^4$ data are sampled on a uniform grid in $[-1,1]^4$ with a mini-batch size of $35^2$ and a learning rate at $10^{-3} \times 0.9^{\lfloor k/6 \rfloor}$ for epochs $k=1,2,\cdots,300$.

% \begin{figure}%[htbp]%[h]
% \begin{minipage}[b]{0.5861\linewidth}
%     \centering
%     \includegraphics[width=0.94569965\textwidth]{figures/HighDim/4DsliceZU0.png}
%     \caption{True function $f(x, y, z, u)$ versus the learned network $h(x, y, z, u)$ with $z=u=0.2$.}  
%     \label{fig:4D:f:vs:h}
% \end{minipage}
% \hfill
% \begin{minipage}[b]{0.3565\linewidth}
%     \centering
%     \includegraphics[width=0.97965\textwidth]{figures/HighDim/logErrorVsEpochw500r12d6fun1.pdf}
%  \caption{Training and test errors (MSE) vs. epoch.}
%     \label{fig:logErrorVsEpoch2}
%     \end{minipage}
% \end{figure}


% %The illustration in Figure~\ref{fig:logErrorVsEpoch2} demonstrates that the network $h(x, y, z)$ successfully approximates the target function $f(x, y, z)$, achieving an impressively low MSE error of $1.02 \times 10^{-5}$ and a maximum error of $0.15$.

% % Our architecture can learn global high frequency, e.g., $\sin(50\pi x)$, 2D $\sum_{i=1}^4 \cos \sin$; 2D star;


% % \subsubsection{3D}

% % Will add later


% % \subsection{Local data sampling}
% % Target function 
% % \[
% % f(x)=\one_{\{|x|<0.02\}} \cdot\sin(
% %         50\pi  x)
% % \]

% % 18000 samples linspace(-0.06,0.06,18000)

% % 2000 samples linspace(-1,1,2000)
% % \begin{figure}%[htbp]%[h]%[htbp!]
% %     \centering	
% %     \foreach \iiidx in {7,8,9,10,15,20,25,30,60,90,...,160}
% %     {
% %     \begin{subfigure}[c]{0.31\textwidth}
% %     \centering            \includegraphics[width=0.998055\textwidth]{figures/LDS1D/epoch\iiidx_idx1.png}
% %     \subcaption{Epoch \iiidx.}
% %     %		\label{subfig:f3}
% %     \end{subfigure}
% %     }
% %     \caption{Illustrations of learned NN and the target function.}
% %     	% \label{fig:}
% %     % 		\vskip 0.02in
% % \end{figure}




% \clearpage
\subsection{Discontinuous functions with porous structures}
\label{sec:discontinuous:func:porous}


Earlier, we tested highly oscillatory functions that are still continuous. To be more inclusive in our experiments, we now consider piecewise smooth functions with porous structures and complicated interfaces. Two target functions are shown in Figure~\ref{fig:discont:funcs}(a,d). The first is a constant function with holes of various shapes removed (piecewise constant), while the second is based on the function in Figure~\ref{fig:vsFCNN:f2D} with holes introduced (piecewise smooth). We choose an MMNN of size (256, 12, 6), denoted MMNN1, and another of size (1024, 32, 6), denoted MMNN2, to learn these two functions, respectively. For training, we sample \(600^2\) data points on a uniform grid in \([-1,1]^2\), using a mini-batch size of 1000 and a learning rate of \(0.001 \times 0.9^{\lfloor k/20 \rfloor}\) for epochs \(k = 1, 2, \dots, 1600\).



\begin{figure}%[htbp]%[h]%[htbp!]
    \centering	
    \begin{subfigure}[b]{0.294311\textwidth}
    \centering            \includegraphics[width=0.998055\textwidth]{figures/Discountinous/Discont_true_funidx1.png}
    \subcaption{Truth.}
    \label{subfig:discont:fun1}
    % 		\label{subfig:vsFCNN:errors:2D}
    \end{subfigure}
        \hfill
              \begin{subfigure}[b]{0.294311\textwidth}
    \centering            \includegraphics[width=0.985\textwidth]{figures/Discountinous/Discont_learned_funidx1.png}
    \subcaption{MMNN1.}
    		% \label{subfig:vsFCNN:errors:1D}
    \end{subfigure}
    \hfill
    \begin{subfigure}[b]{0.294311\textwidth}
    \centering            \includegraphics[width=0.985\textwidth]{figures/Discountinous/Discont_diff_funidx1.png}
    \subcaption{Error.}
    		% \label{subfig:vsFCNN:errors:1D}
    \end{subfigure}
    \\
         \begin{subfigure}[b]{0.294311\textwidth}
    \centering            \includegraphics[width=0.998055\textwidth]{figures/Discountinous/Discont_true_funidx2.png}
    \subcaption{Truth.}
    \label{subfig:discont:fun2}
    \end{subfigure}
    \hfill
    \begin{subfigure}[b]{0.294311\textwidth}
    \centering            \includegraphics[width=0.998055\textwidth]{figures/Discountinous/Discont_learned_funidx2.png}
    \subcaption{MMNN2.}
    
    % 		\label{subfig:vsFCNN:errors:2D}
    \end{subfigure}
    \hfill
    \begin{subfigure}[b]{0.294311\textwidth}
    \centering            \includegraphics[width=0.998055\textwidth]{figures/Discountinous/Discont_diff_funidx2.png}
    \subcaption{Error.}
    
    % 		\label{subfig:vsFCNN:errors:2D}
    \end{subfigure}
    \caption{
   True functions, corresponding learned networks, and their differences (errors). The top and bottom rows correspond to two different target functions.}
  \label{fig:discont:funcs}
    % 		\vskip 0.02in
\end{figure}


\begin{table}[htbp!]
	\centering  
 \setlength{\tabcolsep}{0.68em} 
 \renewcommand{\arraystretch}{1.15}
\caption{
Test errors for two approximation results in Figure~\ref{fig:discont:funcs}.
}
\label{tab:discont:porous}
	\resizebox{0.948\textwidth}{!}{ 
		\begin{tabular}{ccccccccc} 
			\toprule% \toprule[1.2pt]  
			    target function & network &  (width, rank, depth) & {\#parameters (trained / all)}  &      {test error (MSE)} 
    & test error (MAX)  \\
			 % \cmidrule{3-5}
			 % & & MSE & MAE & MAX\\
			\midrule
			 \rowcolor{mygray}
			Figure~\ref{fig:discont:funcs}(a) &  MMNN1  & (256, 12, 6) &  15677 / 33085 & 
$1.35 \times 10^{-3}$  &  $9.94 \times 10^{-1}$ \\	
			% \midrule
   		Figure~\ref{fig:discont:funcs}(d) &	 MMNN2  & (1024, 32, 6) & 165025 / 337057 & 
$1.14 \times 10^{-3}$  &  $1.12\times 10^{0}$
   \\	
			\bottomrule% \bottomrule[1.2pt] 
		\end{tabular} 
	}%%% \resizebox
% 	\vskip -0.1in
\end{table} 
% \red{Not perform we

% As illustrated in Figure~\ref{fig:discont:funcs}, MMNNs demonstrate impressive ability  in localizing and capturing discontinuities, geometries, and oscillations simultaneously. In other words, MMNNs can be adaptive in both space and frequency domains!
% }
As shown in Figure~\ref{fig:discont:funcs}, MMNNs demonstrate an impressive ability to simultaneously localize and capture discontinuities, geometric features, and oscillatory behaviors. This indicates that MMNNs are adaptive in both spatial and frequency domains. 
While the errors presented in Table~\ref{tab:discont:porous} are somewhat larger, they are primarily concentrated near the discontinuous parts, as illustrated in Figure~\ref{fig:discont:funcs}, which is reasonable.


\subsection{Learning dynamics}\label{sec:dynamics}


% In this section, 
Here,
we show some interesting learning dynamics observed during the training process. As the first example in Section~\ref{sec:oscillatory} and the following examples show, the training process not just learns from low frequency first but can also learn feature by feature, i.e., can be localized in both frequency domain and spatial domain. We believe this is due to the combination of MMNN's ``divide and conquer" ability and the Adam optimizer which utilizes momentum. More understanding is needed and will be studied in our future research.

We again start with a one-dimensional example, $f(x)=\sin\big(36\pi |x|^{1.5}\big), x\in[-1,1]$. 
An MMNN of size $(600,30,8)$ produces a good approximation of this highly oscillatory function, as illustrated by the error plot in Figure~\ref{fig:36pi:power1.5-error}. For this test, a total of $1000$ uniformly sampled points in $[-1,1]$ are used with a mini-batch size of $100$ and a learning rate of $10^{-3}\times 0.9^{\lfloor k/200 \rfloor}$, where $k=1,2,\cdots,10000$ is the epoch number. 
As illustrated in Figure~\ref{fig:36pi:power1.5-derivative}, the function is less oscillatory near 0. Therefore, we might anticipate that the network will initially learn the part near 0 and then feature by feature from the middle to the boundary. The experimental results presented in Figure~\ref{fig:sin:36pi:power:1.5:learning:process} agree with our expectations.

%Let's now focus on functions exhibiting global high frequency. We will limit our analysis to one-dimensional and cases because high-dimensional scenarios require substantial computational power if the target function displays global high frequency in such spaces, which is affected by the curse of dimensionality.


% First, 
%We begin with the one-dimensional case, choosing the function $\sin(50\pi x)$ on the interval $[-1,1]$ as our target. The diagram in Figure~\ref{fig:sin:50pi:NN} clearly shows that learning starts at the boundary, consistent with our earlier theoretical analysis.

% \begin{figure}%[htbp]%[h]%[htbp!]
% \centering 
%         \begin{subfigure}[c]{0.19272431\textwidth}
%    \centering	    \includegraphics[width=0.808333\textwidth]{figures/Dynamics/LDfunDerivative.pdf}
%     % \subcaption{}
%     %		\label{subfig:f3}
%     \caption{$f'$.}
%     	\label{fig:36pi:power1.5-derivative}
%     \end{subfigure}
%     \hfill
%     \begin{subfigure}[c]{0.19272431\textwidth}
% \centering\includegraphics[width=0.7859034599\textwidth]{figures/Dynamics/errorsLD1D_w600r30d8.pdf}
%     % \subcaption{}
%     %		\label{subfig:f3}
%     \caption{Errors vs. epoch.}
%     	\label{fig:36pi:power1.5-error}
%     \end{subfigure}
%     \hfill
%     \begin{subfigure}[c]{0.19272431\textwidth}
%         \centering
%         \includegraphics[width=0.99\textwidth]{figures/Dynamics/epoch400_idx0_w600r30d8.pdf}
%         \subcaption{Epoch 400.}
%     \end{subfigure}
%     \hfill
%     \begin{subfigure}[c]{0.19272431\textwidth}
%         \centering
%         \includegraphics[width=0.99\textwidth]{figures/Dynamics/epoch500_idx0_w600r30d8.pdf}
%         \subcaption{Epoch 500.}
%     \end{subfigure}
%     \hfill
%     \begin{subfigure}[c]{0.19272431\textwidth}
%         \centering
%         \includegraphics[width=0.99\textwidth]{figures/Dynamics/epoch600_idx0_w600r30d8.pdf}
%         \subcaption{Epoch 600.}
%     \end{subfigure}
%     \hfill
%     \begin{subfigure}[c]{0.19272431\textwidth}
%         \centering
%         \includegraphics[width=0.99\textwidth]{figures/Dynamics/epoch700_idx0_w600r30d8.pdf}
%         \subcaption{Epoch 700.}
%     \end{subfigure}
% \hfill
%     \begin{subfigure}[c]{0.19272431\textwidth}
%         \centering
%         \includegraphics[width=0.99\textwidth]{figures/Dynamics/epoch1000_idx0_w600r30d8.pdf}
%         \subcaption{Epoch 1000.}
%     \end{subfigure}
%     \hfill
%     \begin{subfigure}[c]{0.19272431\textwidth}
%         \centering
%         \includegraphics[width=0.99\textwidth]{figures/Dynamics/epoch1700_idx0_w600r30d8.pdf}
%         \subcaption{Epoch 1700.}
%     \end{subfigure}
%     \hfill
%     \begin{subfigure}[c]{0.19272431\textwidth}
%         \centering
%         \includegraphics[width=0.99\textwidth]{figures/Dynamics/epoch2400_idx0_w600r30d8.pdf}
%         \subcaption{Epoch 2400.}
%     \end{subfigure}
%     \hfill
%     \begin{subfigure}[c]{0.19272431\textwidth}
%         \centering
%         \includegraphics[width=0.99\textwidth]{figures/Dynamics/epoch10000_idx0_w600r30d8.pdf}
%         \subcaption{Epoch 10000.}
%     \end{subfigure}
    
%     \caption{Illustration of the training process.}
%     \label{fig:sin:36pi:power:1.5:learning:process}
% \end{figure}

\begin{figure}[h]%[htbp!]
\centering
\begin{minipage}[b]{0.4568\linewidth}
    \centering	    \includegraphics[width=0.59\textwidth]{figures/Dynamics/LDfunDerivative.pdf}
    % \subcaption{}
    %		\label{subfig:f3}
    \caption{Derivative of $f$.
    % $f(x)=\sin\big(36\pi |x|^{1.5}\big)$.
    }
    	\label{fig:36pi:power1.5-derivative}
    % 		\vskip 0.02in
    \end{minipage}
    \hspace{8pt}
    % \hfill
    \begin{minipage}[b]{0.4568\linewidth}        \centering\includegraphics[width=0.59\textwidth]{figures/Dynamics/errorsLD1D_w600r30d8.pdf}
    % \subcaption{}
    %		\label{subfig:f3}
    \caption{Errors (in MSE) vs. epoch.}
    	\label{fig:36pi:power1.5-error}
    % 		\vskip 0.02in
    \end{minipage}
\end{figure}


% \begin{figure}%[htbp]%[h]%[htbp!]
%     \centering	
%     \foreach \iiidx in {120,180,240,300,700,1000,1300,5000}
%     {
%     \begin{subfigure}[c]{0.231\textwidth}
%     \centering            \includegraphics[width=0.99\textwidth]{figures/Dynamics/epoch\iiidx_idx0_w600r30d8.pdf}
%     \subcaption{Epoch \iiidx.}
%     %		\label{subfig:f3}
%     \end{subfigure}
%     }
%     \caption{Illustrations of the training process.}
%     	\label{fig:sin:36pi:power:1.5:learning:process}
%     % 		\vskip 0.02in
% \end{figure}
% \begin{figure}%[htbp]%[h]
%     \centering
%     \begin{subfigure}[c]{0.231\textwidth}
%         \centering
%         \includegraphics[width=0.99\textwidth]{figures/Dynamics/epoch120_idx0_w600r30d8.pdf}
%         \subcaption{Epoch 120.}
%     \end{subfigure}
%     \hfill
%     \begin{subfigure}[c]{0.231\textwidth}
%         \centering
%         \includegraphics[width=0.99\textwidth]{figures/Dynamics/epoch180_idx0_w600r30d8.pdf}
%         \subcaption{Epoch 180.}
%     \end{subfigure}
%     \hfill
%     \begin{subfigure}[c]{0.231\textwidth}
%         \centering
%         \includegraphics[width=0.99\textwidth]{figures/Dynamics/epoch240_idx0_w600r30d8.pdf}
%         \subcaption{Epoch 240.}
%     \end{subfigure}
%     \hfill
%     \begin{subfigure}[c]{0.231\textwidth}
%         \centering
%         \includegraphics[width=0.99\textwidth]{figures/Dynamics/epoch300_idx0_w600r30d8.pdf}
%         \subcaption{Epoch 300.}
%     \end{subfigure}
% \\
%     \begin{subfigure}[c]{0.231\textwidth}
%         \centering
%         \includegraphics[width=0.99\textwidth]{figures/Dynamics/epoch700_idx0_w600r30d8.pdf}
%         \subcaption{Epoch 700.}
%     \end{subfigure}
%     \hfill
%     \begin{subfigure}[c]{0.231\textwidth}
%         \centering
%         \includegraphics[width=0.99\textwidth]{figures/Dynamics/epoch1000_idx0_w600r30d8.pdf}
%         \subcaption{Epoch 1000.}
%     \end{subfigure}
%     \hfill
%     \begin{subfigure}[c]{0.231\textwidth}
%         \centering
%         \includegraphics[width=0.99\textwidth]{figures/Dynamics/epoch1300_idx0_w600r30d8.pdf}
%         \subcaption{Epoch 1300.}
%     \end{subfigure}
%     \hfill
%     \begin{subfigure}[c]{0.231\textwidth}
%         \centering
%         \includegraphics[width=0.99\textwidth]{figures/Dynamics/epoch5000_idx0_w600r30d8.pdf}
%         \subcaption{Epoch 5000.}
%     \end{subfigure}
    
%     \caption{Illustration of the training process.}
%     \label{fig:sin:36pi:power:1.5:learning:process}
% \end{figure}

\begin{figure}%[htbp]%[h]
    \centering
    \begin{subfigure}[c]{0.2431\textwidth}
        \centering
        \includegraphics[width=0.99\textwidth]{figures/Dynamics/epoch400_idx0_w600r30d8.pdf}
        \subcaption{Epoch 400.}
    \end{subfigure}
    \hfill
    \begin{subfigure}[c]{0.2431\textwidth}
        \centering
        \includegraphics[width=0.99\textwidth]{figures/Dynamics/epoch500_idx0_w600r30d8.pdf}
        \subcaption{Epoch 500.}
    \end{subfigure}
    \hfill
    \begin{subfigure}[c]{0.2431\textwidth}
        \centering
        \includegraphics[width=0.99\textwidth]{figures/Dynamics/epoch600_idx0_w600r30d8.pdf}
        \subcaption{Epoch 600.}
    \end{subfigure}
    \hfill
    \begin{subfigure}[c]{0.2431\textwidth}
        \centering
        \includegraphics[width=0.99\textwidth]{figures/Dynamics/epoch700_idx0_w600r30d8.pdf}
        \subcaption{Epoch 700.}
    \end{subfigure}
\\
    \begin{subfigure}[c]{0.2431\textwidth}
        \centering
        \includegraphics[width=0.99\textwidth]{figures/Dynamics/epoch1000_idx0_w600r30d8.pdf}
        \subcaption{Epoch 1000.}
    \end{subfigure}
    \hfill
    \begin{subfigure}[c]{0.2431\textwidth}
        \centering
        \includegraphics[width=0.99\textwidth]{figures/Dynamics/epoch1700_idx0_w600r30d8.pdf}
        \subcaption{Epoch 1700.}
    \end{subfigure}
    \hfill
    \begin{subfigure}[c]{0.2431\textwidth}
        \centering
        \includegraphics[width=0.99\textwidth]{figures/Dynamics/epoch2400_idx0_w600r30d8.pdf}
        \subcaption{Epoch 2400.}
    \end{subfigure}
    \hfill
    \begin{subfigure}[c]{0.2431\textwidth}
        \centering
        \includegraphics[width=0.99\textwidth]{figures/Dynamics/epoch10000_idx0_w600r30d8.pdf}
        \subcaption{Epoch 10000.}
    \end{subfigure}
    
    \caption{Illustration of the training process.}
    \label{fig:sin:36pi:power:1.5:learning:process}
\end{figure}


% \begin{figure}%[htbp]%[h]%[htbp!]
%     \centering	
%     \foreach \iiidx in {120,180,240,300,700,1000,1300,5000}
%     {
%     \begin{subfigure}[c]{0.31\textwidth}
%     \centering            \includegraphics[width=0.99\textwidth]{figures/Dynamics/epoch\iiidx_idx0_w600r30d8.pdf}
%     \subcaption{Epoch \iiidx.}
%     %		\label{subfig:f3}
%     \end{subfigure}
%     }
%     \begin{subfigure}[c]{0.31\textwidth}
%     \centering            \includegraphics[width=0.99\textwidth]{figures/Dynamics/errorsLD1D_w600r30d8.pdf}
%     \subcaption{Errors.}
%     %		\label{subfig:f3}
%     \end{subfigure}
%     \caption{Illustrations of the training process.}
%     	\label{fig:sin:36pi:power:1.5}
%     % 		\vskip 0.02in
% \end{figure}
% \begin{figure}%[htbp]%[h]%[htbp!]
%     \centering	
%     \includegraphics[width=0.599\textwidth]{figures/Dynamics/errorsLD1D_w600r30d8.pdf}
%     % \subcaption{}
%     %		\label{subfig:f3}
%     \caption{Illustrations of training and test errors measured in MSE.}
%     	\label{fig:36pi:power1.5-error}
%     % 		\vskip 0.02in
% \end{figure}

% \begin{figure}%[htbp]%[h]%[htbp!]
% \begin{minipage}[b]{0.68\linewidth}
%     \centering	
%     \foreach \iiidx in {120,180,240,300,700,1000,1300,5000}
%     {
%     \begin{subfigure}[c]{0.31\textwidth}
%     \centering            \includegraphics[width=0.99\textwidth]{figures/Dynamics/epoch\iiidx_idx0_w600r30d8.pdf}
%     \subcaption{Epoch \iiidx.}
%     %		\label{subfig:f3}
%     \end{subfigure}
%     }
%     \caption{Illustrations of the training process.}
%     	\label{fig:sin:36pi:power:1.5}
%     % 		\vskip 0.02in
%     \end{minipage}
%     \hfill
%     \begin{minipage}[b]{0.268\linewidth}
%     \centering	
%     \includegraphics[width=0.99\textwidth]{figures/Dynamics/errorsLD1D_w600r30d8.pdf}
%     % \subcaption{}
%     %		\label{subfig:f3}
%     \caption{Illustrations of training and test errors measured in MSE.}
%     	\label{fig:36pi:power1.5-error}
%     % 		\vskip 0.02in
%     \end{minipage}
% \end{figure}




Now we show an example of two-dimensional function $f(r, \theta)$ (see Figure~\ref{fig:LD:fun2D}) defined in polar coordinates $(r, \theta)$ as
\begin{equation*}
    f(r,\theta)=\begin{cases}
        0 & \tn{if  } 0.5+5\rho-5r\le 0,\\
        1 & \tn{if  } 0.5+5\rho-5r\ge 1,\\
        0.5+5\rho-5r & \tn{otherwise},\\
    \end{cases}
    \quad \tn{where}\quad \rho=0.5+0.1\cos(\pi^2 \theta^2).
\end{equation*}

\begin{figure}%[htbp]%[h]%[htbp!]
\begin{minipage}[b]{0.6357604\linewidth}
    \centering	
    \begin{subfigure}[c]{0.458\textwidth}
    \centering            \includegraphics[height=0.795\textwidth]{figures/Dynamics/LDfun2D.png}
    % \subcaption{}
    %		\label{subfig:f3}
    \end{subfigure}
    \hfill
    \begin{subfigure}[c]{0.458\textwidth}
    \centering            \includegraphics[height=0.7925\textwidth]{figures/Dynamics/LDfun3D.pdf}
    % \subcaption{}
    %		\label{subfig:f3}
    \end{subfigure}
    \caption{Illustration of the target function.}
    	\label{fig:LD:fun2D}
    % 		\vskip 0.02in
    \end{minipage}
    \hfill
\begin{minipage}[b]{0.318\linewidth}
    \centering	
    \includegraphics[width=0.95\textwidth]{figures/Dynamics/errorsLD2D_w500r20d8.pdf}
    % \subcaption{}
    %		\label{subfig:f3}
    \caption{Training and test errors (in MSE) vs. epoch.}
    	\label{fig:errors:LD:fun2D}
    % 		\vskip 0.02in
    \end{minipage}
\end{figure}

Our MMNN is of a compact size $(500, 20, 8)$. For this test, a total of $600^2$ data are sampled on a uniform grid in $[-1,1]^2$ with a mini-batch size of $1000$ and a learning rate of $0.001 \times 0.9^{\lfloor k/6 \rfloor}$ for epochs $k=1,2,\cdots,300$.
Figure~\ref{fig:errors:LD:fun2D} gives the error plot. The training process shown in Figure~\ref{fig:2D-dynamics} illustrates that an overall coarse scale or low-frequency component of the shape is learned first and then localized features are learned one by one from coarse to fine.



% \clearpage
%%%%%%%%%%%%%%%%%%%%%%%%%%%%%%%%%
%%%%%%%%%%%%%%%%%%%%%%%%%%%%%%%%%%%%%
\subsection{Tests in three dimensions and higher}\label{sec:higher}
% In this section, 
Here,
we test a few examples in three and four dimensions. Even sampling an interesting function becomes challenging as the dimension becomes higher. Although our examples are limited by our computation power using a laptop, our tests show that MMNN performs well and is more effective than a fully connected network. 

The first example is a three-dimensional function a level set of which is shown in Figure~\ref{fig:f3D:spikyBall}. Using polar coordinates \((r, \theta, \phi)\), $\theta \in [0,\pi]$, $\phi\in[0,2\pi)$,
%\[ \begin{cases}
%    x = r\sin(\theta)\cos(\phi), \\
%    y = r\sin(\theta)\sin(\phi), \\
%    z = r\cos(\theta),
%\end{cases} 
%\quad \theta \in [0,\pi), \phi\in[0,2\pi)].
%\]
%where \(\theta\) ranges from \([0, \pi)\) and \(\phi\) from \([0, 2\pi)\).
the target function \( f(x, y, z) \) is defined as:
\begin{equation*}
    f(r, \theta, \phi) = \begin{cases}
        0 & \text{if } 0.5 + 5\rho - 5r \leq 0, \\
        1 & \text{if } 0.5 + 5\rho - 5r \geq 1, \\
        0.5 + 5\rho - 5r & \text{otherwise},
    \end{cases}
\end{equation*}
% \begin{equation*}
%     f(r, \theta, \phi) = \begin{cases}
%         0 & \text{if } 0.5 + 10\rho - 10r \leq 0, \\
%         1 & \text{if } 0.5 + 10\rho - 10r \geq 1, \\
%         0.5 + 10\rho - 10r & \text{otherwise},
%     \end{cases}
% \end{equation*}
% where \( \rho \) denotes a parameter related to the density or distribution within the spiky ball structure.
where
\begin{equation*}
    \rho=\rho(\theta,\phi)=0.5+0.2\sin(6  \theta)\cos(6 \phi)\sin^2(\theta).
\end{equation*}
% \red{For examples, $f(r,\theta,\phi)=0.5$ corresponds to $r=r(\theta,\phi)=0.5+0.2\sin(6  \theta)\cos(6 \phi)\sin^2(\theta)$, i.e.,
% \begin{equation*}
%     \tn{$x = r\sin(\theta)\cos(\phi)$,\quad
%    $y = r\sin(\theta)\sin(\phi)$,\quad
%    $z = r\cos(\theta)$\qquad for $\theta \in [0,\pi]$,\quad $\phi\in[0,2\pi)$.
%    }
% \end{equation*}
% % and $f(r,\theta,\phi)=0.1$.
% % and $f(r,\theta,\phi)=0.9$ correspond to $r=0.54s+0.2\sin(6  \theta)\cos(6 \phi)\sin^2(\theta)$ and $r=0.46+0.2\sin(6  \theta)\cos(6 \phi)\sin^2(\theta)$, respectively, similar to that of $f(r,\theta,\phi)=0.5$.
% It is easy to verify that the surface $f(r,\theta,\phi)=v$ for $v\in (0,1)$ is also a spiky ball, similar to that of $f(r,\theta,\phi)=0.5$.}


\begin{figure}%[htbp]%[h]
    \begin{minipage}[b]{0.30\linewidth}
    \centering    \includegraphics[width=0.9585\textwidth]{figures/HighDim/f3DspikyBall.pdf}
    \caption{Surface of the levelset $f(r,\theta,\phi)=0.5$.}
    \label{fig:f3D:spikyBall}
\end{minipage}
\hfill
    \begin{minipage}[b]{0.30\linewidth}
    \centering    \includegraphics[width=0.9585\textwidth]{figures/HighDim/f3DspikyBall.pdf}
    \caption{Surface of the levelset $h(r,\theta,\phi)=0.5$.}
    \label{fig:h3D:spikyBall}
\end{minipage}
\hfill
\begin{minipage}[b]{0.33\linewidth}
    \centering
    \includegraphics[width=0.969985\textwidth]{figures/HighDim/logErrorVsEpochw600r20d8fun2.pdf}
 \caption{Training and test errors (MSE) vs. epoch.}
    \label{fig:3DlogErrorVsEpoch1}
    \end{minipage}
\end{figure}

Our MMNN is of a compact size $(600, 20, 8)$. For this test, a total of $111^3$ data are sampled on a uniform grid in $[-1,1]^3$ with a mini-batch size of $999$ and a learning rate of $0.0005 \times 0.9^{\lfloor k/6 \rfloor}$ for epochs $k=1,2,\cdots,300$.
Figure~\ref{fig:3DlogErrorVsEpoch1} gives the error plot. 
As shown in Figures~\ref{fig:f3D:spikyBall} and \ref{fig:h3D:spikyBall}, the levelsets corresponding to the target function $f$ and the learned MMNN approximation $h$ are nearly identical.
To visually demonstrate the quality of the approximation and complex structure of the three-dimensional function, we present several slices of the target function and the MMNN approximation by fixing either $x$, $y$, or $z$ in Figure~\ref{fig:3Dslices}.
%Specifically, we fix the value of $z$, such as $z = 0, 0.2, 0.4, 0.6$, and consider $h(x,y,z)$ and $f(x,y,z)$ as functions of $x$ and $y$. By examining multiple values of $z$, we can clearly observe how well $h(x,y,z)$ approximates the target function $f(x,y,z)$. 
% \begin{figure}%[htbp]%[h]
%     \centering
%     \includegraphics[width=0.5\textwidth]{figures/HighDim/logErrorVsEpochw600r20d8fun2.pdf}
%  \caption{Log of training and test errors (MSE) versus epoch.}
%     \label{fig:3DlogErrorVsEpoch1}
% \end{figure}


% \begin{figure}%[htbp]%[h]%[htbp!]
%     \centering	
%     \begin{subfigure}[c]{0.48\textwidth}
%     \centering            \includegraphics[width=0.9598055\textwidth]{figures/HighDim/3DsliceX1.png}
%     \subcaption{$x=0.2$.}
%     		% \label{subfig:f3}
%     \end{subfigure}
%     \hfill
%                 \begin{subfigure}[c]{0.48\textwidth}
%     \centering            \includegraphics[width=0.9598055\textwidth]{figures/HighDim/3DsliceX2.png}
%     \subcaption{$x=0.4$.}
%     		% \label{subfig:f3}
%     \end{subfigure}\\
%     \begin{subfigure}[c]{0.48\textwidth}
%     \centering            \includegraphics[width=0.9598055\textwidth]{figures/HighDim/3DsliceY1.png}
%     \subcaption{$y=0.2$.}
%     		% \label{subfig:f3}
%     \end{subfigure}
%     \hfill
%                 \begin{subfigure}[c]{0.48\textwidth}
%     \centering            \includegraphics[width=0.9598055\textwidth]{figures/HighDim/3DsliceY2.png}
%     \subcaption{$y=0.4$.}
%     		% \label{subfig:f3}
%     \end{subfigure}
%    \\
%     \begin{subfigure}[c]{0.48\textwidth}
%     \centering            \includegraphics[width=0.9598055\textwidth]{figures/HighDim/3DsliceZ1.png}
%     \subcaption{$z=0.2$.}
%     		% \label{subfig:f3}
%     \end{subfigure}
%     \hfill
%                 \begin{subfigure}[c]{0.48\textwidth}
%     \centering            \includegraphics[width=0.9598055\textwidth]{figures/HighDim/3DsliceZ2.png}
%     \subcaption{$z=0.4$.}
%     		% \label{subfig:f3}
%     \end{subfigure}
%     \caption{Slices of the true function $f(x, y, z)$ vs. those of the MMNN approximation $h(x, y, z)$.}   \label{fig:3Dslices}
%     % 	
% % \caption{Comparison among our network, interpolation and FEM method. Interpolation employs $72\times 72=5184$ uniform grid points and FEM method use that many bases. Our network architecture is designed with $(100+1) \times 10 \times (6-1) + (100+1) = 5151$ degrees of freedom (parameters). The maximum error for neural networks is approximately $0.05$, for interpolation it is around $0.31$, for FEM is  $0.38$. For MSE, $0.85\times 10^{-6}$ vs 1.95e-04 vs 1.45e-04}    		\vskip 0.02in
% \end{figure}

\begin{figure}%[htbp]%[h]%[htbp!]
    \centering	
       \begin{subfigure}[c]{0.3248\textwidth}
    \centering            \includegraphics[width=0.9869598055\textwidth]{figures/HighDim/3DsliceZ0.png}
    \subcaption{$z=0$.}
    		% \label{subfig:f3}
    \end{subfigure}
    \hfill
    \begin{subfigure}[c]{0.3248\textwidth}
    \centering            \includegraphics[width=0.9869598055\textwidth]{figures/HighDim/3DsliceZ1.png}
    \subcaption{$z=0.1$.}
    		% \label{subfig:f3}
    \end{subfigure}\hfill             
    \begin{subfigure}[c]{0.3248\textwidth}
    \centering            \includegraphics[width=0.9869598055\textwidth]{figures/HighDim/3DsliceY0.png}
    \subcaption{$y=0.2$.}
    		% \label{subfig:f3}
    \end{subfigure}\\
       \begin{subfigure}[c]{0.3248\textwidth}
    \centering            \includegraphics[width=0.9869598055\textwidth]{figures/HighDim/3DsliceY1.png}
    \subcaption{$y=0.3$.}
    		% \label{subfig:f3}
    \end{subfigure}
    \hfill
        \begin{subfigure}[c]{0.3248\textwidth}
    \centering            \includegraphics[width=0.9869598055\textwidth]{figures/HighDim/3DsliceX1.png}
    \subcaption{$x=0.4$.}
    		% \label{subfig:f3}
    \end{subfigure}
    \hfill
                \begin{subfigure}[c]{0.3248\textwidth}
    \centering            \includegraphics[width=0.9869598055\textwidth]{figures/HighDim/3DsliceX0.png}
    \subcaption{$x=0.5$.}
    		% \label{subfig:f3}
    \end{subfigure}
    \caption{Slices of the true function $f(x, y, z)$ vs. those of the MMNN approximation $h(x, y, z)$.}   \label{fig:3Dslices}
\end{figure}


% \sz{Include illustrative graphs in subsequent revisions, e.g.,   surface $$g(x,y,z)=f(r,\theta, \phi)=0.3$$
% $$g(x,y,z)=f(r,\theta, \phi)=0.7$$}




% We first consider the pdf of gaussion (normal) distribution.  

% The multivariate normal distribution is said to be ``non-degenerate" when the symmetric covariance matrix \(\Sigma\) is positive definite. In this case the distribution has \textit{density}
% \[
% f(\bmx) = f(x_1, \dots, x_k) = \frac{\exp\left(-\frac{1}{2} (x-\mu)^\ts  \Sigma^{-1} (x-\mu)\right)}{\sqrt{(2\pi)^k \det(\Sigma)}}
% \]

% Let the target function be the pdf of gaussion (normal) distribution.  
% For simplicity we mannuly give the inverse of  symmetric covariance matrix \(\Sigma\), which could avoid computing inverse of a matrix. That is, 
% set $\bmmu=\bmzero$ 
% and $\Sigam^{-1}=$ 3d example.
Next, we consider the probability density function (PDF) of a Gaussian (normal) distribution in four dimensions
\[
f(\bm{x}) = f(x_1, \dots, x_4) = \frac{\exp\left(-\frac{1}{2} (\bm{x}-\boldsymbol{\mu})^\ts  \bmSigma^{-1} (\bm{x}-\bm{\mu})\right)}{\sqrt{(2\pi)^k \det(\bmSigma)}}
\]
where \(\bmSigma\) is the covariance matrix. 
% We begin by examining the probability density function (pdf) of the Gaussian (normal) distribution. The multivariate normal distribution is classified as ``non-degenerate" when the symmetric covariance matrix \(\Sigma\) is positive definite. Under this condition, the distribution possesses the following \textit{density}:
% \[
% f(\bm{x}) = f(x_1, \dots, x_k) = \frac{\exp\left(-\frac{1}{2} (\bm{x}-\boldsymbol{\mu})^\ts  \Sigma^{-1} (\bm{x}-\boldsymbol{\mu})\right)}{\sqrt{(2\pi)^k \det(\Sigma)}}
% \]
We set \(\bm{\mu} = \bm{0}\)  and
% $    \bmSigma^{-1}=20\left[\begin{smallmatrix}
%     1.0 & 0.9 & 0.8 & 0.7 \\
%     0.9 & 2.0 & 1.9 & 1.8 \\
%     0.8 & 1.9 & 3.0 & 2.9 \\
%     0.7 & 1.8 & 2.9 & 4.0 \\
% \end{smallmatrix}\right].$
\begin{equation*}
    \bmSigma^{-1}=20\begin{bmatrix}
    1.0 & 0.9 & 0.8 & 0.7 \\
    0.9 & 2.0 & 1.9 & 1.8 \\
    0.8 & 1.9 & 3.0 & 2.9 \\
    0.7 & 1.8 & 2.9 & 4.0 \\
\end{bmatrix}.
\end{equation*}
We remark that the eigenvalues of 
$\bmSigma^{-1}$ are $6.82,  9.93, 25.28, 158.05$ which means that the distribution is quite anisotropic and concentrated near the center.
% follows:
% \begin{itemize}
%     \item For a 3-dimensional (3D) case:
% \[
% \bmSigma^{-1} = 0.05\begin{bmatrix}
% 1.0 & 0.9 & 0.8 \\
% 0.9 & 2.0 & 1.9 \\
% 0.8 & 1.9 & 3.0 \\
% \end{bmatrix}
% \]
% \item
% For a 5-dimensional (5D) case:
% \[
% \bmSigma^{-1} = 0.05\begin{bmatrix}
% 1.0 & 0.9 & 0.8 & 0.7 & 0.6 \\
% 0.9 & 2.0 & 1.9 & 1.8 & 1.7 \\
% 0.8 & 1.9 & 3.0 & 2.9 & 2.8 \\
% 0.7 & 1.8 & 2.9 & 4.0 & 3.9 \\
% 0.6 & 1.7 & 2.8 & 3.9 & 5.0 \\
% \end{bmatrix}
% \]
% \end{itemize}

A compact MMNN with size of $(500,12,6)$ produces a good approximation as shown in the error plot Figure~\ref{fig:logErrorVsEpoch2}. Figure~\ref{fig:4D:f:vs:h} compares the true function $f(x, y, z, u)$ and  the MMNN approximation $h(x, y, z, u)$ with $z=u=0.2$. For this test a total of $35^4$ data are sampled on a uniform grid in $[-1,1]^4$ with a mini-batch size of $35^2$ and a learning rate at $10^{-3} \times 0.9^{\lfloor k/6 \rfloor}$ for epochs $k=1,2,\cdots,300$.

\begin{figure}%[htbp]%[h]%[ht]
\begin{minipage}[b]{0.5861\linewidth}
    \centering
    \includegraphics[width=0.94569965\textwidth]{figures/HighDim/4DsliceZU0.png}
    \caption{True function $f(x, y, z, u)$ versus the learned network $h(x, y, z, u)$ with $z=u=0.2$.}  
    \label{fig:4D:f:vs:h}
\end{minipage}
\hfill
\begin{minipage}[b]{0.3565\linewidth}
    \centering
    \includegraphics[width=0.97965\textwidth]{figures/HighDim/logErrorVsEpochw500r12d6fun1.pdf}
 \caption{Training and test errors (MSE) vs. epoch.}
    \label{fig:logErrorVsEpoch2}
    \end{minipage}
\end{figure}





% \begin{figure}%[htbp]%[h]%[htbp!]
%     \centering	
%     \foreach \iiidx in {1,3,5,7,14,22,30,300}
%     {
%     \begin{subfigure}[c]{0.23\textwidth}
%     \centering            \includegraphics[width=0.998055\textwidth]{figures/Dynamics/epoch\iiidx_idx1_w500r20d8_f.png}
%     \subcaption{Epoch \iiidx.}
%     %		\label{subfig:f3}
%     \end{subfigure}
%     }
%     \caption{Illustrations of the learned neural network in different epochs.}
%     	% \label{fig:}
%     % 		\vskip 0.02in
% \end{figure}

\begin{figure}[htbp!]%[h]
    \centering
    \begin{subfigure}[c]{0.23\textwidth}
        \centering
        \includegraphics[width=0.998055\textwidth]{figures/Dynamics/epoch1_idx1_w500r20d8_f.png}
        \subcaption{Epoch 1.}
    \end{subfigure}
    \hfill
    \begin{subfigure}[c]{0.23\textwidth}
        \centering
        \includegraphics[width=0.998055\textwidth]{figures/Dynamics/epoch3_idx1_w500r20d8_f.png}
        \subcaption{Epoch 3.}
    \end{subfigure}
    \hfill
    \begin{subfigure}[c]{0.23\textwidth}
        \centering
        \includegraphics[width=0.998055\textwidth]{figures/Dynamics/epoch5_idx1_w500r20d8_f.png}
        \subcaption{Epoch 5.}
    \end{subfigure}
    \hfill
    \begin{subfigure}[c]{0.23\textwidth}
        \centering
        \includegraphics[width=0.998055\textwidth]{figures/Dynamics/epoch7_idx1_w500r20d8_f.png}
        \subcaption{Epoch 7.}
    \end{subfigure}
\\
\begin{subfigure}[c]{0.23\textwidth}
        \centering
        \includegraphics[width=0.998055\textwidth]{figures/Dynamics/epoch14_idx1_w500r20d8_f.png}
        \subcaption{Epoch 14.}
    \end{subfigure}
    \hfill
    \begin{subfigure}[c]{0.23\textwidth}
        \centering
        \includegraphics[width=0.998055\textwidth]{figures/Dynamics/epoch22_idx1_w500r20d8_f.png}
        \subcaption{Epoch 22.}
    \end{subfigure}
    \hfill
    \begin{subfigure}[c]{0.23\textwidth}
        \centering
        \includegraphics[width=0.998055\textwidth]{figures/Dynamics/epoch30_idx1_w500r20d8_f.png}
        \subcaption{Epoch 30.}
    \end{subfigure}
    \hfill
    \begin{subfigure}[c]{0.23\textwidth}
        \centering
        \includegraphics[width=0.998055\textwidth]{figures/Dynamics/epoch300_idx1_w500r20d8_f.png}
        \subcaption{Epoch 300.}
    \end{subfigure}
    
    \caption{Illustration of the learning dynamics.}
    \label{fig:2D-dynamics}
\end{figure}

% \clearpage
% \red{
% Note that all settings for FMMNN remain identical to those of ReLU MMNN, except that the activation function is replaced from ReLU to sine.

% FMMNN vs ReLU MMNN:

% \begin{figure}%[htbp]%[h]%[ht]%[htbp!]
%     \centering	
%     \includegraphics[width=0.5295\textwidth]{figures/FMMNNerrorsLD2D_w500r20d8.pdf}
%     % \subcaption{}
%     %		\label{subfig:f3}
%     \caption{FMMNN: Errors (in MSE) vs. epoch.  }
%     	% \label{fig:errors:LD:fun2D}
%     % 		\vskip 0.02in
% \end{figure}
% \begin{figure}%[htbp]%[h]%[ht]%[htbp!]
%     \centering	
%     \includegraphics[width=0.5295\textwidth]{figures/FMMNNvsReLUMMNNerrorsLD2D_w500r20d8.pdf}
%     % \subcaption{}
%     %		\label{subfig:f3}
%     \caption{FMMNN vs ReLU MMNN: test errors (in MSE) vs. epoch.  FMMNN MSE $1.04\times 10^{-6}$
%     vs ReLU MMNN $1.97\times 10^{-6}$.}
%     	% \label{fig:errors:LD:fun2D}
%     % 		\vskip 0.02in
% \end{figure}

% In this example, FMMNN converges fast than ReLU MMNN. 
% \begin{figure}%[htbp]%[h]%[ht]
%     \centering
%     \begin{subfigure}[c]{0.23\textwidth}
%         \centering
%         \includegraphics[width=0.998055\textwidth]{figures/FMMNNepoch1_idx1_w500r20d8_f.png}
%         \subcaption{Epoch 1.}
%     \end{subfigure}
%     \hfill
%     \begin{subfigure}[c]{0.23\textwidth}
%         \centering
%         \includegraphics[width=0.998055\textwidth]{figures/FMMNNepoch3_idx1_w500r20d8_f.png}
%         \subcaption{Epoch 3.}
%     \end{subfigure}
%     \hfill
%     \begin{subfigure}[c]{0.23\textwidth}
%         \centering
%         \includegraphics[width=0.998055\textwidth]{figures/FMMNNepoch5_idx1_w500r20d8_f.png}
%         \subcaption{Epoch 5.}
%     \end{subfigure}
%     \hfill
%     \begin{subfigure}[c]{0.23\textwidth}
%         \centering
%         \includegraphics[width=0.998055\textwidth]{figures/FMMNNepoch7_idx1_w500r20d8_f.png}
%         \subcaption{Epoch 7.}
%     \end{subfigure}
% \\
% \begin{subfigure}[c]{0.23\textwidth}
%         \centering
%         \includegraphics[width=0.998055\textwidth]{figures/FMMNNepoch14_idx1_w500r20d8_f.png}
%         \subcaption{Epoch 14.}
%     \end{subfigure}
%     \hfill
%     \begin{subfigure}[c]{0.23\textwidth}
%         \centering
%         \includegraphics[width=0.998055\textwidth]{figures/FMMNNepoch22_idx1_w500r20d8_f.png}
%         \subcaption{Epoch 22.}
%     \end{subfigure}
%     \hfill
%     \begin{subfigure}[c]{0.23\textwidth}
%         \centering
%         \includegraphics[width=0.998055\textwidth]{figures/FMMNNepoch30_idx1_w500r20d8_f.png}
%         \subcaption{Epoch 30.}
%     \end{subfigure}
%     \hfill
%     \begin{subfigure}[c]{0.23\textwidth}
%         \centering
%         \includegraphics[width=0.998055\textwidth]{figures/FMMNNepoch300_idx1_w500r20d8_f.png}
%         \subcaption{Epoch 300.}
%     \end{subfigure}
    
%     \caption{FMMNN: Illustration of the learning dynamics.}
%     % \label{fig:2D-dynamics}
% \end{figure}


% }


% \begin{figure}%[htbp]%[h]%[htbp!]
%     \centering	
%     \foreach \iiidx in {5,10,15,1000}
%     {
%     \begin{subfigure}[c]{0.23\textwidth}
%     \centering            \includegraphics[width=0.998055\textwidth]{figures/LHF2D2/epoch\iiidx_idx1_w100r10d6_f.png}
%     \subcaption{Epoch \iiidx.}
%     %		\label{subfig:f3}
%     \end{subfigure}
%     }
%     \caption{The difference between the learned neural network and the target function.}
%     	% \label{fig:}
%     % 		\vskip 0.02in
% \end{figure}





% Next, we consider another function, which is defined as
% \begin{equation*}
%     f(r,\theta)=\begin{cases}
%         0 & \tn{if  } 0.5+5\rho-5r\le 0,\\
%         1 & \tn{if  } 0.5+5\rho-5r\ge 1,\\
%        0.5+5\rho-5r & \tn{otherwise},\\
%     \end{cases}
%     \quad \tn{where}\quad \rho=0.5+0.1\cos(\pi  \theta^2).
% \end{equation*}
% Refer to Figure~\ref{fig:f3D:local:HF} for visual representations of this function.


% We utilize a compact network for approximation, configured with a width of 100, a rank of 10, and a depth of 6 layers. The following hyperparameters are used during the training of our neural networks:
% \begin{itemize}
% \item \textbf{Learning Rate}: The learning rate for epoch-$k$ is set at $0.001 \times 0.9^{\lfloor k/25 \rfloor}$.
% \item \textbf{Mini-Batch Size}: 1000.
% \item \textbf{Sampling}: We employ uniform grid sampling, with $400\times  400$ samples across the domain $[-1,1]^2$, totaling 160,000 samples.
% \end{itemize}



% \begin{figure}%[htbp]%[h]%[htbp!]
%     \centering	
%     \begin{subfigure}[c]{0.48\textwidth}
%     \centering            \includegraphics[height=0.8\textwidth]{figures/LHF2D/LHF2D.png}
%     % \subcaption{}
%     %		\label{subfig:f3}
%     \end{subfigure}
%     \hfill
%     \begin{subfigure}[c]{0.48\textwidth}
%     \centering            \includegraphics[height=0.8\textwidth]{figures/LHF2D/LHFf3D.pdf}
%     % \subcaption{}
%     %		\label{subfig:f3}
%     \end{subfigure}
%     \caption{Illustrations of the target function.}
%     	\label{fig:f3D:local:HF}
%     % 		\vskip 0.02in
% \end{figure}


% \begin{figure}%[htbp]%[h]%[htbp!]
%     \centering	
%     \foreach \iiidx in {10,20,50,1000}
%     {
%     \begin{subfigure}[c]{0.23\textwidth}
%     \centering            \includegraphics[width=0.998055\textwidth]{figures/LHF2D/epoch\iiidx_idx1_w100r10d6_diff.png}
%     \subcaption{Epoch \iiidx.}
%     %		\label{subfig:f3}
%     \end{subfigure}
%     }
%     \caption{Illustrations of the learned neural network in different epochs.}
%     	% \label{fig:}
%     % 		\vskip 0.02in
% \end{figure}

% \begin{figure}%[htbp]%[h]%[htbp!]
%     \centering	
%     \foreach \iiidx in {10,20,50,1000}
%     {
%     \begin{subfigure}[c]{0.23\textwidth}
%     \centering            \includegraphics[width=0.998055\textwidth]{figures/LHF2D/epoch\iiidx_idx1_w100r10d6_f.png}
%     \subcaption{Epoch \iiidx.}
%     %		\label{subfig:f3}
%     \end{subfigure}
%     }
%     \caption{The difference between the learned neural network and the target function.}
%     	% \label{fig:}
%     % 		\vskip 0.02in
% \end{figure}

% \begin{figure}%[htbp]%[h]%[htbp!]
%     \centering	
%     \includegraphics[width=0.8\textwidth]{figures/LHF2D/errorsLHF2D_w100r10d6.pdf}
%     % \subcaption{}
%     %		\label{subfig:f3}
%     \caption{Illustrations of training and test errors in difference epochs, measured in MSE.}
%     	% \label{fig:}
%     % 		\vskip 0.02in
% \end{figure}






% \subsubsection{Higher dim cases}
 
% \subsection{Dual-phase learning approach}

% \begin{equation*}
%     f(x)=(1-x^2)\sin(20\pi x)
% \end{equation*}
% Learning such a target function
% $f$ is challenging, as previously mentioned. To bridge this gap, we suggest a dual-phase learning strategy, grounded in the observation that if a neural network (NN) can learn a local feature effectively, it can then rapidly acquire other features.

% Our method involves two stages:

% \begin{itemize}
%     \item Initial Stage: This focuses on identifying a basic feature by implementing a mask ($\one_{\{|x|<0.1\}}$).
% \item
% Subsequent Stage: In this phase, we proceed with standard training, this time without the mask.
% \end{itemize}

% See GIFs for illustrations (200 epochs in total; our dual-phase method using mask for first 60 epochs): 
% \blue{output\_idx0.gif} for the normal method; 
% \blue{output\_idx1.gif} for our dual-phase method.

% \bigskip
% We want to note that although the mask $\exp(-100x^2)$ is functional, it does not perform as efficiently as $\one_{\{|x|<0.1\}}$.

% To test the effectiveness of the new method on complex functions, we experiment with 
% $f(x)=(1-x^2)\sin(40\pi x)$. Initially, we use a mask for the first 100 epochs, followed by 400 epochs without a mask ($\one_{\{|x|<0.05\}}$). See \blue{output\_idx2.gif} for visual illustrations.


% \section{Further discussion}
% Here we make a few general comments about MMNNs. 
% \begin{itemize}
%     \item The two key differences between a standard fully connected network (FCNN) or multi-layer perceptrons (MLP)  and an MMNN are 1) the introduction of the weights $\bmA, \bmc$ for different linear combinations of hidden neurons (or perceptrons) as the multi-components in each layer, and 2) the training strategy that fixes those randomly initialized $\bmW,\bmb$ (random features) in the hidden neurons. Hence it is extremely easy to modify a FCNN or MLP to MMNN. 
%     \item 
%     MMNNs are much more effective than FCNNs in terms of representation, training, and accuracy especially for complex functions. In comparison, as shown in those experiments in Section~\ref{sec:comparison}, MMNNs 1) have much fewer training parameters, 2) converge much faster in training, 3) achieve much better accuracy. Moreover, experiments show that training process of MMNNs converges not only faster but also with a consistant rate while FCNNs saturates pretty early to a quite low accuracy, as commonly observed in practices. These nice  behaviors of MMNNs are due to their balanced structure for smooth decomposition as well as the training strategy. In practice, 
%     the introduction of $\bmA, \bmc$ in MMNNs provides an important balance between the network width, which is the number of hidden neurons (basis functions) and can be very large, and the dimension of the input space, which is the number of components from the previous layer and can be much smaller than the network width. In other words, using a few linear combinations of the basis functions can capture smooth structures in the input space well. On the other hand, for FCNNs the two are the same and no balance is exerted.  
%     \item 
%     There are three hyperparameters for the configuration of MMNNs, the network width, the number of components (rank), and the number of layers (depth). Here are the general guidelines based on our mathematical construction and extensive experiments:\\
%     1. The network width should provide enough resolution to resolve fine details of the target function. This means that the width should be at least comparable to the size of an adaptive mesh that can approximate the target function well. \\
%     2. The number of components (rank) is related to the overall complexity  of the target function which depends on its spatial domain and Fourier domain representation as well as the input dimension. As indicated by our mathematical multi-component construction, it is related to the "divide and conquer" strategy. \\
%     3. The number of layers (depth) is also related to the overall complexity of the target function as for the number of components. Rank and depth are complementary but work together effectively for a smooth decomposition of the target function. The rule of thumb for depth is similar to that for the rank. \\

%     Here we use more concrete examples to illustrate the guidelines. For simplicity we fix the input dimension and domain of interest. As the domain size and dimension increases, the network size needs to increase correspondingly. For a smooth target function, a compact MMNN in terms of width, rank, and depth is enough and easy training process can render accurate results. Larger MMNNs are needed for target functions with localized rapid changes. Even with a relative compact size, the training process can allocate resources adaptive to the target function and render good approximation. The most difficult situation is to approximate globally highly oscillatory functions with diverse Fourier modes for which large MMNNs are needed. For instance, if the oscillation frequency doubles, the network width should increase by $2^d$ where $d$ is the dimension. In general the network width needs to deal with the curse of dimensionality just like a mesh based method. However, the growth of the number of components and layers with the increase of complexity seems to be relative mild (maybe polylogarithmic suggested by our mathematical construction).
%     \item
%     Overall, for a given target function, MMNNs can work well with quite a large range of configuration with a trade-off between the network size and training process. For example, the training process for a network more on the compact size with respect to the complexity of a given target function may become more subtle and challenging, e.g., choosing the appropriate learning rate and batch size, due to the lack of flexibility (or redundancy) of the representation. On the other hand, a network of too large size (or redundancy) with respect to the complexity of a given target function requires unnecessarily expensive training cost.
%     \item
%     The most advantageous situation for using MMNNs is when approximating a function in relative high dimension which is mostly smooth except for localized fine features, e.g., a distribution in high dimensions concentrated on a low dimensional manifold. Through training, MMNNs can provide an automatic adaptive approximation of the underlying structure which is very challenging for a traditional mesh based method.  

    
%     \item 
%     We also tried using different activation functions for MMNNs. In general, \texttt{ReLU} provides the best results. However, in situations where a smooth approximation is needed, one might consider using smooth alternatives to \ReLU\ such as \GELU\ \cite{2016arXiv160608415H} or \Swish\ \cite{2017arXiv171005941R}, which generally yield results comparable to \ReLU. In contrast, other popular S-shaped activation functions like \Sigmoid\ and \Tanh\ have shown poor performance in our tests possibly due to the vanishing derivative problem. 
% \end{itemize}
    


% \clearpage

\section{Further discussion}
\label{sec:further:discussion}

% Here we make a few general comments about MMNNs. 
In this section, we provide a few more comments about MMNNs. First, in Section~\ref{sec:advantages:MMNN:vs:FCNN}, we discuss the advantages of MMNNs over fully connected networks (FCNNs) or multi-layer perceptrons (MLPs). Next, in Section~\ref{sec:guideline:MMNN:size}, we offer practical guidelines for determining the appropriate MMNN size based on our theoretical understanding and extensive numerical experiments. Finally in Section~\ref{sec:other:activation:functions}, we discuss the use of alternative activation functions beyond \ReLU\ in MMNNs.


\subsection{Advantages compared to FCNNs or MLPs}
\label{sec:advantages:MMNN:vs:FCNN}

%\begin{itemize}
The two key differences between a standard FCNN or MLP  and an MMNN are (1) the introduction of the weights $\bmA, \bmc$ for different linear combinations of hidden neurons (or perceptrons) as the multi-components in each layer, and (2) the training strategy that fixes those randomly initialized $\bmW,\bmb$ (random features) in the hidden neurons. Hence it is extremely easy to modify a FCNN or MLP to an MMNN. 

    MMNNs are much more effective than FCNNs in terms of representation, training, and accuracy especially for complex functions. In comparison, as shown in those experiments in Section~\ref{sec:comparison}, MMNNs (1) have much fewer training parameters, (2) converge much faster in training, (3) achieve much better accuracy. Moreover, experiments show that training process of MMNNs converges not only faster but also with a steady rate while FCNNs saturates pretty early to a quite low accuracy, as commonly observed in practices. These nice  behaviors of MMNNs are due to their balanced structure for smooth decomposition as well as the training strategy. In practice, 
    the introduction of $\bmA, \bmc$ in MMNNs provides an important balance between the network width, which is the number of hidden neurons (basis functions) and can be very large, and the dimension of the input space, which is the number of components from the previous layer and can be much smaller than the network width. In other words, using a few linear combinations of the basis functions can capture smooth structures in the input space well. On the other hand, for FCNNs the two are the same and no balance is exerted.  
%\end{itemize}



\subsection{Practical guidelines for MMNNs}
\label{sec:guideline:MMNN:size}

    There are three hyperparameters for the configuration of MMNN sizes, the network width, the number of components (rank), and the number of layers (depth). Here are the general guidelines based on our mathematical construction and extensive experiments:
%    \\
    \begin{enumerate}
   \item The network width should provide enough resolution to capture fine details of the target function. This means that the width should be at least comparable to the size of an adaptive mesh that can approximate the target function well. 
\item The number of components (rank) is related to the overall complexity of the target function which depends on its spatial domain and Fourier domain representation as well as the input dimension. As indicated by our mathematical multi-component construction, it is related to the ``divide and conquer" strategy. 
\item The number of layers (depth) is also related to the overall complexity of the target function as for the number of components. Rank and depth are complementary but work together effectively for a smooth decomposition of the target function. The rule of thumb for depth is similar to that for the rank. 
    \end{enumerate}
    % 1. The network width should provide enough resolution to resolve fine details of the target function. This means that the width should be at least comparable to the size of an adaptive mesh that can approximate the target function well. \\
    % 2. The number of components (rank) is related to the overall complexity of the target function which depends on its spatial domain and Fourier domain representation as well as the input dimension. As indicated by our mathematical multi-component construction, it is related to the ``divide and conquer" strategy. \\
    % 3. The number of layers (depth) is also related to the overall complexity of the target function as for the number of components. Rank and depth are complementary but work together effectively for a smooth decomposition of the target function. The rule of thumb for depth is similar to that for the rank. \\

    Here we use more concrete examples to illustrate the guidelines. For simplicity we fix the input dimension and domain of interest. As the domain size and dimension increases, the network size needs to increase correspondingly. For a smooth target function, a compact MMNN in terms of width, rank, and depth is enough and easy training process can render accurate results. Larger MMNNs are needed for target functions with localized rapid changes. Even with a relative compact size, the training process can allocate resources adaptive to the target function and render good approximation. The most difficult situation is to approximate globally highly oscillatory functions with diverse Fourier modes for which large MMNNs are needed. For instance, if the oscillation frequency doubles, the network width should increase by $2^d$ where $d$ is the dimension. In general the network width needs to deal with the curse of dimensionality just like a mesh based method. However, the growth of the number of components and layers with the increase of complexity seems to be relative mild (maybe polylogarithmic suggested by our mathematical construction).
%\begin{itemize}

    Overall, for a given target function, MMNNs can work well with quite a large range of configuration with a trade-off between the network size and training process. For example, the training process for a network more on the compact size with respect to the complexity of a given target function may become more subtle and challenging, e.g., choosing the appropriate learning rate and batch size, due to the lack of flexibility (or redundancy) of the representation. On the other hand, a network of too large size (or redundancy) with respect to the complexity of a given target function requires unnecessarily expensive training cost. There is a trade-off between representation and optimization one needs to balance in practice. An important question for future research is how to develop a self-adaptive strategy to adjust the network size.
%\item
    
    The most advantageous situation for using MMNNs is when approximating a function in relative high dimension which is mostly smooth except for localized fine features, e.g., a distribution in high dimensions concentrated on a low dimensional manifold. Through training, MMNNs can provide an automatic adaptive approximation of the underlying structure which can be challenging for a mesh based method. 
 %   \item 
 
    We would like to remark that learning rate scheduler can be a subtle and important issue for all training process in practice. For all our training process, the Step Learning Rate suffices. However, one could consider using other learning rate schedulers, such as the Cosine Scheduler \cite{loshchilov2017sgdr} or the gradual warm-up strategy \cite{2017arXiv170602677G}. Exploring and designing a more efficient learning rate scheduler with some automatic restart mechanism is a potential interesting topic for future work.

    
    % \item 
    % We also tried using different activation functions for MMNNs. In general, \texttt{ReLU} provides the best results. However, in situations where a smooth approximation is needed, one might consider using smooth alternatives to \ReLU\ such as \GELU\ \cite{2016arXiv160608415H} or \Swish\ \cite{2017arXiv171005941R}, which generally yield results comparable to \ReLU. In contrast, other popular S-shaped activation functions like \Sigmoid\ and \Tanh\ have shown poor performance in our tests possibly due to the vanishing derivative problem.    
%\end{itemize}


\subsection{Beyond \texttt{ReLU} to other activation functions}
\label{sec:other:activation:functions}

We also tried using different activation functions for MMNNs, e.g., \GELU\ \cite{2016arXiv160608415H}, \Swish\ \cite{2017arXiv171005941R}, \Sigmoid, and \Tanh. In general, \texttt{ReLU} provides the overall best results for various target functions. 
However, in situations where a smooth (e.g., $C^s$ or real analytic) approximation  is needed, one might consider using smooth alternatives to \ReLU\ such as \GELU\  or \Swish, which generally yield results comparable to \ReLU.

% \begin{itemize}
%     \item For target functions that are $C^s$ or even real analytic (and can be highly oscillatory), such as $f(x) = \cos(36\pi x^2) - 0.6\cos(12\pi x^2)$, \GELU\ (or \Swish) tends to perform slightly better than \ReLU.
%     \item For target functions with non-differentiable points, such as $f(x) = \one_{\{|x| < 0.02\}} \cdot \sin(50\pi x)$, \GELU\ (or \Swish) generally performs slightly worse than \ReLU.
%     \item The use of \GELU\ (or \Swish) typically results in slightly longer training time compared to \ReLU.
% \end{itemize}
Additionally, other popular S-shaped activation functions like \Sigmoid\ and \Tanh\ have demonstrated poor performance in our tests, possibly due to the vanishing gradient problem. For highly oscillatory target functions, when using \Sigmoid\ or \Tanh\ training errors did not even decrease  during the training process.



\section{Conclusion}
\label{sec:conclusion}

In this work, we introduced the Multi-component and Multi-layer Neural Network (MMNN) and demonstrated its effectiveness in approximating complex functions. By incorporating the principles of structured and balanced decomposition, the MMNN architecture addresses the limitations of shallow networks, particularly in capturing high-frequency components and localized fine features. Our proposed network structure as confirmed by extensive numerical experiments can approximate highly oscillatory functions and functions with rapid transitions efficiently and accurately. Additionally, we highlight the advantages of our training strategy, which optimizes only the linear combination weights of basis functions for each component while keeping the parameters within the activation (basis) functions fixed, leading to a more efficient and stable training process.

The theoretical underpinnings and practical implementations presented in this paper suggest that MMNNs offer a promising direction for constructing neural networks capable of handling complex tasks with fewer parameters and reduced computational overhead. Future research can explore further generalizations and applications of MMNNs, as well as investigate the interplay between representation and optimization in more depth.

% \clearpage
%%%%%%%%%%%%%%%%%%%%%%%%%%%%%%%%%%%%%%%%%%%%%
\section*{Acknowledgments}

S. Zhang was partially supported by the 
start-up fund P0053092
from The Hong Kong Polytechnic University.
H. Zhao was partially supported by NSF grants DMS-2309551, and DMS-2012860. Y. Zhong was partially supported by NSF grant DMS-2309530. H. Zhou was partially supported by NSF grant DMS-2307465.

% S. Zhang was partially supported by the start-up fund P0053092
% from The Hong Kong Polytechnic University.
% H. Zhao was partially supported by NSF grants DMS-2309551, and DMS-2012860. Y. Zhong was partially supported by NSF grant DMS-2309530, H. Zhou was partially supported by NSF grant DMS-2307465.


% %%%%%%%%%%%%%%%%%%%%%%%%%%%%%%%%%%%%%%%%%%%%%%%%%%%%%%%%%%%%%%%%%
% \bibliographystyle{siamplain}
% \bibliography{references}
{
\small
\let\oldhref\href
\def\href#1#2{\oldhref{#1}{\nolinkurl{#2}}}
%%%\clearpage
% \bibliographystyle{siam}
% \bibliographystyle{plainnat}
\bibliographystyle{plain}   
\bibliography{references}
}
%%%%%%%%%%%%%%%%%%%%%%%%%%%%%%%%%%%

% \end{document}

% \clearpage
% \appendix 

% \section{Proofs of Theorems~\ref{thm:decomposition:1D} and \ref{thm:decomposition:2D}}

% \begin{proof}[Proof of Theorem~\ref{thm:decomposition:1D}]
% By definition of $\psi_i$ in Equation~\eqref{eq:1D:psii}, 
% it is easy to check
% \begin{equation*}%\label{eq:1D:psi}
%     \psi_i(x)=\begin{cases}
%         b_i &\tn{if  } x>x_{i},\\
%         \calL_i^{-1}(x) &\tn{if  }x\in [x_{i-1},x_i],\\
%         a_i &\tn{if  }x<x_{i-1},
%     \end{cases}
%     \quad \Longrightarrow \quad 
%         \psi_i(x)=\begin{cases}
%         b_i &\tn{if  }i\le j-1,\\
%         \calL_j^{-1}(x) &\tn{if  }i= j ,\\
%         a_i &\tn{if  }i\ge j+1,
%     \end{cases}
% \end{equation*}
% for a fixed $j\in \{1,2,\cdots,n\}$ and any $x\in [x_{j-1}, x_j]$.
% % Fixing $j$, for any $x\in [x_{j-1}, x_j]$, we have
% % \begin{equation*}
% %     \psi_i(x)=\begin{cases}
% %         b_i &\tn{if  }i\le j-1,\\
% %         \calL_j^{-1}(x) &\tn{if  }i= j ,\\
% %         a_i &\tn{if  }i\ge j+1,
% %     \end{cases}
% % \end{equation*}
% % Then $\psi_i(x)=x_i$ if $i\le j-1$
% % and $\psi_i(x)=x_{i-1}$ if $i\ge j+1$.
% % Then we have
% % from which we deduce
% % \begin{equation*}
% %     \begin{split}
% %         \sum_{i=1}^n f_i\circ \psi_i(x)=\sum_{i=1}^n f\circ \calL_i\circ \psi_i(x)
% %     &=\sum_{i=1}^{j-1} f\circ \calL_i\circ \psi_i(x)
% %     +f\circ \calL_j\circ \psi_j(x)
% %     +\sum_{i=j+1}^n f\circ \calL_i\circ \psi_i(x)
% %     \\& = \sum_{i=1}^{j-1} f\circ \calL_i(b_i)
% %     +f\circ \calL_j\circ\calL_j^{-1}(x) 
% %     +\sum_{i=j+1}^n f\circ \calL_i(a_i)
% %     \\& = \sum_{i=1}^{j-1} f (x_i)
% %     +f(x)
% %     +\sum_{i=j+1}^n f (x_{i-1})
% %     =f(x)+\underbrace{\sum_{i=1}^{n-1}f(x_i)}_{\tn{constant}}.
% %     \end{split}
% % \end{equation*}
% It follows that
% \begin{equation*}
%     \begin{split}
%         & \phantom{=\;\;}
%         \sum_{i=1}^n f_i\circ \psi_i(x)
%         \\&=\sum_{i=1}^n f\circ \calL_i\circ \psi_i(x)
%     =\sum_{i=1}^{j-1} f\circ \calL_i\circ \psi_i(x)
%     +f\circ \calL_j\circ \psi_j(x)
%     +\sum_{i=j+1}^n f\circ \calL_i\circ \psi_i(x)
%     \\& = \sum_{i=1}^{j-1} f\circ \calL_i(b_i)
%     +f\circ \calL_j\circ\calL_j^{-1}(x) 
%     +\sum_{i=j+1}^n f\circ \calL_i(a_i)
%     \\& = \sum_{i=1}^{j-1} f (x_i)
%     +f(x)
%     +\sum_{i=j+1}^n f (x_{i-1})
%     =f(x)+\underbrace{\sum_{i=1}^{n-1}f(x_i)}_{\tn{constant}}.
%     \end{split}
% \end{equation*}
% It follows that
% \begin{equation*}
%     \begin{split}
%         f(x)=\sum_{i=1}^n f_i\circ \psi_i(x)
%     -\underbrace{\sum_{i=1}^{n-1}f(x_i)}_{\tn{constant}}\quad \tn{for any $x\in [x_{j-1},x_j]$.}
%     \end{split}
% \end{equation*}
% Since $j$ is arbitrary, the above equation holds for all
% $x=\cup_{j=1}^n[x_{j-1}, x_j]= [x_0,x_n]$. 
% \end{proof}

% \begin{proof}[Proof of Theorem~\ref{thm:decomposition:2D}]
% % of Equation~\eqref{eq:highD:decomposition}.
% Fixing $(k,j)$, for any $(x,y)\in [x_{k-1}, x_k]\times [y_{\ell-1}, y_\ell]$, we have
% \begin{equation*}
%     \psi_i(x)=\begin{cases}
%         b_i &\tn{if  }i\le k-1,\\
%         \calL_{1,k}^{-1}(x) &\tn{if  }i= k ,\\
%         a_i &\tn{if  }i\ge k+1 
%     \end{cases}
%         \quad \tn{and}\quad 
%         \phi_j(y)=\begin{cases}
%         d_j &\tn{if  } j \le \ell-1,\\
%         \calL_{2,\ell}^{-1}(y) &\tn{if  } j=\ell,\\
%         c_j &\tn{if  } j\ge \ell+1.
%     \end{cases}
% \end{equation*}
% It follows that 
% \begin{equation*}
%     \begin{split}
%           &\phantom{=\;}\sum_{i=1}^n f_{i,0}\Big( \psi_i(x), \, y\Big)
%         = \sum_{i=1}^n f\Big( \calL_{1,i}\circ \psi_i(x), \, y\Big)
%         \\
%         &=
%         \sum_{i=1}^{k-1} f\Big( \calL_{1,i}\circ \psi_i(x), \, y\Big)
%         +f\Big( \calL_{1,k}\circ \psi_k(x), \, y\Big)
%         +
%         \sum_{i=k+1}^n f\Big( \calL_{1,i}\circ \psi_i(x), \, y \Big)
%         \\
%         &=
%         \sum_{i=1}^{k-1} f\Big( \calL_{1,i}(b_i), \, y\Big)
%         +f\Big( \calL_{1,k}\circ \calL_{1,k}^{-1}(x), \, y\Big)
%         +
%         \sum_{i=k+1}^n f\Big( \calL_{1,i}(a_i), \, y\Big)
%         \\
%         &=
%         \sum_{i=1}^{k-1} f( x_i, \, y)
%         +f( x, \, y)
%         +
%         \sum_{i=k+1}^n f( x_{i-1}, \, y)
% =
%         f( x, \, y)
%         +\sum_{i=1}^{n-1} f( x_i, \, y),
%     \end{split}
% \end{equation*}
% implying 
% \begin{equation*}
%     \begin{split}          
%         f( x, \, y)=\sum_{i=1}^n f_{i,0}\Big( \psi_i(x), \, y\Big)-  \sum_{i=1}^{n-1} f( x_i, \, y).
%     \end{split}
% \end{equation*}
% For each \(i\), using the one-dimensional decomposition technique described in Section~\ref{sec:decomposition:1D}, we find the decompositions for 
% \(f_{i,0}\big( \psi_i(x), \, y\big)\) and \(f( x_i, \, y)\).
% We have 
% % \begin{equation*}
% %     \begin{split}
% %           &\phantom{=\;}\sum_{j=1}^m f_{i,j}\Big( \psi_i(x), \, \phi_j(y)\Big)
% %         = \sum_{j=1}^m f_{i,0}\Big(\psi_i(x),\, \calL_{2,j}\circ \phi_j(y)\Big)\\
% %         &=
% %         \sum_{j=1}^{\ell-1} f_{i,0}\Big(\psi_i(x),\, \calL_{2,j}\circ \phi_j(y)\Big)
% %         +f_{i,0}\Big(\psi_i(x),\, \calL_{2,\ell }\circ \phi_\ell(y)\Big)
% %         +
% %         \sum_{j=\ell+1}^m f_{i,0}\Big(\psi_i(x),\, \calL_{2,j}(\circ \phi_j(y))\Big)
% %         \\
% %         &=
% %         \sum_{j=1}^{\ell-1} f_{i,0}\Big(\psi_i(x),\, \calL_{2,j}(d_j)\Big)
% %         +f_{i,0}\Big(\psi_i(x),\, \calL_{2,\ell }\circ \calL_{2,\ell }^{-1}(y)\Big)
% %         +
% %         \sum_{j=\ell +1}^m f_{i,0}\Big(\psi_i(x),\, \calL_{2,j}(c_j)\Big)
% %         \\
% %         &=
% %         \sum_{j=1}^{\ell-1} f_{i,0}\Big(\psi_i(x),\, y_j\Big)
% %         +f_{i,0}\Big(\psi_i(x),\, y\Big)
% %         +
% %          \sum_{j=\ell+1}^m f_{i,0}\Big(\psi_i(x),\, y_{j-1}\Big)\\
% %         &=f_{i,0}\Big(\psi_i(x),\, y\Big)
% %         +
% %         \sum_{j=1}^{m-1} f_{i,0}\Big(\psi_i(x),\, y_{j-1}\Big),
% %     \end{split}
% % \end{equation*}
% \begin{equation*}
%     \begin{split}
%           &\phantom{=\;}\sum_{j=1}^m f_{i,j}\Big( \psi_i(x), \, \phi_j(y)\Big)
%         = \sum_{j=1}^m f_{i,0}\Big(\psi_i(x),\, \calL_{2,j}\circ \phi_j(y)\Big)\\
%         &=
%         \sum_{j=1}^{\ell-1} f_{i,0}\Big(\psi_i(x),\, \calL_{2,j}\circ \phi_j(y)\Big)
%         +f_{i,0}\Big(\psi_i(x),\, \calL_{2,\ell }\circ \phi_\ell(y)\Big)
%         +
%         \sum_{j=\ell+1}^m f_{i,0}\Big(\psi_i(x),\, \calL_{2,j}(\circ \phi_j(y))\Big)
%         \\
%         &=
%         \sum_{j=1}^{\ell-1} f_{i,0}\Big(\psi_i(x),\, \calL_{2,j}(d_j)\Big)
%         +f_{i,0}\Big(\psi_i(x),\, \calL_{2,\ell }\circ \calL_{2,\ell }^{-1}(y)\Big)
%         +
%         \sum_{j=\ell +1}^m f_{i,0}\Big(\psi_i(x),\, \calL_{2,j}(c_j)\Big)
%         \\
%         &=
%         \sum_{j=1}^{\ell-1} f_{i,0}\Big(\psi_i(x),\, y_j\Big)
%         +f_{i,0}\Big(\psi_i(x),\, y\Big)
%         +
%          \sum_{j=\ell+1}^m f_{i,0}\Big(\psi_i(x),\, y_{j-1}\Big)
%          \\&=f_{i,0}\Big(\psi_i(x),\, y\Big)
%         +
%         \sum_{j=1}^{m-1} f_{i,0}\Big(\psi_i(x),\, y_{j-1}\Big),
%     \end{split}
% \end{equation*}
% implying 
% \begin{equation*}
%     \begin{split}
%           f_{i,0}\Big(\psi_i(x),\, y\Big)
%         =\sum_{j=1}^m f_{i,j}\Big( \psi_i(x), \, \phi_j(y)\Big)-
%         \sum_{j=1}^{m-1} f_{i,0}\Big(\psi_i(x),\, y_{j}\Big).
%     \end{split}
% \end{equation*}
% Moreover, 
%  \begin{equation*}
%     \begin{split}
%           &\phantom{=\;}\sum_{j=1}^m f_{0,j}\Big( x_i, \, \phi_j(y)\Big)
%         = \sum_{j=1}^m f \Big(x_i,\, \calL_{2,j}\circ \phi_j(y)\Big)\\
%         &=
%         \sum_{j=1}^{\ell-1} f \Big(x_i,\, \calL_{2,j}\circ \phi_j(y)\Big)
%         +f \Big(x_i,\, \calL_{2,\ell}\circ \phi_\ell(y)\Big)
%         +
%          \sum_{j=\ell+1}^mf \Big(x_i,\, \calL_{2,j}\circ \phi_j(y)\Big)
%         \\
%         &=
%         \sum_{j=1}^{\ell-1} f \Big(x_i,\, \calL_{2,j}(d_j)\Big)
%         +f \Big(x_i,\, \calL_{2,\ell }\circ \calL_{2,\ell }^{-1}(y)\Big)
%         +
%         \sum_{j=\ell+1}^m f \Big(x_i,\, \calL_{2,j}(c_j)\Big)
%         \\
%         &=
%         \sum_{j=1}^{\ell-1} f (x_i,\, y_j)
%         +f (x_i,\, y)
%         +
%         \sum_{j=\ell+1}^m f (x_i,\, y_{j-1})        =f (x_i,\, y)
%         +
%         \sum_{j=1}^{m-1} f (x_i,\, y_{j}),
%     \end{split}
% \end{equation*}
% implying 
% \begin{equation*}
%     \begin{split}
%           f (x_i,\, y)
%         =\sum_{j=1}^m f_{0,j}\Big( x_i, \, \phi_j(y)\Big)-
%         \sum_{j=1}^{m-1} f (x_i,\, y_{j}).
%     \end{split}
% \end{equation*}
% Therefore, for any $(x,y)\in [x_{k-1}, x_k]\times [y_{\ell-1}, y_\ell]$,
% \begin{equation*}
%     \begin{split}          
%         f( x, \, y)
%         &=\sum_{i=1}^n f_{i,0}\Big( \psi_i(x), \, y\Big)-  \sum_{i=1}^{n-1} f( x_i, \, y)
%         \\& =\sum_{i=1}^n 
%         \Bigg(\sum_{j=1}^m f_{i,j}\Big( \psi_i(x), \, \phi_j(y)\Big)-
%         \sum_{j=1}^{m-1} f_{i,0}\Big(\psi_i(x),\, y_{j}\Big)\Bigg)
%         \\ & \quad\   -  \sum_{i=1}^{n-1}\Bigg(\sum_{j=1}^m f_{0,j}\Big( x_i, \, \phi_j(y)\Big)-
%         \sum_{j=1}^{m-1} f (x_i,\, y_{j})\Bigg)
%         \\& =\sum_{i=1}^n \sum_{j=1}^m f_{i,j}\Big( \psi_i(x), \, \phi_j(y)\Big)-
%         \sum_{i=1}^n\sum_{j=1}^{m-1} f_{i,0}\Big(\psi_i(x),\, y_{j}\Big)
%         \\ & \quad\   -  \sum_{i=1}^{n-1} \sum_{j=1}^m f_{0,j}\Big( x_i, \, \phi_j(y)\Big)+
%         \sum_{i=1}^{n-1}\sum_{j=1}^{m-1} f (x_i,\, y_{j}).
%     \end{split}
% \end{equation*}
% % \begin{equation*}
% %     \begin{split}          
% %         &\phantom{=\;\;}f( x, \, y)
% %         =\sum_{i=1}^n f_{i,0}\Big( \psi_i(x), \, y\Big)-  \sum_{i=1}^{n-1} f( x_i, \, y)
% %         \\& =\sum_{i=1}^n 
% %         \Bigg(\sum_{j=1}^m f_{i,j}\Big( \psi_i(x), \, \phi_j(y)\Big)-
% %         \sum_{j=1}^{m-1} f_{i,0}\Big(\psi_i(x),\, y_{j}\Big)\Bigg)
% %           -  \sum_{i=1}^{n-1}\Bigg(\sum_{j=1}^m f_{0,j}\Big( x_i, \, \phi_j(y)\Big)-
% %         \sum_{j=1}^{m-1} f (x_i,\, y_{j})\Bigg)
% %         \\& =\sum_{i=1}^n \sum_{j=1}^m f_{i,j}\Big( \psi_i(x), \, \phi_j(y)\Big)-
% %         \sum_{i=1}^n\sum_{j=1}^{m-1} f_{i,0}\Big(\psi_i(x),\, y_{j}\Big)
% %            -  \sum_{i=1}^{n-1} \sum_{j=1}^m f_{0,j}\Big( x_i, \, \phi_j(y)\Big)+
% %         \sum_{i=1}^{n-1}\sum_{j=1}^{m-1} f (x_i,\, y_{j}).
% %     \end{split}
% % \end{equation*} 
% Since $k$ and $\ell$ are arbitrary, the above equation holds for all
% $(x,y)=\cup_{k=1}^n\cup_{\ell=1}^m[x_{k-1}, x_k]\times [y_{\ell-1}, y_\ell]= [x_0,x_n]\times [y_0,y_m]$.
% \end{proof}

% \section{Supplementary for experiments in Section~\ref{sec:local}}

% In the appendix, 
% % we first present an additional example for Section~\ref{sec:local} in Section~\ref{sec:one:more:local}. Then, 
% we provide numerical experiments in three and higher dimensions.
% in Section~\ref{sec:higher}.


% \section{One more example for Section~\ref{sec:local}}
% \label{sec:one:more:local}



% \subsection{Training process of 2D example in Section~\ref{sec:local}}


% \begin{figure}%[htbp]%[h]%[ht]
%     \centering
%     \begin{subfigure}[c]{0.23\textwidth}
%         \centering
%         \includegraphics[width=0.998055\textwidth]{figures/Local2D/epoch10_idx1_w100r10d6_f.png}
%         \subcaption{Epoch 10.}
%     \end{subfigure}
%     \hfill
%     \begin{subfigure}[c]{0.23\textwidth}
%         \centering
%         \includegraphics[width=0.998055\textwidth]{figures/Local2D/epoch50_idx1_w100r10d6_f.png}
%         \subcaption{Epoch 50.}
%     \end{subfigure}
%     \hfill
%     \begin{subfigure}[c]{0.23\textwidth}
%         \centering
%         \includegraphics[width=0.998055\textwidth]{figures/Local2D/epoch100_idx1_w100r10d6_f.png}
%         \subcaption{Epoch 100.}
%     \end{subfigure}
%     \hfill
%     \begin{subfigure}[c]{0.23\textwidth}
%         \centering
%         \includegraphics[width=0.998055\textwidth]{figures/Local2D/epoch1000_idx1_w100r10d6_f.png}
%         \subcaption{Epoch 1000.}
%     \end{subfigure}
%     \caption{Illustration of the training process for the 2D example in Section~\ref{sec:local}.}
%     \label{fig:2D-2}
% \end{figure}




%The illustration in Figure~\ref{fig:logErrorVsEpoch2} demonstrates that the network $h(x, y, z)$ successfully approximates the target function $f(x, y, z)$, achieving an impressively low MSE error of $1.02 \times 10^{-5}$ and a maximum error of $0.15$.

% Our architecture can learn global high frequency, e.g., $\sin(50\pi x)$, 2D $\sum_{i=1}^4 \cos \sin$; 2D star;


% \subsubsection{3D}

% Will add later


% \subsection{Local data sampling}
% Target function 
% \[
% f(x)=\one_{\{|x|<0.02\}} \cdot\sin(
%         50\pi  x)
% \]

% 18000 samples linspace(-0.06,0.06,18000)

% 2000 samples linspace(-1,1,2000)
% \begin{figure}%[htbp]%[h]%[htbp!]
%     \centering	
%     \foreach \iiidx in {7,8,9,10,15,20,25,30,60,90,...,160}
%     {
%     \begin{subfigure}[c]{0.31\textwidth}
%     \centering            \includegraphics[width=0.998055\textwidth]{figures/LDS1D/epoch\iiidx_idx1.png}
%     \subcaption{Epoch \iiidx.}
%     %		\label{subfig:f3}
%     \end{subfigure}
%     }
%     \caption{Illustrations of learned NN and the target function.}
%     	% \label{fig:}
%     % 		\vskip 0.02in
% \end{figure}


\end{document}